% Is There an Optimal Decision Theory? -- LuaLaTeX.  Build: sh build.sh  (or: lualatex, biber, lualatex x2)
\documentclass[11pt,a4paper,oneside]{scrartcl}

% ---- page and type ------------------------------------------------------------
\usepackage[textwidth=31pc,textheight=47\baselineskip,top=1.1in,headsep=2.2\baselineskip,
            footskip=3\baselineskip,heightrounded]{geometry}
\usepackage{amsmath,amsthm,mathtools}
\usepackage{unicode-math}
\setmainfont{Libertinus Serif}[Numbers=OldStyle]
\setsansfont{Libertinus Sans}
\setmathfont{Libertinus Math}
\setmathfont{Latin Modern Math}[range={\mathcal,\mathbfcal,\mathscr}]
\usepackage[babel=true]{microtype}
\usepackage[english]{babel}
\usepackage[autostyle]{csquotes}
\usepackage{booktabs,tabularx,array,enumitem,etoolbox}
\AtBeginEnvironment{tabular}{\upshape} % tables inside theorem statements stay upright
\usepackage{flafter} % floats never precede their place in the text
\usepackage[backend=biber,style=authoryear-comp,useprefix=true,maxcitenames=2,maxbibnames=10,giveninits=true,
            uniquename=false,uniquelist=false,doi=false,url=false,isbn=false,eprint=true,
            dashed=false]{biblatex}
\addbibresource{refs.bib}
\renewcommand*{\bibfont}{\small}
\setlength{\bibhang}{1.4em}
\renewbibmacro{in:}{\ifentrytype{article}{}{\printtext{\bibstring{in}\intitlepunct}}}
\usepackage{scrlayer-scrpage}
\deffootnote[1.4em]{1.4em}{1em}{\makebox[1.4em][l]{\thefootnotemark.}}
\setlength{\footnotesep}{.7\baselineskip}
\setkomafont{footnote}{\small}
\setkomafont{caption}{\small}
\setkomafont{captionlabel}{\scshape}
\setcapindent{0pt}
\renewcommand*{\captionformat}{.\enskip}
\usepackage[hidelinks,pdfusetitle]{hyperref}

\linespread{1.05}
\setlength{\parindent}{1.4em}
\widowpenalty=4000 \clubpenalty=4000 \interfootnotelinepenalty=0
\setlength{\emergencystretch}{1.5em}
\setlist{itemsep=.25em,topsep=.5em,parsep=0pt}

% ---- running heads --------------------------------------------------------------
\newcommand{\shorttitle}{Is There an Optimal Decision Theory?}
\newcommand{\theauthor}{Un crack}
\clearpairofpagestyles
\ohead{\textsc{\MakeLowercase{\shorttitle}}}
\ihead{\textsc{\MakeLowercase{\theauthor}}}
\cfoot*{\pagemark}
\setkomafont{pageheadfoot}{\normalfont\footnotesize}
\setkomafont{pagenumber}{\normalfont\small}
\pagestyle{scrheadings}

% ---- headings ---------------------------------------------------------------------
\setkomafont{disposition}{\normalfont}
\setkomafont{section}{\normalfont\large\scshape}
\setkomafont{subsection}{\normalfont\normalsize\itshape}
\RedeclareSectionCommand[beforeskip=-2\baselineskip plus -.4\baselineskip,
                         afterskip=.8\baselineskip]{section}
\RedeclareSectionCommand[beforeskip=-1.2\baselineskip plus -.3\baselineskip,
                         afterskip=.4\baselineskip]{subsection}
\setkomafont{part}{\normalfont\Large\scshape}
\setkomafont{partnumber}{\normalfont\Large\scshape}
\RedeclareSectionCommand[beforeskip=2.4\baselineskip,afterskip=1.4\baselineskip]{part}
\renewcommand*{\partformat}{\partname~\thepart.\enskip}
\renewcommand*{\sectionformat}{\thesection.\enskip}
\renewcommand*{\subsectionformat}{\thesubsection.\enskip}

% ---- statements ---------------------------------------------------------------------
\newtheoremstyle{dtplain}{.8\baselineskip}{.8\baselineskip}{\itshape}{}{\scshape}{.}{.5em}{}
\newtheoremstyle{dtdef}{.8\baselineskip}{.8\baselineskip}{\upshape}{}{\scshape}{.}{.5em}{}
\theoremstyle{dtplain}
\newtheorem{theorem}{Theorem}[section]
\newtheorem{proposition}[theorem]{Proposition}
\newtheorem{lemma}[theorem]{Lemma}
\newtheorem{corollary}[theorem]{Corollary}
\theoremstyle{dtdef}
\newtheorem{definition}[theorem]{Definition}
\newtheorem{example}[theorem]{Example}
\newtheorem{remark}[theorem]{Remark}
\renewcommand{\proofname}{\scshape Proof}
\renewcommand{\qedsymbol}{\openbox}

% ---- notation ----------------------------------------------------------------------
\newcommand{\Rl}{\mathbb{R}}
\newcommand{\E}{\mathbb{E}}
\newcommand{\one}{\mathbf{1}}
\newcommand{\Pol}{\Pi}                       % policy space
\newcommand{\cls}{\mathcal{F}}               % class of problems
\newcommand{\ax}[1]{\textup{\textsf{#1}}}    % cross-problem axiom
\newcommand{\cri}[1]{\textup{C#1}}           % within-problem criterion
\newcommand{\Cpr}{\ax{C}\ensuremath{_{\mathrm{pr}}}}
\newcommand{\Ccol}{\ax{C}\ensuremath{_{\mathrm{col}}}}
\newcommand{\Cref}{\ax{C}\ensuremath{_{\mathrm{ref}}}}
\newcommand{\Bpr}{\ax{B}\ensuremath{_{\mathrm{pr}}}}
\newcommand{\Bcol}{\ax{B}\ensuremath{_{\mathrm{col}}}}
\newcommand{\Mon}{\ax{Mon}}
\newcommand{\QJ}{\ax{Q}\ensuremath{_{J}}}
\newcommand{\dev}[3]{#1[#2{\leftarrow}#3]} % local deviation pi[h <- a]
\newcommand{\Reach}{\operatorname{Reach}}
\newcommand{\supp}{\operatorname{supp}}
\newcommand{\cont}{\mathfrak{c}}
\DeclareMathOperator*{\argmax}{arg\,max}
\DeclareMathOperator*{\argmin}{arg\,min}
\DeclareMathOperator{\first}{first}
\newcommand{\Prov}{\mdlgwhtsquare} % provability box

\hypersetup{pdftitle={Is There an Optimal Decision Theory?},pdfauthor={Un crack},
            pdfsubject={Possibility and impossibility theorems for decision theories: within-problem
                        criteria, cross-problem axioms, the semantics of menus, and logical counterfactuals}}

% ===================================================================================
\begin{document}

\thispagestyle{plain.scrheadings}
\begin{center}
  \vspace*{.2\baselineskip}
  {\LARGE Is There an Optimal Decision Theory?\par}
  \vspace{.7\baselineskip}
  {\large\itshape Within-Problem Conflicts, the Semantics of Menus,\\ and Counterfactual Alternatives\par}
  \vspace{1.4\baselineskip}
  {\scshape \MakeLowercase{\theauthor}\par}
  \vspace{.3\baselineskip}
  {\small October 2026\par}
  \vspace{1.6\baselineskip}
\end{center}

\begin{center}
\begin{minipage}{27pc}
  \small
  \noindent{\scshape abstract.}\enspace
  It is often suggested that no decision theory can be optimal on every fair decision problem,
  because the standards by which theories are judged conflict. This paper makes the claim precise
  for three successively stronger notions of a decision theory and determines when it holds. If a
  theory is any assignment of policies to finite problems whose structural representation supplies
  every counterfactual, ex-ante optimality is always attainable, and every inconsistency among
  ex-ante optimality, causal dominance, ratifiability, commitment stability, and an evidential
  criterion has a core of one or two criteria; the conflict between ex-ante optimality and dominance
  is possible exactly when the state law depends on the prescribed action. If a theory is a rule on
  causal matrices, whose rows are the belief states induced by the agent's policies and whose
  diagonal is ex-ante value, common affine invariance again yields only small cores, while
  invariance under row translations, with neutrality and contraction, forces a rule to choose every
  option on exactly the domains satisfying a reachability condition, and principal binary
  independence characterises ex-ante evaluation; both results fail if a smaller menu keeps the belief
  model of the larger one. If a theory is a computable procedure that reads the source code of its
  environment and must infer the dependence itself, a uniform optimiser exists under a designated
  substitution semantics and a termination promise; the negative results come from fixed copies of
  the agent, from the uncomputability of exact selection, and from the ways of recognising the agent
  in code. An axiomatisation of counterfactual operators shows that calibration together with an
  extensionality that identifies computations by their actual behaviour forces causal evaluation and
  is inconsistent with ratifiability, while an extensionality that compares computations under every
  substitution admits ex-ante evaluation satisfying every desideratum. Each impossibility concerns how
  the agent's alternatives are specified, not how to choose among them.
\end{minipage}
\end{center}


\vspace{.8\baselineskip}
{\small\setcounter{tocdepth}{1}\renewcommand{\contentsname}{\normalsize\scshape contents}\tableofcontents}

% -----------------------------------------------------------------------------------
\section{Introduction}\label{sec:intro}

Is there a decision theory that gets every decision problem right? For half a century the question
has been pursued case by case. A theory is proposed; a problem is constructed on which its verdict
seems wrong; the theory is amended or a rival proposed; and the rival meets its own
counterexample. The impression this history leaves is that no theory can be optimal everywhere,
because the standards by which theories are judged pull in different directions. This paper tries to
turn that impression into theorems. It asks, for precisely defined criteria of optimality and
precisely defined classes of problems, when a decision theory satisfying them exists and when it
does not, and it tries to close the question in the sense of determining every inconsistency among
the criteria considered and exhibiting a witness for every consistent combination.

\subsection{From dominance to Newcomb's problem}

The expected-utility theories of \textcite{ramsey1931}, \textcite{vonneumannmorgenstern1944}, and
\textcite{savage1954} model an agent's options as functions from states of the world to
consequences, with the states not depending on what the agent does. Under that assumption, the
principle of dominance---prefer an option that is better in every state---is a consequence of
maximising expected utility, and the two principles never conflict. \textcite{jeffrey1965} replaced
the act-independent states with a single algebra of propositions, valuing an act by the news it
would bring; \textcite{bolker1966,bolker1967} supplied the representation theorem. In Jeffrey's theory
the probability of a state may depend on the act, and dominance and expected utility can come apart.

\textcite{nozick1969} made the possibility vivid. A reliable predictor has placed a large sum in an
opaque box if and only if it predicted that the agent would take only that box; a transparent box
contains a small sum. Taking both boxes dominates, since the contents are fixed; taking only the
opaque box maximises expected utility conditional on the act. Nozick credited the problem to the
physicist William Newcomb and proposed restricting dominance to cases where the states are
probabilistically independent of the acts. \textcite{gibbardharper1978} gave the conflict its
canonical form by distinguishing two expected utilities: one computed with the probabilities of
counterfactual conditionals, which measures what an act would causally bring about, and Jeffrey's,
which measures what it indicates. Causal decision theory, as the first came to be called, takes both
boxes; evidential decision theory takes one. Gibbard and Harper also introduced Death in Damascus, in
which Death, having predicted the agent's destination, waits there: whichever city the agent
chooses, the other looks better. \textcite{lewis1979} observed that the prisoner's dilemma played
against a sufficiently similar partner is a Newcomb problem, and \textcite{lewis1981cdt} unified the
causal theories of Gibbard and Harper, \textcite{skyrms1980,skyrms1982}, and Sobel \autocite[see][]{sobel1994} by means of
dependency hypotheses. Against the charge that causal decision theorists end up poorer, Lewis
replied that Newcomb's problem rewards irrationality \autocite{lewis1981rich}. On the evidential side,
the tickle defence of \textcite{eells1982} argued that a self-aware agent's own beliefs and desires screen
off the correlations on which one-boxing relies, and \textcite{spohn1977} argued that deliberation
crowds out prediction, so that an agent cannot coherently assign probabilities to its own current
options. \textcite{joyce1999} gave causal decision theory a representation theorem, and by the end of
the century it had become the orthodox view \autocite{weirich2024sep}.

\subsection{Instability, ratifiability, and the counterexamples}

Death in Damascus exposed a second fault line. \textcite{jeffrey1981,jeffrey1983} proposed that a
rational choice should be \emph{ratifiable}: optimal relative to the beliefs the agent would have
upon making it. In Death in Damascus no pure option is ratifiable. \textcite{harper1986} proposed
maximising causal expected utility among ratifiable options and admitting mixed acts so that such
options exist; \textcite{richter1984}, \textcite{weirich1985}, and \textcite{rabinowicz1988} studied the
instability that makes every option look worse once chosen; and \textcite{skyrms1990dynamics}
modelled deliberation as a dynamical process whose rest points are deliberational equilibria,
connecting ratifiability with equilibrium in games \autocite{skyrms1990ratif}. \textcite{sobel1994}
collected much of this work.

\textcite{egan2007} reopened the debate with cases in which an act is evidence about the very causal
structure that determines its outcome---shooting as evidence of a lesion that would make one miss,
pressing a button that kills all psychopaths as evidence that one is a psychopath---and in which
causal decision theory seems to give the wrong verdict. The responses define the current landscape.
Deliberational theories recommend equilibria of a process of self-updating
\autocite{arntzenius2008,joyce2012,joyce2018}; benchmark theory assesses each option against a
state-dependent benchmark \autocite{wedgwood2013}; \textcite{spohn2012} argued that a properly
formulated causal theory one-boxes; \textcite{ahmed2012,ahmed2014book,ahmed2014dicing} revived the
evidential theory with cases in which causal reasoning fails; \textcite{harehedden2016} classified
decisions as self-reinforcing or self-frustrating and criticised the stability-based remedies;
\textcite{gallow2020} proposed a causal theory that manages the news an act carries; and
\textcite{gallow2026} argues for comparativism, on which instrumental value is fundamentally a matter
of pairwise comparison. The argument from wealth was reconsidered as well: \textcite{bales2018} and
\textcite{spencerwells2019} argued that it compares agents who face different problems, and
\textcite{wells2019} constructed a problem with equal opportunities in which causal reasoners
predictably end up richer. The collection edited by \textcite{ahmed2018ed} surveys the state of play.

\subsection{Policies, plans, and fairness}

A third line of work evaluates not acts but policies. Kuhn's analysis of extensive games introduced
perfect recall and the equivalence of mixed and behavioural strategies under it
\autocite{kuhn1953}. \textcite{hammond1988} derived expected utility from consequentialism and dynamic
consistency in decision trees; \textcite{mcclennen1990} defended resolute choice, the execution of an
ex-ante plan even when it is not optimal ex post; and \textcite{gauthier1986,gauthier1994} argued that
dispositions which are rational to adopt can make it rational to carry out assurances and threats.
\textcite{piccionerubinstein1997} separated planning-stage from action-stage optimality in the
absent-minded driver, and \textcite{aumannetal1997} showed that the planning-optimal strategy is
action-optimal although the converse fails; \textcite{arntzeniusetal2004} described sequences of
individually rational choices whose sum is disastrous and from which only binding escapes.

In artificial intelligence the policy view became a family of decision theories.
Updateless decision theory \autocite{dai2009} selects the input--output map that is optimal under
the agent's prior; timeless decision theory \autocite{yudkowsky2010} treats a decision as the output
of an abstract computation; and \textcite{soaresfallenstein2015} and \textcite{yudkowskysoares2017}
developed these ideas into functional decision theory, with a peer-reviewed presentation by
\textcite{levinsteinsoares2020}. Functional decision theory is claimed to do well on \emph{fair}
problems, those in which the outcome depends only on the agent's behaviour in the problem at hand;
its authors leave a formal criterion of fairness as future work and suggest that a general-purpose
theory may match or outperform all rivals on fair problems if one can be found.
\textcite{macaskill2019} objected with the Bomb, a problem in which the policy-optimal action is to take
a box known to contain a bomb, violating the principle that under certainty one should choose the
best outcome. \textcite{oesterheldconitzer2021} showed how to extract money from causal decision
theorists by predictive offers; \textcite{everittetal2015} formulated sequential, act-based and
policy-based versions of the causal and evidential theories; \textcite{belletal2021} showed that
value-based reinforcement learners in Newcomb-like environments converge, if at all, to ratifiable
policies, and that continuous environments always admit one; \textcite{halpern2020} gave conditions for
the existence of ratifiable mixed acts; \textcite{macdermottetal2023} characterised the theories with
mechanised causal graphs; and \textcite{oesterheldconitzer2024dese} and \textcite{cooperetal2024}
studied when self-locating beliefs make act-level causal or evidential reasoning agree with the
ex-ante optimal policy. \textcite{wolpertbenford2013} argued that the two Newcomb arguments presuppose
incompatible formalisations of how choice and prediction are related.

\subsection{Programs, proofs, and logical counterfactuals}

The functional theories treat a decision as the output of a computation, and so inherit a question
that the earlier theories could leave implicit: what would have happened had that computation
produced a different output? An idea of this kind is already in Hofstadter's \emph{superrationality},
on which agents who know that they reason identically should cooperate in a symmetric prisoner's
dilemma \autocite{hofstadter1983}. Game theory made programs the objects of choice.
\textcite{tennenholtz2004} introduced program equilibrium, in which players submit programs that can
read one another's code, and \textcite{barasz2014} showed that agents defined in provability logic
cooperate robustly with one another, while \textcite{critch2019} proved a resource-bounded version of
Löb's theorem underlying such results.

Within the functional programme, \textcite{soaresfallenstein2015} identified reasoning about
counterpossibles---conditionals whose antecedent is a non-actual output of a deterministic
computation---as the central open problem, and gave a proof-based formalisation that queries a
provability oracle. Such agents can act on spurious counterfactuals, provable only because the agent's own action is
provably different. The difficulty was raised in discussion of timeless decision theory
\autocite{drescher2009}, given an explicit Löbian form for proof-searching agents
\autocite{slepnev2010,slepnev2012}, and analysed by \textcite{bensontilsen2014}; it is now known as the
\enquote{five-and-ten} problem \autocite{fallenstein2013fiveandten,demskigarrabrant2019}. The
\enquote{chicken rule}, on which an agent takes any action it can prove it will not take, was
proposed to block such proofs \autocite{slepnev2011}. Fallenstein developed
modal versions of updateless decision theory, proved that for every provably extensional problem
some sound extension of arithmetic makes it optimal \autocite{fallenstein2014optimality,
fallenstein2014improving}, constructed problems on which every fixed modal theory fails
\autocite{fallenstein2014evil}, and proposed a partial semantics of \enquote{third-person
counterfactuals} \autocite{fallenstein2015third,fallenstein2015chicken}. Logical induction
\autocite{garrabrant2016li} gave a computable criterion for assigning probabilities to logical
sentences; its use for decisions requires exploration and meets obstacles of its own
\autocite{garrabrant2017obstacles}. \textcite{demskigarrabrant2019} placed these problems within a
wider account of agents embedded in the world they reason about, and \textcite{taylor2020}
proposed alternatives to logical counterfactuals altogether. In general reinforcement learning,
\textcite{leikehutter2015} showed that natural notions of optimality for universal agents are either
dependent on the prior or trivial. Two classical results of computability theory frame all of this:
Rice's theorem \autocite{rice1953}, that every non-trivial property of the functions that programs
compute is undecidable, and Kleene's second recursion theorem \autocite{kleene1938,moschovakis2010}, which
allows a program to refer to its own code.

In philosophy, the same conditionals are \emph{counterpossibles}. On the orthodox semantics of
\textcite{lewis1973counterfactuals} they are all vacuously true; non-vacuists evaluate them at impossible
worlds \autocite{nolan1997,brogaardsalerno2013}; and the exchange between \textcite{williamson2015,williamson2018}
and \textcite{bertoetal2018} concerns whether a semantics can treat them non-vacuously while respecting
the identifications that ordinary counterfactuals respect. \textcite{lewis1976} proved the classical
triviality theorem for a related identification, of conditional probabilities with probabilities of
conditionals.

\subsection{Impossibility results}

Four results are of the form this paper pursues. \textcite{skalse2021} proposed a general
counterexample to any decision theory: an opaque box that is empty exactly when the target theory
permits taking it. \textcite{fallenstein2014evil} proved the corresponding result in provability logic:
every fixed modal decision theory fails on some provably extensional problem. \textcite{briggs2010} proved that no completely general decision theory satisfies
both a Pareto condition across the evaluations of an agent's possible post-decision selves and a
self-sovereignty condition on which the comparison of two options depends only on how those two
options evaluate all options, modelling decision-theoretic paradoxes on voting paradoxes and Sen's
liberal paradox. \textcite{podgorski2022} proved that, for theories assigning each option an overall
value, independence of irrelevant alternatives forces the comparison of two options to depend only
on their values by their own lights, and so singles out the evidential theory; his tournament decision
theory escapes the result by producing no overall value.

These results draw on the theory of social choice, whose methods the present paper also uses.
\textcite{arrow1951} established that no aggregation rule satisfies a handful of individually
innocuous conditions on an unrestricted domain. The contraction condition on choice functions goes
back to \textcite{chernoff1954}; \textcite{arrow1959} showed that choice from all finite subsets is
rationalised by an ordering exactly when it satisfies the weak axiom of revealed preference; and
\textcite{sen1970,sen1971,sen1977,sen1993} named and analysed the conditions $\alpha$, $\beta$, and $\gamma$
and the limits of internal consistency of choice. Contraction consistency conflicts with majority
choice through the binary cycle, a fact central to the theory of tournament solutions
\autocite{moulin1986,laslier1997,brandtharrenstein2011,brandtbrillharrenstein2016}, and it continues to
generate impossibilities \autocite{blairetal1976,basumallikbhattacharya2026}.

\subsection{The question and the answer}

The thesis this paper set out to prove was that no decision theory is optimal on all fair decision
problems, the obstruction coming from mutually incompatible criteria of optimality rather than from
any demand for a general rule. Whether it is true depends on what a decision theory is required to
be, and the paper answers the question three times, for three successively stronger requirements.
Each strengthening moves the obstruction, and at each stage the obstruction turns out to lie in the
specification of the agent's alternatives rather than in the choice among them.

\emph{A theory as an assignment} (Part~\ref{part:single}). Section~\ref{sec:model} defines a class of
finite problems with prediction in which every counterfactual the criteria need is supplied by a
structural representation: for every policy, the law of the world's response is part of the problem.
A decision theory is any assignment of policies to problems, lookup tables included. The thesis is
then false: choosing an ex-ante maximiser in each problem is a decision theory that is ex-ante optimal
everywhere (Theorem~\ref{thm:universal}). The answer is cheap because the model hands the agent exactly
what causal, evidential, and functional theories disagree about. What remains true is shallow. Among
ex-ante optimality, structural causal dominance, ratifiability, commitment stability, and an evidential
criterion, every inconsistency has a core of one or two criteria (Theorems~\ref{thm:map15}
and~\ref{thm:map16}), the central pair being Nozick's. A conflict between ex-ante optimality and
dominance is possible exactly when the state law depends on the prescribed action
(Theorem~\ref{thm:characterisation}); it survives imperfect prediction and randomisation, with an exact
trade-off between the two kinds of regret (Proposition~\ref{prop:tradeoff}). Every within-problem
conflict reduces to the choice of the law under which an alternative is evaluated: the law induced by
the alternative, or the law induced by the actual choice.

\emph{A theory as a rule} (Part~\ref{part:across}). Requiring one rule across problems needs a common
description of them. Each problem is reduced to its causal matrix, whose rows are the belief states
induced by the agent's possible policies and whose diagonal is ex-ante value; every real square matrix
arises (Theorem~\ref{thm:realise}). With common affine invariance, the cross-problem axioms again
reduce to singleton and pairwise cores, one of which is Briggs's theorem (Theorem~\ref{thm:pairwise}).
With invariance under row translations---the matrix form of evaluating options against the baseline
of each belief state---a genuine four-axiom impossibility appears: neutrality, row-translation
invariance, and contraction force a rule to choose every option (Theorem~\ref{thm:circulant}), on
exactly those domains satisfying a reachability condition (Theorem~\ref{thm:DR}). Ex-ante evaluation,
in turn, is characterised by principal binary independence together with neutrality, affine
invariance, ordering, non-triviality, and monotonicity (Theorem~\ref{thm:selfesteem} and
Corollary~\ref{cor:three}), whose core is Podgorski's theorem. Both results depend on how a smaller
menu is modelled. If removing options also removes the beliefs that those options would induce, the
results hold; if the belief model is held fixed, row-baseline-invariant evaluation becomes menu
coherent, and contraction no longer forces ex-ante evaluation (Section~\ref{sec:menus} and
Theorem~\ref{thm:lambda}). The impossibilities concern the preservation of comparisons when the belief
perspectives associated with removed options disappear, not the preservation of choices when options
disappear.

\emph{A theory as a procedure} (Part~\ref{part:programs}). Finally, the supplied dependence is removed.
An agent is a computable function that receives the source code of its environment and must itself
determine what would happen under each of its outputs. With a designated substitution interface and a
termination promise, one algorithm is optimal everywhere (Theorem~\ref{thm:promiseopt}), the uniform
counterpart of Theorem~\ref{thm:universal}. The negative results separate into three kinds that have
nothing to do with counterfactual inference: queries about other problems
(Proposition~\ref{prop:foreign}); fixed copies of the agent, which defeat every agent that halts on all
inputs (Theorem~\ref{thm:fixedcopy}, Skalse's construction made well defined); and the uncomputability
of exact selection among computable expected values, which holds with no predictor at all
(Theorem~\ref{thm:exactsel}). The genuinely counterfactual question is how an agent should identify the
parts of a closed program that depend on its output. Literal, local, and global matching of the agent's
code each fail in a characteristic way (Theorems~\ref{thm:localmatch}--\ref{thm:haltred}), and
proof-based counterfactuals are spurious or undefined for ordinary programs (Theorem~\ref{thm:fiveten}
and Proposition~\ref{prop:guardedarith}). Axiomatising counterfactual operators gives a complete
consistency map (Theorem~\ref{thm:cfmap}). Calibration on independent problems, together with an
extensionality that identifies computations by their actual behaviour, forces causal evaluation
(Lemma~\ref{lem:projection}), and with ratifiability it forms a minimal inconsistent triple whose
witness is Death in Damascus (Theorem~\ref{thm:calibration}). Under an extensionality that compares
computations under every substitution, a single operator, ex-ante evaluation, satisfies every
desideratum (Corollary~\ref{cor:parametric}). The two evaluations of Part~\ref{part:single} thus
reappear as two notions of when computations are the same.

The conclusion of the paper unifies the three answers. \emph{No impossibility found here is an
impossibility of choosing well among specified alternatives. Each concerns how the alternatives are
specified: under which law an alternative is evaluated, which belief perspectives survive when options
are removed, and how a change in the agent's output is transported through the computations that
depend on it.}

Several results are known and are presented as such: the Newcomb conflict (Nozick), the absence of
ratifiable pure acts (Gibbard and Harper; Jeffrey), mixed existence under continuity (Halpern; Bell
and co-authors), the Pareto--self-sovereignty conflict (Briggs), the core of the binary
characterisation (Podgorski), the fixed-copy diagonal (Skalse; Fallenstein in modal form), the
spurious-proof mechanism behind five-and-ten, modal optimality (Fallenstein), and the recursion-theoretic
ingredients of the exact-selection theorem. The consistency maps, the characterisation of when
dominance and ex-ante optimality can conflict, the trade-off identity, the analysis of noisy
commitment, the circulant theorem and its exact domain criterion, the comparison of menu semantics,
the cardinality results on families of rules, the matching theorems, and the consistency map for
counterfactual operators were not found in the literature searched; their ingredients are classical,
and no claim of priority is made.

Part~\ref{part:single} (Sections~\ref{sec:model}--\ref{sec:robust}) treats single problems: the model,
the criteria and the possibility theorem, the witness problems, commitment and the within-problem
maps, and robustness. Part~\ref{part:across} (Sections~\ref{sec:matrix}--\ref{sec:mixed}) treats rules:
the causal matrix, the pairwise reduction, row translations and contraction, menu semantics, the
characterisation of ex-ante evaluation, problem types, and mixed policies. Part~\ref{part:programs}
(Sections~\ref{sec:programs}--\ref{sec:counterpossibles}) treats procedures: problems as programs,
computable optimisation and its limits, the inference of dependence from source, the axioms for
counterfactual operators, and the witness problems revisited together with the literature on
counterpossibles. Section~\ref{sec:conclusion} states the final thesis and the questions left open, and
Appendix~\ref{app:stipulations} lists every stipulation made.


\part{Single problems}\label{part:single}
% -----------------------------------------------------------------------------------
\section{Decision problems with prediction}\label{sec:model}

This section fixes the class of problems on which decision theories are compared. The aim is a
model rich enough to contain the problems that drive the debate (Newcomb's problem, the twin
prisoner's dilemma, Death in Damascus, the Bomb, counterfactual mugging) and precise enough that
every criterion to be considered has a determinate verdict on every problem. Each place where an
informal notion---a counterfactual, \enquote{the same computation}, \enquote{holding the state
fixed}---has to be settled by stipulation is marked as such; the stipulations are collected in
Appendix~\ref{app:stipulations}.

\subsection{Problems, policies, and structural plans}

\begin{definition}[Decision problem]\label{def:problem}
A \emph{decision problem} $D$ is a finite extensive-form tree with three kinds of non-terminal
nodes and real payoffs at its leaves:
\begin{enumerate}[label=(\roman*)]
  \item \emph{nature nodes}, with fixed probability distributions over their successors;
  \item \emph{prediction nodes}, nature nodes whose distribution may depend on the agent's policy,
        but only through the actions that the policy prescribes at specified information sets;
  \item \emph{agent nodes}, partitioned into a finite set $H_D$ of information sets, each $h\in H_D$
        carrying a finite nonempty set $A_h$ of actions.
\end{enumerate}
It is required that no information set occurs twice on any path of the tree.
\end{definition}

Perfect recall implies the last requirement, which is all that will be used. A \emph{deterministic
policy} for $D$ is a choice of action at every information set, and the policy space is
\[
  \Pol_D=\prod_{h\in H_D}A_h,
\]
a finite nonempty set. Mixed policies are deferred to Sections~\ref{sec:robust}
and~\ref{sec:mixed}.

The policy-dependent kernels of the prediction nodes determine the distribution of actual play,
but they do not by themselves settle the causal counterfactuals that several criteria require:
what would have happened had the agent acted otherwise, \emph{with the environment's responses
held fixed}. To make these counterfactuals explicit, the tree is supplemented with a structural
description of the environment.

\begin{definition}[Structural representation]\label{def:structural}
A \emph{structural representation} of $D$ consists of
\begin{enumerate}[label=(\roman*)]
  \item a finite nonempty set $S_D$ of \emph{environment plans}, each of which specifies the response
        of every nature and prediction node, including nodes on branches not taken and responses that
        depend on earlier physical actions;
  \item for every policy $\pi\in\Pol_D$, a probability distribution $\mu_\pi\in\Delta(S_D)$;
  \item the payoff $U_D(s,\pi)$ obtained by traversing the tree with plan $s$ and policy $\pi$.
\end{enumerate}
The \emph{ex-ante value} of $\pi$ is
\[
  V_D(\pi)=\sum_{s\in S_D}\mu_\pi(s)\,U_D(s,\pi).
\]
\end{definition}

The representation separates two kinds of dependence. \emph{Ordinary causal dependence} of the
environment on earlier actions is carried by the responses written into each plan $s$.
\emph{Predictive dependence} on the policy is carried by the map $\pi\mapsto\mu_\pi$. A tree with
policy-dependent kernels can always be given a structural representation by supplying potential
outcomes for its nature and prediction nodes; how the outcomes on alternative branches are coupled
is an additional modelling choice, which matters for the causal criteria and is fixed by the
representation. Conversely, every structural representation is implemented by a tree whose root
is a prediction node selecting $s$ according to $\mu_\pi$, followed by deterministic traversal.

\subsection{Well-definedness and fairness}

Two conditions on problems are distinguished.

\begin{definition}[Well-definedness]\label{def:welldefined}
$D$ is \emph{well defined} if $\mu_\pi$ is a probability distribution on $S_D$ for \emph{every}
$\pi\in\Pol_D$.
\end{definition}

\begin{definition}[Fairness]\label{def:fair}
$D$ is \emph{fair} if $S_D$, $U_D$, and the entire map $\pi\mapsto\mu_\pi$ are fixed features of $D$:
the evaluation of an agent may depend on the actions its policy prescribes, but not on the
identity, name, or source code of the decision theory that produced the policy. Fixed algorithms
may occur as constants of the environment (\enquote{the predictor simulates the algorithm $g_0$}),
provided $g_0$ does not vary with the agent being evaluated.
\end{definition}

Fairness excludes the trivial diagonal problem that pays an agent exactly when its action differs
from that of a given theory $T$; such a problem defeats every $T$ and shows nothing about decision
theory. Well-definedness is a separate requirement. A verbally described \enquote{perfect
predictor} may satisfy an action-based fairness condition while its promised behaviour is
logically inconsistent for some policy; Section~\ref{sec:bomb} gives an example. The presence of a
fixed algorithm in the environment does not by itself make a self-referential specification well
defined (Remark~\ref{rem:skalse}). Fixed algorithms are harmless here because a problem is a finite
object whose data include their outputs; when problems are given as source code and the agent must
compute those outputs, they are not (Section~\ref{sec:computation}).

Write $\cls^*$ for the class of well-defined, fair finite problems with a fixed structural
representation, closed under the commitment construction of Section~\ref{sec:commit}.

\begin{definition}[Decision theory]\label{def:theory}
A \emph{decision theory} on a class $\cls$ of problems assigns to every $D\in\cls$ a policy
$T_D\in\Pol_D$. No uniformity, computability, or generality is required: every lookup table is a
decision theory.
\end{definition}

Totality is part of the definition. A proposed criterion that recommends nothing on some problem
therefore fails to define a decision theory on any class containing that problem; it does not
define a theory that is silent there. The permissive definition is the weakest sense of
\enquote{decision theory}, and it is the first of three. Part~\ref{part:across} requires a theory to be a
single rule on a common reduced form of problems, and Part~\ref{part:programs} requires it to be a
computable procedure that receives a description of the problem and must itself work out what the
structural representation supplies here.

\subsection{Reach and local causal deviation}

For a policy $\pi$ and an information set $h$, let
\[
  R_h(\pi)=\{s\in S_D: h\text{ is reached in the traversal with }(s,\pi)\},
  \qquad
  r_\pi(h)=\mu_\pi\bigl(R_h(\pi)\bigr).
\]
The number $r_\pi(h)$ is the actual probability that $h$ is reached under $\pi$; $h$ is
\emph{positively reached} under $\pi$ if $r_\pi(h)>0$. For $a\in A_h$, let $\dev{\pi}{h}{a}$ denote the
policy that agrees with $\pi$ everywhere except at $h$, where it prescribes $a$.

\begin{lemma}\label{lem:local}
For every $\pi$, $h$, and $a\in A_h$:
\begin{enumerate}[label=(\roman*)]
  \item $R_h(\dev{\pi}{h}{a})=R_h(\pi)$;
  \item if $s\notin R_h(\pi)$, then $U_D(s,\dev{\pi}{h}{a})=U_D(s,\pi)$.
\end{enumerate}
\end{lemma}

\begin{proof}
With the plan $s$ fixed, the path through the tree is determined by the prescriptions at the
information sets it visits. Up to its first visit to $h$, the path under $(s,\dev{\pi}{h}{a})$ uses
only prescriptions at information sets other than $h$, which are unchanged; so it coincides with the
path under $(s,\pi)$, and $h$ is reached under one exactly when it is reached under the other. If
$h$ is not reached, the two paths coincide entirely. Since no information set recurs on a path, a
visit to $h$ is the first and only one.
\end{proof}

The \emph{causal value} of $a$ at a positively reached $h$, with beliefs induced by $\pi$, is
\[
  Q_\pi(h,a)=\frac{1}{r_\pi(h)}\sum_{s\in R_h(\pi)}\mu_\pi(s)\,U_D\bigl(s,\dev{\pi}{h}{a}\bigr).
\]
This fixes the causal counterfactual completely: the distribution over environmental responses is
the one the actual policy induces, the action at $h$ is changed, and the policy's other
prescriptions are kept. The contrast with ex-ante evaluation lies in a single subscript:
$Q_\pi(h,a)$ averages with $\mu_\pi$, whereas $V_D(\dev{\pi}{h}{a})$ averages with
$\mu_{\dev{\pi}{h}{a}}$. Every within-problem conflict studied below is a consequence of that
difference.

% -----------------------------------------------------------------------------------
\section{Optimality criteria and the possibility of ex-ante optimality}\label{sec:criteria}

\subsection{Within-problem criteria}

The criteria below are the precise forms of the principles that the main decision theories are
usually taken to satisfy. Each is a condition on the policy $T_D$ that a theory $T$ selects in a
problem $D$; the commitment criterion \cri{4} and the evidential criterion \cri{6} are introduced in
Section~\ref{sec:commit}, once the constructions they require are available.

\begin{definition}[Ex-ante optimality, \cri{1}]
$T$ satisfies \cri{1} on $D$ if \[ T_D\in\argmax_{\pi\in\Pol_D}V_D(\pi). \]
\end{definition}

\begin{definition}[Structural causal dominance, \cri{2}]
$T$ satisfies \cri{2} on $D$ if, at every information set $h$ positively reached under $\pi=T_D$,
there is no $a\in A_h$ such that
\[
  U_D\bigl(s,\dev{\pi}{h}{a}\bigr)>U_D\bigl(s,\dev{\pi}{h}{\pi(h)}\bigr)
  \quad\text{for every }s\in R_h(\pi).
\]
\end{definition}

The quantifier ranges over \emph{structural} states: all plans that reach $h$, including those to
which $\mu_\pi$ assigns probability zero. This is a deliberate resolution of an ambiguity in
\enquote{every state compatible with $h$}, and it matters. If compatibility meant only the states in
the support of $\mu_\pi$, then in Death in Damascus with an infallible predictor (Section~\ref{sec:did})
the single supported state places Death in the chosen city, so switching would strictly improve the
payoff in every supported state and \cri{2} itself would be unsatisfiable. Likewise, on trees where
prediction nodes follow $h$, \enquote{hold fixed the outcomes of the nodes preceding $h$} leaves open
how the later prediction nodes are to be evaluated; freezing their policy dependence while
averaging over their outcomes can make dominance collapse into an unsatisfiable best-response
requirement. The structural plan fixes both matters at once.

\begin{definition}[Ratifiability, \cri{3}]
$T$ satisfies \cri{3} on $D$ if, at every information set $h$ positively reached under $\pi=T_D$,
\[
  \pi(h)\in\argmax_{a\in A_h}Q_\pi(h,a).
\]
\end{definition}

\begin{definition}[Existence, \cri{5}]
$T$ satisfies \cri{5} on $D$ if its own criterion is satisfiable there, that is, if $T_D$ is defined.
\end{definition}

Since decision theories are total by Definition~\ref{def:theory}, \cri{5} holds for every theory
and imposes no restriction; it records the requirement that a criterion be satisfiable on every
problem of the class. Were partial theories admitted, totality would become an independent axiom
and the consistency maps of Section~\ref{sec:commit} would change.

\begin{proposition}\label{prop:c3c2}
\cri{3} implies \cri{2}.
\end{proposition}

\begin{proof}
Suppose $h$ is positively reached under $\pi=T_D$ and some $a$ satisfies
$U_D(s,\dev{\pi}{h}{a})>U_D(s,\pi)$ for every $s\in R_h(\pi)$. Since $r_\pi(h)>0$, the conditional
distribution $\mu_\pi(\cdot\mid R_h(\pi))$ is a probability measure carried by $R_h(\pi)$, and a strict
pointwise inequality on $R_h(\pi)$ gives $Q_\pi(h,a)>Q_\pi(h,\pi(h))$. Thus $\pi(h)$ is not a maximiser
of $Q_\pi(h,\cdot)$, and \cri{3} fails.
\end{proof}

\subsection{Ex-ante optimality is always attainable}

The opening question admits a direct answer once \enquote{decision theory} is taken in the
permissive sense of Definition~\ref{def:theory}.

\begin{theorem}[Universal ex-ante selection]\label{thm:universal}
On every class of well-defined finite problems there is a decision theory satisfying \cri{1} on
every problem of the class.
\end{theorem}

\begin{proof}
Fix a tie-breaking order on each $\Pol_D$ and set
$T^{\mathrm{ex}}_D=\first\argmax_{\pi\in\Pol_D}V_D(\pi)$. Each $\Pol_D$ is finite and nonempty, and
$V_D$ is real valued on it because $D$ is well defined, so the maximum exists.
\end{proof}

\begin{remark}[Self-referential problems]\label{rem:skalse}
\textcite{skalse2021} proposes, for a fixed target theory $X$, a choice between a transparent box
containing $1000$ and an opaque box whose content is $0$ if $X$ permits choosing it in that problem and
$2000$ otherwise, and argues that another theory can outperform $X$. The target theory is an
environmental constant, so this is not discrimination by the identity of the evaluated agent.
Translated into the present model, however, the construction requires a fixed point: for an ex-ante
optimiser, an empty opaque box makes it choose the transparent box, which requires the opaque box to
be full, and conversely. Neither content yields a well-defined instance. In the same way, the
diagonal family $D_b$ with $u(a)=\one\{a\neq b\}$, $b\in\{0,1\}$, would need $b=T(D_b)$, but an ex-ante
optimiser has $T(D_0)=1$ and $T(D_1)=0$. The schema does not refute Theorem~\ref{thm:universal} on
well-defined problems; its additional self-referential realisability assumption is where it leaves
that class. \textcite{lewis1981rich} analysed a related failure: an announced reward for supposedly
irrational choice can render the proposed scenario inconsistent. The verdict depends on the problem
being a finite object whose data include the constant's output. When problems are given as programs,
so that the agent must compute that output itself, the same construction is well defined and defeats
every agent that halts on all inputs (Theorem~\ref{thm:fixedcopy}).
\end{remark}

The unqualified thesis that no decision theory is optimal on all fair problems is therefore false.
What can be true is that particular \emph{combinations} of criteria are jointly unsatisfiable, and
the remaining sections determine exactly which. A first observation shows how little such a result
can depend on generality or uniformity. Call a criterion $C$ \emph{local} if whether a theory $T$
satisfies $C$ on $D$ depends only on $T_D$; \cri{1}, \cri{2}, \cri{3}, and \cri{5} are local, while
\cri{4} is not. For a local criterion $C_i$ let $A_i(D)\subseteq\Pol_D$ be the set of policies that
satisfy it on $D$.

\begin{lemma}[Pointwise reduction]\label{lem:pointwise}
For local criteria $C_1,\dots,C_k$ and a class $\cls$,
\[
  \exists T\;\forall D\in\cls\;\bigwedge_{i}C_i(T,D)
  \quad\Longleftrightarrow\quad
  \forall D\in\cls\;\;\bigcap_{i}A_i(D)\neq\varnothing.
\]
\end{lemma}

\begin{proof}
From left to right, $T_D$ lies in the intersection. From right to left, let $T_D$ be the first
element of the intersection under a fixed order.
\end{proof}

Because lookup tables are admitted, an impossibility for local criteria is never more than the
exhibition of a single problem on which their acceptable sets are disjoint. No difficulty about
general rules, computability, or self-reference is involved; those enter only when a theory must be
a single rule on a common description of problems (Part~\ref{part:across}) or a computable procedure
that reads descriptions (Part~\ref{part:programs}). This is the precise sense in which the
within-problem impossibilities of the next two sections are shallow; non-local criteria, and the
cross-problem axioms of Sections~\ref{sec:matrix}--\ref{sec:exante}, are where deeper structure
can appear.

% -----------------------------------------------------------------------------------
\section{The witness problems}\label{sec:witnesses}

This section works through the problems that separate the criteria. Payoffs are scaled for
arithmetic convenience; nothing depends on the particular numbers beyond the inequalities stated.

\subsection{Newcomb's problem}\label{sec:newcomb}

Let $O$ denote taking only the opaque box and $B$ taking both boxes. A predictor of symmetric error
$\varepsilon\in[0,1]$ forecasts the agent's prescription at its single information set, which is
reached with probability one; the opaque box is full if and only if the forecast is $O$. Payoffs are

\begin{center}
\small
\begin{tabular}{@{}lcc@{}}
  \toprule
  & opaque box full & opaque box empty \\
  \midrule
  $O$ & 2 & 0 \\
  $B$ & 3 & 1 \\
  \bottomrule
\end{tabular}
\end{center}
with $\Pr(\text{full}\mid O)=1-\varepsilon$ and $\Pr(\text{full}\mid B)=\varepsilon$. The structural
states are the two contents of the opaque box.

\begin{proposition}[\cite{nozick1969}]\label{prop:newcomb}
In Newcomb's problem $V(O)=2(1-\varepsilon)$ and $V(B)=1+2\varepsilon$, so $V(O)-V(B)=1-4\varepsilon$
and \cri{1} uniquely requires $O$ whenever $0\le\varepsilon<\tfrac14$. Since $U(s,B)-U(s,O)=1$ in both
structural states, \cri{2} and \cri{3} uniquely require $B$ for every $\varepsilon$. Hence no theory
satisfies \cri{1} and \cri{2} on any class containing Newcomb's problem with $\varepsilon<\tfrac14$.
\end{proposition}

With opaque prize $m$ and transparent bonus $c$, the same conflict holds exactly when
$m(1-2\varepsilon)>c$. Any symmetric predictor better than chance therefore produces it for suitable
stakes. The conflict is Nozick's original one between maximising expected utility and dominance;
the content of Proposition~\ref{prop:newcomb} is only that the formal criteria reproduce it.

The proposition shows that a conflict \emph{can} arise. The next result characterises exactly when
it can, and shows that neither accurate prediction nor large stakes is needed.

\begin{theorem}[Policy-dependent state laws characterise the conflict]\label{thm:characterisation}
Let $A$ be a finite set of actions with $|A|\ge2$, $S$ a finite set of states, and $\mu_a\in\Delta(S)$
the state law induced by prescribing $a$ at a single always-reached choice. The following are
equivalent:
\begin{enumerate}[label=(\roman*)]
  \item for every payoff function $u\colon A\times S\to[0,1]$, some ex-ante maximising action is not
        strictly dominated state by state;
  \item the laws $\mu_a$, $a\in A$, are all equal.
\end{enumerate}
\end{theorem}

\begin{proof}
(ii)$\Rightarrow$(i). Under a common law $\mu$, if $a$ strictly dominates $b$ in every state then
$\E_\mu u(a,\cdot)>\E_\mu u(b,\cdot)$, so no ex-ante maximiser is strictly dominated; a maximiser
exists because $A$ is finite.

(i)$\Rightarrow$(ii). Suppose $\mu_a\neq\mu_b$. There is an event $E\subseteq S$ on which they differ;
exchanging the names of $a$ and $b$ if necessary,
$\delta=\mu_b(E)-\mu_a(E)>0$. Choose $0<c<\delta/(1+\delta)$ and set
\[
  u(b,s)=(1-c)\,\one_E(s),\qquad u(a,s)=(1-c)\,\one_E(s)+c,
\]
with payoff $0$ for every other action. All payoffs lie in $[0,1]$, and $a$ strictly dominates $b$ by
$c$ in every state. Yet
\[
  V(b)-V(a)=(1-c)\mu_b(E)-(1-c)\mu_a(E)-c=(1-c)\delta-c>0,
\]
and $V(a)\ge c>0$ exceeds the value $0$ of every other action. So $b$ is the unique ex-ante
maximiser, and it is strictly dominated.
\end{proof}

Any non-constant dependence of the preceding state law on the prescription thus suffices for some
bounded payoff assignment to separate \cri{1} from \cri{2}; for fixed payoffs, the size of the
dependence determines whether the conflict occurs.

\subsection{The twin prisoner's dilemma}\label{sec:twin}

Two copies of the same policy play a prisoner's dilemma with temptation $3$, reward $2$, punishment
$1$, and sucker's payoff $0$. The twin is represented as a prediction node that copies the agent's
prescription at its single information set; it is \emph{not} a second, independently assignable
coordinate of $\Pol_D$, which would remove the stipulated identity of the two decisions. Then
$V(\text{cooperate})=2$ and $V(\text{defect})=1$, so \cri{1} requires cooperation, while with the
twin's action held fixed, defection gains one unit whether the twin cooperates or defects, so
\cri{2} and \cri{3} require defection. The connection between this problem and Newcomb's is due to
\textcite{lewis1979}. What the tree does not supply is the identity of the two decisions:
\enquote{the same computation} enters only as a stipulated dependence of the prediction node on
the agent's prescription. Sections~\ref{sec:inference} and~\ref{sec:operators} ask whether an agent
given the twin's source code could infer that dependence, and Section~\ref{sec:counterpossibles}
records how each candidate answer decides this problem.

\subsection{Death in Damascus}\label{sec:did}

Death predicts which of two cities the agent will choose and waits there
\autocite{gibbardharper1978}. Let actions and states both be $\{1,2\}$, with
$U(s,a)=\one\{a\neq s\}$ and $\mu_a=\delta_a$ (infallible prediction).

\begin{proposition}\label{prop:did}
In Death in Damascus no deterministic policy satisfies \cri{3}, whereas both deterministic policies
satisfy \cri{2}. The failure of \cri{3} persists for every symmetric error $\varepsilon<\tfrac12$.
\end{proposition}

\begin{proof}
Under the policy $a$, the induced state is $a$; staying has causal value $0$ and switching has causal
value $1$. With error $\varepsilon$, staying has causal value $\varepsilon$ and switching $1-\varepsilon$,
and $\varepsilon<1-\varepsilon$. For \cri{2}, neither city is better than the other in \emph{both}
structural states, so no action is dominated.
\end{proof}

Consequently \cri{3}, read as a requirement on deterministic policies, is individually
unsatisfiable on any class containing this problem: it fails to define a decision theory at all.
The absence of ratifiable pure acts in this problem is well known from the ratifiability
literature \autocite{skyrms1990ratif}; Section~\ref{sec:mixedexist} records how far randomisation
repairs it.

\subsection{The Bomb}\label{sec:bomb}

\textcite{macaskill2019} describes a predictor of error $10^{-24}$ that places a bomb in the left
box if and only if it predicts that the agent will take the right box, and leaves a note saying so.
The agent sees the note and the bomb. Taking the left box is free; taking the right box costs $1$;
burning to death costs $M=10^{9}$ in the scaling used here.

Once the agent observes whether a bomb is present, there are \emph{two} information sets,
$h_{\mathrm{bomb}}$ and $h_{\mathrm{clear}}$, and four policies
$(a_{\mathrm{bomb}},a_{\mathrm{clear}})\in\{L,R\}^2$. This has a consequence for well-definedness.

\begin{proposition}[The Bomb as stated is not well defined]\label{prop:bombill}
No predictor whose forecast determines the observation (\enquote{predicts $R$: bomb; predicts $L$:
no bomb}) and whose target is the agent's \emph{eventual} action can have an error bound below one
uniformly over the four policies.
\end{proposition}

\begin{proof}
Under the policy $(L,R)$, a forecast of $R$ produces a bomb and hence the action $L$; a forecast of $L$
produces no bomb and hence the action $R$. The forecast is wrong in every case.
\end{proof}

A coherent version is obtained by letting the predictor target the fixed coordinate
$\pi(h_{\mathrm{bomb}})$, whether or not $h_{\mathrm{bomb}}$ is eventually reached:
\[
  \Pr\bigl(\text{bomb}\mid a_{\mathrm{bomb}}=L\bigr)=\varepsilon,\qquad
  \Pr\bigl(\text{bomb}\mid a_{\mathrm{bomb}}=R\bigr)=1-\varepsilon.
\]

\begin{proposition}\label{prop:bomb}
In the repaired Bomb with death cost $M>1$, the ex-ante values are
\begin{center}
\small
\begin{tabular}{@{}lc@{}}
  \toprule
  policy $(a_{\mathrm{bomb}},a_{\mathrm{clear}})$ & $V$ \\
  \midrule
  $(L,L)$ & $-M\varepsilon$ \\
  $(L,R)$ & $-M\varepsilon-(1-\varepsilon)$ \\
  $(R,L)$ & $-(1-\varepsilon)$ \\
  $(R,R)$ & $-1$ \\
  \bottomrule
\end{tabular}
\end{center}
For $0<\varepsilon<1/(M+1)$, \cri{1} uniquely selects $(L,L)$, while \cri{2} and \cri{3} uniquely
select $(R,L)$.
\end{proposition}

\begin{proof}
The values follow by conditioning on the presence of the bomb. $(L,L)$ beats $(R,L)$ exactly when
$M\varepsilon<1-\varepsilon$, that is $\varepsilon<1/(M+1)$; it beats $(L,R)$ and $(R,R)$ a fortiori. For
$\varepsilon>0$ both information sets are positively reached under every policy. At
$h_{\mathrm{bomb}}$, every structural state reaching it contains the bomb, and $R$ yields $-1>-M$;
at $h_{\mathrm{clear}}$, $L$ yields $0>-1$. Hence $(R,L)$ is the only policy satisfying \cri{2},
and, since it is causally optimal at both information sets, it satisfies \cri{3}.
\end{proof}

MacAskill's numbers satisfy the inequality. Two qualifications are essential. At $\varepsilon=0$,
$h_{\mathrm{bomb}}$ is unreached under $(L,L)$, and the reached-only criteria \cri{2} and \cri{3} do
not oppose that policy. And if the problem is instead taken to be the already-conditioned decision
\enquote{a bomb is visibly present}, its two options have values $-M$ and $-1$, and \cri{1} itself
selects $R$. The visible-bomb challenge is MacAskill's; the complete contingent-policy kernel above
is the additional assumption under which it becomes a theorem.

\subsection{Counterfactual mugging}\label{sec:mugging}

A fair coin is tossed. On heads the agent receives $r$ if and only if an infallible predictor
forecasts that the agent's prescription on tails is to pay; on tails the agent is asked to pay $c$,
with $r>c>0$ \autocite{nesov2009}. There is one information set, reached on tails. Then
\[
  V(\text{Pay})=\tfrac12(r-c)>0=V(\text{Refuse}),
\]
so \cri{1} requires paying, while at the tails information set paying costs $c$ in every structural
state reaching it, so \cri{2} and \cri{3} require refusing. The problem separates ex-ante from
ex-post evaluation by a prediction about an information set that is unreached on half the paths; it
is used in Section~\ref{sec:evidential} to locate evidential criteria and in
Section~\ref{sec:targets} to test restrictions on what predictors may target.

% -----------------------------------------------------------------------------------
\section{Commitment, evidence, and the within-problem map}\label{sec:commit}

\subsection{The commitment extension}

Reflective stability asks what a theory would do if, before any prediction is made, it could choose
which policy it will run. To state this as a criterion, the choice of a policy must itself be a
decision problem.

\begin{definition}[Commitment extension]\label{def:commit}
Let $D$ be a problem and $\lambda$ a \emph{coupling} of the laws $(\mu_q)_{q\in\Pol_D}$, that is, a
probability distribution on $\prod_{q\in\Pol_D}\supp\mu_q$ whose marginal on the $q$-th coordinate is
$\mu_q$. The \emph{commitment extension} $K_\lambda(D)$ is the problem in which
\begin{enumerate}[label=(\roman*)]
  \item the agent, at a single information set, chooses a policy $q\in\Pol_D$;
  \item this choice irrevocably installs $q$;
  \item the environment draws $\omega=(s_q)_{q\in\Pol_D}$ from $\lambda$, independently of the choice, and
        the agent receives $U_{K}(\omega,q)=U_D(s_q,q)$.
\end{enumerate}
Unless stated otherwise the structural states of $K_\lambda(D)$ are
$S_{K}=\prod_{q}\supp\mu_q$, and $K(D)$ denotes the product coupling
$\lambda=\bigotimes_q\mu_q$. The constructor is tagged and well founded, so that $K(D)$,
$K(K(D))$, \dots{} have unambiguous decompositions, and $\cls^*$ is closed under it.
\end{definition}

The extension models a \emph{causal} choice of policy made before prediction: changing the
commitment changes the environment's subsequent response, because the coordinate $s_q$ that is
paid out depends on $q$, while the law $\lambda$ of the whole vector does not.

\begin{proposition}\label{prop:commitvalue}
For every coupling $\lambda$ and every $q\in\Pol_D$, $V_{K_\lambda(D)}(q)=V_D(q)$.
\end{proposition}

\begin{proof}
$V_{K_\lambda(D)}(q)=\E_\lambda\bigl[U_D(s_q,q)\bigr]=\E_{\mu_q}\bigl[U_D(s,q)\bigr]=V_D(q)$, since the
$q$-th marginal of $\lambda$ is $\mu_q$.
\end{proof}

\begin{definition}[Commitment stability, \cri{4}]
$T$ satisfies \cri{4} on $D$ if the commitment $q=T_{K(D)}$ that $T$ chooses in $K(D)$ satisfies
$q(h)=T_D(h)$ at every information set $h$ positively reached under $T_D$.
\end{definition}

Unlike the other criteria, \cri{4} relates a theory's behaviour on two different problems. It is
the first cross-problem condition in this paper, and the only one stated without passing to the
reduced form of Section~\ref{sec:matrix}.

\subsection{The map of \texorpdfstring{\cri{1}--\cri{5}}{C1-C5}}

The consistent combinations need explicit total witnesses, not merely examples on a few problems.

\begin{proposition}[A coherent ex-ante theory]\label{prop:coherent}
There is a decision theory on $\cls^*$ satisfying \cri{1} and \cri{4}.
\end{proposition}

\begin{proof}
On problems that are not commitment extensions, let $T_D=T^{\mathrm{ex}}_D$. Recursively, in $K(D)$
let $T$ commit to the policy $T_D$ already selected for $D$. By Proposition~\ref{prop:commitvalue},
$T_D$ maximises $V_{K(D)}$, so \cri{1} holds on $K(D)$; \cri{4} holds with exact agreement.
\end{proof}

\cri{1} alone does not imply \cri{4}: on a problem in which both actions yield zero, an ex-ante
maximising lookup table may choose different tied actions in $D$ and in $K(D)$.

\begin{proposition}[A dominance-respecting theory]\label{prop:nu}
For each $D$ fix a full-support distribution $\nu_D$ on $S_D$, for instance the uniform one, and let
\[
  T^{\nu}_D=\first\argmax_{\pi\in\Pol_D}\;\sum_{s\in S_D}\nu_D(s)\,U_D(s,\pi).
\]
Then $T^{\nu}$ satisfies \cri{2} on every problem.
\end{proposition}

\begin{proof}
Let $\pi=T^\nu_D$ and suppose that at an information set $h$ positively reached under $\pi$ some $a$
satisfies $U_D(s,\dev{\pi}{h}{a})>U_D(s,\pi)$ for all $s\in R_h(\pi)$. By Lemma~\ref{lem:local}(ii)
the payoff is unchanged on $S_D\setminus R_h(\pi)$. The set $R_h(\pi)$ is nonempty because
$r_\pi(h)>0$, and $\nu_D$ gives it positive weight, so $\dev{\pi}{h}{a}$ strictly improves the
objective that $\pi$ maximises, a contradiction.
\end{proof}

$T^\nu$ is an existence witness, not a recommended theory: its fixed weights ignore the
predictive dependence altogether.

\begin{proposition}\label{prop:c2c4}
No theory satisfies \cri{2} and \cri{4} on a class containing Newcomb's problem with $\varepsilon=0$
and its commitment extension.
\end{proposition}

\begin{proof}
In $D$, \cri{2} requires $B$ (Proposition~\ref{prop:newcomb}). At $\varepsilon=0$ every support is a
singleton, so $K(D)$ has the single structural state (opaque box full under $O$, empty under $B$),
and committing to $O$ yields $2$ while committing to $B$ yields $1$. Committing to $O$ therefore
strictly dominates, and \cri{2} in $K(D)$ requires it; \cri{4} requires the commitment to agree with
$B$ at the always-reached information set.
\end{proof}

\begin{theorem}[Within-problem map]\label{thm:map15}
On $\cls^*$ the inclusion-minimal inconsistent subsets of $\{\cri{1},\dots,\cri{5}\}$ are exactly
\[
  \{\cri{3}\},\qquad\{\cri{1},\cri{2}\},\qquad\{\cri{2},\cri{4}\},
\]
and the maximal consistent subsets are exactly $\{\cri{1},\cri{4},\cri{5}\}$, witnessed by
Proposition~\ref{prop:coherent}, and $\{\cri{2},\cri{5}\}$, witnessed by
Proposition~\ref{prop:nu}.
\end{theorem}

\begin{proof}
The three sets are inconsistent by Propositions~\ref{prop:did}, \ref{prop:newcomb}, and
\ref{prop:c2c4}; their proper subsets are consistent by the two witnesses (\cri{5} holds for every
theory). A consistent set cannot contain \cri{3}. If it contains \cri{2}, it contains neither \cri{1}
nor \cri{4}, so it lies in $\{\cri{2},\cri{5}\}$. Otherwise its members lie in
$\{\cri{1},\cri{4},\cri{5}\}$.
\end{proof}

Among \cri{1}--\cri{5}, then, there is no minimally inconsistent family of three or more
individually satisfiable criteria. The requested minimality is attained only by the pairs
$\{\cri{1},\cri{2}\}$ and $\{\cri{2},\cri{4}\}$, whose singletons have explicit total witnesses, and
the first is a direct pointwise disagreement in the sense of Lemma~\ref{lem:pointwise}.

\subsection{An evidential criterion}\label{sec:evidential}

An evidential criterion requires a joint law of acts and states. Under a deterministic policy,
conditioning on an alternative action is ordinarily conditioning on a null event, so the expression
$\E[\text{payoff}\mid h,\text{agent takes }a]$ cannot be read off the actual distribution of play.
The following completion makes it determinate.

\begin{definition}[Reference-policy evidential criterion, \cri{6}$^\kappa$]
For each $D$ fix a full-support distribution $\kappa_D$ on $\Pol_D$ and define the joint law
$J_D(\Pi=\pi,S=s)=\kappa_D(\pi)\mu_\pi(s)$, with $Y=U_D(S,\Pi)$. For $h\in H_D$ and $a\in A_h$ let
\[
  e_h(a)=\E_{J_D}\bigl[Y\mid h\text{ reached},\ \Pi(h)=a\bigr],
\]
with $e_h(a)=-\infty$ when the conditioning event is null. $T$ satisfies \cri{6}$^\kappa$ on $D$ if
$T_D(h)\in\argmax_a e_h(a)$ at every $h$ positively reached under $T_D$.
\end{definition}

This is a specified evidential criterion, not an identification with every version of evidential
decision theory; in sequential problems it evaluates continuation behaviour with the reference
distribution over policies. It is individually satisfiable: choose an $e_h$-maximiser at every $h$
with $J_D(h\text{ reached})>0$ and anything elsewhere. If $h$ is positively reached under $T_D$, then
$J_D(h\text{ reached})\ge\kappa_D(T_D)\,r_{T_D}(h)>0$, so the constraint is met, and some action has
a finite evidential value there.

\begin{proposition}\label{prop:c6}
\begin{enumerate}[label=(\roman*)]
  \item In Newcomb's problem $e(O)=V(O)$ and $e(B)=V(B)$, so \cri{6}$^\kappa$ requires $O$ for
        $\varepsilon<\tfrac14$ and conflicts with \cri{2}.
  \item In counterfactual mugging \cri{6}$^\kappa$ requires refusing, since $e(\text{Pay})=-c<0=e(\text{Refuse})$
        at the tails information set, and conflicts with \cri{1}.
  \item In the commitment extension of counterfactual mugging, $e(q)=V_D(q)$, so \cri{6}$^\kappa$
        commits to paying while refusing in the original problem, and conflicts with \cri{4}.
\end{enumerate}
\end{proposition}

\begin{proof}
(i) With one always-reached information set, conditioning on the action identifies the policy, so
$e(a)=\E_{\mu_a}U(\cdot,a)=V(a)$. (ii) Conditional on tails and on the tails prescription being
\enquote{pay}, the payoff is $-c$; on \enquote{refuse} it is $0$. (iii) In $K(D)$ the commitment
is the policy variable and $\lambda$ does not depend on it, so $e(q)=\E_\lambda U_D(s_q,q)=V_D(q)$, and
$V_D(\text{Pay})>V_D(\text{Refuse})$; the tails information set is reached with probability
$\tfrac12$ under the refusing policy.
\end{proof}

\begin{theorem}[Map with the evidential criterion]\label{thm:map16}
On $\cls^*$ the inclusion-minimal inconsistent subsets of $\{\cri{1},\dots,\cri{5},\cri{6}^\kappa\}$
are
\[
  \{\cri{3}\},\quad\{\cri{1},\cri{2}\},\quad\{\cri{2},\cri{4}\},\quad
  \{\cri{1},\cri{6}^\kappa\},\quad\{\cri{2},\cri{6}^\kappa\},\quad\{\cri{4},\cri{6}^\kappa\},
\]
and the maximal consistent subsets are $\{\cri{1},\cri{4},\cri{5}\}$, $\{\cri{2},\cri{5}\}$, and
$\{\cri{6}^\kappa,\cri{5}\}$.
\end{theorem}

\begin{proof}
Among \cri{1}, \cri{2}, \cri{4}, \cri{6}$^\kappa$, every pair except $\{\cri{1},\cri{4}\}$ is
inconsistent by Propositions~\ref{prop:newcomb}, \ref{prop:c2c4}, and~\ref{prop:c6}, and each listed
maximal set has a total witness.
\end{proof}

A different evidential semantics can change this table; Section~\ref{sec:evsem} compares two. The
theorem is a statement about the displayed definitions.

\subsection{Noisy commitment and couplings}\label{sec:couplings}

Proposition~\ref{prop:c2c4} uses exact prediction and, implicitly, a choice of coupling. Both
matter, and they matter differently for structural and for expectational comparisons. Let the
opaque prize be $m$ and the transparent bonus $c$, with $m>c>0$, and let the symmetric error be
$\varepsilon\in[0,1]$. Write $g=m(1-2\varepsilon)-c$ for the ex-ante advantage of $O$.

Let $X$ indicate a full opaque box under commitment to $O$, and $Y$ a full opaque box under commitment
to $B$, so that $\Pr(X=1)=1-\varepsilon$ and $\Pr(Y=1)=\varepsilon$. Every coupling is determined by
$t=\Pr(X=1,Y=1)$, which ranges over $0\le t\le\min(\varepsilon,1-\varepsilon)$:

\begin{center}
\small
\begin{tabular}{@{}ccccc@{}}
  \toprule
  $(X,Y)$ & probability & payoff of $O$ & payoff of $B$ & difference $O-B$ \\
  \midrule
  $(1,1)$ & $t$ & $m$ & $m+c$ & $-c$ \\
  $(1,0)$ & $1-\varepsilon-t$ & $m$ & $c$ & $m-c$ \\
  $(0,1)$ & $\varepsilon-t$ & $0$ & $m+c$ & $-m-c$ \\
  $(0,0)$ & $t$ & $0$ & $c$ & $-c$ \\
  \bottomrule
\end{tabular}
\end{center}

\begin{proposition}[Noisy commitment in Newcomb's problem]\label{prop:noisycommit}
\begin{enumerate}[label=(\roman*)]
  \item For every coupling, the expected causal value of committing to $q$ is $V_D(q)$ whatever the
        agent's beliefs about its own commitment; every criterion that ranks commitments by expected
        causal value requires commitment to $O$ exactly when $g>0$, while the same criterion requires
        $B$ in $D$.
  \item With structural states $S_K=\supp\lambda$, commitment to $O$ strictly dominates commitment to $B$
        if and only if $\varepsilon=0$; commitment to $B$ strictly dominates if and only if
        $\varepsilon\ge\tfrac12$ and $t=1-\varepsilon$.
  \item With structural states $S_K=\prod_q\supp\mu_q$, neither commitment dominates for
        $0<\varepsilon<1$, whatever the coupling.
\end{enumerate}
Consequently the \cri{2}--\cri{4} conflict of Proposition~\ref{prop:c2c4} occurs, on this pair of
problems, only at $\varepsilon=0$ and under every coupling, while the reversal of expectational
comparisons between $D$ and $K_\lambda(D)$ occurs whenever $g>0$, independently of $\lambda$.
\end{proposition}

\begin{proof}
(i) Since $\lambda$ does not depend on the commitment, the causal and the ex-ante value of committing
to $q$ coincide and equal $V_D(q)$ by Proposition~\ref{prop:commitvalue}. In $D$ the bonus $c$ is
gained in every state, so every expectational causal comparison favours $B$. (ii) $O$ dominates
only if the three rows with negative difference are null: $t=0$ and $\varepsilon-t=0$, so
$\varepsilon=0$; conversely, at $\varepsilon=0$ only the row $(1,0)$ remains. Directly: at positive error
$O$ pays $0$ with positive probability whereas $B$ always pays at least $c>0$, so no coupling makes
$O$ even weakly dominant. $B$ dominates exactly when the row $(1,0)$ is null, $t=1-\varepsilon$, which
the constraint $t\le\min(\varepsilon,1-\varepsilon)$ allows only for $\varepsilon\ge\tfrac12$. (iii) For
$0<\varepsilon<1$ all four rows are structural states, and the differences have both signs.
\end{proof}

Changing a probability law cannot change \cri{2} while its structural state set is held fixed; the
coupling matters for \cri{2} only through the convention $S_K=\supp\lambda$. The three couplings one
might call natural give
\begin{center}
\small
\begin{tabular}{@{}lc@{}}
  \toprule
  coupling & $t$ \\
  \midrule
  independent (product) & $\varepsilon(1-\varepsilon)$ \\
  shared predictor-correctness bit & $0$ \\
  comonotone realised payoffs & $\min(\varepsilon,1-\varepsilon)$ \\
  \bottomrule
\end{tabular}
\end{center}
and none of them restores dominance of $O$ at positive error. Couplings do matter in other problems:
for a shared Bernoulli variable $Z$, the payoffs $U=1+100Z$ and $W=100Z$ satisfy $U>W$ everywhere on
the support of the shared coupling, whereas the product coupling admits the reversing pair $U=1$,
$W=100$. Which couplings produce dominance is settled in general by the following criterion, a
finite form of Strassen's theorem on stochastic order \autocite{strassen1965}.

\begin{proposition}[Dominance under some coupling]\label{prop:strassen}
Let $\nu_U,\nu_W$ be finitely supported laws on $\Rl$ with distribution functions $F_U,F_W$.
\begin{enumerate}[label=(\roman*)]
  \item There is a coupling with $U\ge W$ almost surely if and only if $F_U(x)\le F_W(x)$ for all $x$.
  \item There is a coupling with $U>W$ almost surely if and only if $F_U(x)\le F_W(x^-)$ for every
        $x\in\supp\nu_U$.
\end{enumerate}
\end{proposition}

\begin{proof}
Necessity follows from the inclusions $\{U\le x\}\subseteq\{W\le x\}$ and
$\{U\le x\}\subseteq\{W<x\}$, which hold almost surely under the respective couplings. For
sufficiency, let $V$ be uniform on $(0,1]$ and set $U=F_U^{-1}(V)$, $W=F_W^{-1}(V)$, with
$F^{-1}(v)=\min\{y:F(y)\ge v\}$. In case (i), $F_U\le F_W$ gives $F_W^{-1}\le F_U^{-1}$ pointwise.
In case (ii), let $x=F_U^{-1}(v)\in\supp\nu_U$; then $F_W(x^-)\ge F_U(x)\ge v$, and since $\nu_W$ is
finitely supported, $F_W(y)\ge v$ for some $y<x$, so $F_W^{-1}(v)<x$.
\end{proof}

A payoff coupling lifts to a coupling of the underlying states by distributing mass within payoff
fibres according to the conditional state laws. For a general problem, the residual content of
\cri{4} is then an agreement condition: some \cri{2}-permitted original policy must agree, on the
information sets it positively reaches, with a commitment that no other commitment structurally
dominates. Proposition~\ref{prop:strassen} decides which commitments are dominated; it does not by
itself decide whether \cri{2} and \cri{4} are consistent on a whole class of problems with
imperfect prediction, which remains open (Section~\ref{sec:conclusion}).

% -----------------------------------------------------------------------------------
\section{Robustness of the within-problem conflict}\label{sec:robust}

The central conflict between \cri{1} and \cri{2} was obtained with deterministic policies, a single
always-reached choice, and a predictor targeting that choice. This section asks which of these
features it needs.

\subsection{Imperfect prediction}

By Proposition~\ref{prop:newcomb}, the conflict holds for every $0\le\varepsilon<\tfrac14$ at the
scaled payoffs, and for every $\varepsilon<\tfrac12$ at suitable stakes, since $m(1-2\varepsilon)>c$ can
be arranged whenever $\varepsilon<\tfrac12$. At $\varepsilon=\tfrac12$ the symmetric predictor carries no
policy dependence and this mechanism disappears. Theorem~\ref{thm:characterisation} gives the
general statement: what is needed is not accuracy but any difference between the state laws that
two prescriptions induce.

\subsection{Randomisation does not reconcile ex-ante and causal evaluation}

Let the agent one-box with probability $p$. Suppose first that the predictor observes the
distribution, filling the box with probability $\varepsilon+(1-2\varepsilon)p$, and that the action is
drawn independently given $p$. Expanding,
\[
  V(p)=p\cdot2\bigl(\varepsilon+(1-2\varepsilon)p\bigr)+(1-p)\bigl(2\varepsilon+2(1-2\varepsilon)p+1\bigr)
      =1+2\varepsilon+(1-4\varepsilon)p .
\]
The same formula holds if the predictor instead observes the sampled action and predicts it with
symmetric error $\varepsilon$. For $\varepsilon<\tfrac14$, \cri{1} requires $p=1$ and \cri{2} requires $p=0$.
The conflict also admits an exact quantitative form. Put $g=1-4\varepsilon>0$, let
$R_1(p)=\max_qV(q)-V(p)$ be ex-ante regret, and let $R_{\mathrm{causal}}(p)$ be the expected regret
relative to the best action in the realised state.

\begin{proposition}[Exact trade-off]\label{prop:tradeoff}
In randomised Newcomb, $R_1(p)=g(1-p)$ and $R_{\mathrm{causal}}(p)=p$, so that
\[
  \frac{R_1(p)}{g}+R_{\mathrm{causal}}(p)=1 .
\]
Simultaneous tolerances $R_1\le\delta_1$ and $R_{\mathrm{causal}}\le\delta_2$ are achievable if and only
if $\delta_1/g+\delta_2\ge1$. Under infallible prediction, every policy has regret at least $\tfrac12$
by one of the two measures.
\end{proposition}

\begin{proof}
$V$ is affine with slope $g$ and maximal at $p=1$, which gives $R_1$. In every realised state
two-boxing is better by exactly one unit, so the causal regret is $1$ when the agent one-boxes and
$0$ otherwise. The tolerances require $1-\delta_1/g\le p\le\delta_2$, a nonempty interval exactly
under the stated inequality. At $\varepsilon=0$, $g=1$ and $\max(R_1,R_{\mathrm{causal}})\ge\tfrac12$.
\end{proof}

\subsection{Existence of optimal and ratifiable mixed policies}\label{sec:mixedexist}

Admitting mixed policies does not by itself make either \cri{1} or \cri{3} satisfiable; regularity
of the predictor's dependence on the distribution is what matters.

\begin{proposition}\label{prop:mixedfail}
\begin{enumerate}[label=(\roman*)]
  \item Let a binary reward, independent of the realised action, succeed with probability $q(p)=p$ for
        $p<1$ and $q(1)=0$. Then $V(p)=q(p)$ has supremum $1$ and no maximiser.
  \item Let Death choose city $1$ if and only if $p\ge\tfrac12$, where $p$ is the probability of city $1$,
        with payoff $1$ for avoiding Death. No mixture is ratifiable, although Death observes only the
        distribution.
\end{enumerate}
\end{proposition}

\begin{proof}
(i) is immediate. (ii) If $p<\tfrac12$, Death is in city $0$, and city $1$ is the unique causal best
response, so a ratifiable $p$ would put all its mass on city $1$; if $p\ge\tfrac12$, city $0$ is the
unique best response. Either way $p$ is not supported on best responses to itself.
\end{proof}

\begin{theorem}[Existence under continuity]\label{thm:mixedexist}
Let actions and states be finite, let the action be sampled independently given the mixed policy
$p\in\Delta(A)$, and let $p\mapsto\mu_p\in\Delta(S)$ be continuous. Then a mixed policy satisfying
\cri{1} exists, and so does a mixed policy all of whose supported actions are causal best responses
to $\mu_p$.
\end{theorem}

\begin{proof}
$V(p)=\sum_ap_a\sum_s\mu_p(s)u(s,a)$ is continuous on the compact simplex, so it attains its maximum.
For ratifiability, let $Q(a;p)=\sum_s\mu_p(s)u(s,a)$ and $B(p)=\Delta\bigl(\argmax_aQ(a;p)\bigr)$. The
correspondence $B$ has nonempty, compact, convex values, and its graph is closed because the
finitely many functions $Q(a;\cdot)$ are continuous. By Kakutani's theorem
\autocite{kakutani1941} there is $p\in B(p)$.
\end{proof}

Theorem~\ref{thm:mixedexist} is the one-shot core of two published existence results with
stronger structure: Halpern's theorem on ratifiable acts, which assumes costless statewise mixing
and conditional state laws that are affine in the mixture \autocite{halpern2020}, and the theorem of
\textcite{belletal2021} that every continuous Newcomblike decision process has a strongly ratifiable
policy. Neither identifies ratifiable policies with ex-ante maximisers, and
Proposition~\ref{prop:tradeoff} shows that they cannot be identified.

\begin{remark}[Predicting a realised sample]\label{rem:samples}
If Death independently samples the agent's distribution, the usual symmetric model admits
$p=\tfrac12$. If Death instead predicts the agent's realised private sample, the verdict depends on
further semantics: whether the agent knows its sample, and whether a deviation preserves or
replaces the randomness the predictor tracks. Conditioning on a revealed, perfectly predicted
sample restores the deterministic instability. A marginal test of support-optimality can still
accept $p=\tfrac12$ while ignoring the correlation between the prediction and the sample, and should
not be confused with a criterion that preserves that correlation when the unchanged policy is
evaluated. Skyrms already cautioned against an unconditional randomisation rescue
\autocite{skyrms1990ratif}.
\end{remark}

\subsection{Restrictions on what predictors may target}\label{sec:targets}

Some responses to Newcomb-type problems restrict prediction to information sets that are in some
sense actual. Newcomb's problem has one information set, reached with probability one under every
policy, so its conflict survives every such restriction except the last:

\begin{center}
\small
\begin{tabular}{@{}>{\raggedright\arraybackslash}p{19em}c@{}}
  \toprule
  admissible prediction targets & Newcomb conflict survives \\
  \midrule
  information sets positively reached under the evaluated policy & yes \\
  information sets reached under every admissible policy & yes \\
  information sets on the realised path & yes \\
  information sets already reached before the prediction occurs & no \\
  \bottomrule
\end{tabular}
\end{center}

The last restriction removes anticipatory dependence altogether: a \enquote{prediction} of an action
already taken is an ordinary causal response, which belongs in the structural plan, so $\mu_\pi$ no
longer depends on $\pi$. In that case ex-ante and causal evaluation agree.

\begin{proposition}\label{prop:causaltree}
If $\mu_\pi=\mu$ does not depend on $\pi$, then every \cri{1}-optimal policy satisfies \cri{3},
and therefore \cri{2}.
\end{proposition}

\begin{proof}
By Lemma~\ref{lem:local},
\begin{align*}
  V_D(\dev{\pi}{h}{a})-V_D(\pi)
  &=\sum_{s\in R_h(\pi)}\mu(s)\bigl(U_D(s,\dev{\pi}{h}{a})-U_D(s,\pi)\bigr)\\
  &=r_\pi(h)\bigl(Q_\pi(h,a)-Q_\pi(h,\pi(h))\bigr).
\end{align*}
If $r_\pi(h)>0$ and $\pi(h)$ is not causally optimal at $h$, the deviation raises the ex-ante value.
\end{proof}

Removing predictions about unreached information sets eliminates particular witnesses, such as
counterfactual mugging, and so changes the evidential part of the map in Theorem~\ref{thm:map16}.
It leaves the central Newcomb incompatibility intact.


\part{Rules across problems}\label{part:across}
% -----------------------------------------------------------------------------------
\section{The causal matrix and cross-problem axioms}\label{sec:matrix}

Theorem~\ref{thm:map15} shows that within-problem criteria yield only singleton and pairwise
obstructions. If a non-trivial impossibility in the style of Arrow exists, it must therefore come
from requirements that relate a theory's choices on \emph{different} problems: invariance,
independence, contraction. Such requirements need a common description of problems, which the
following reduced form provides.

\subsection{Reduced form}

\begin{definition}[Causal matrix]\label{def:matrix}
The \emph{causal matrix} of a problem $D$ is the square matrix indexed by $\Pol_D$,
\[
  M_D[\pi,\sigma]=\sum_{s\in S_D}\mu_\pi(s)\,U_D(s,\sigma),
\]
the value of running $\sigma$ while the environment responds as it would to $\pi$.
\end{definition}

Row $\pi$ is the causal evaluation of every policy under the beliefs that $\pi$ induces; column
$\sigma$ records how $\sigma$ fares under every such belief state; and the diagonal is the ex-ante
value, $M_D[\pi,\pi]=V_D(\pi)$. The matrix compresses the problem: it retains expectations under
each induced law, and it forgets the structural states. Two consequences should be kept in view. The
column-dominance axiom below compares expected values, whereas \cri{2} compares structural outcomes,
including states of probability zero; and \enquote{$\pi$ maximises its own row} is a global form of
ratifiability, which differs from the local, information-set form \cri{3} in multi-stage problems.
For problems with a single choice, the matrix and the structural criteria agree, and all the
witnesses below are of that kind.

\begin{theorem}[Realisability]\label{thm:realise}
Every finite real square matrix is the causal matrix of a fair, well-defined problem with a single
information set. Deleting a set of indices together with their rows and columns is realised by
deleting the corresponding actions while keeping the remaining laws.
\end{theorem}

\begin{proof}
Given $M\in\Rl^{I\times I}$, let the agent choose $i\in I$ at a single information set; let the states
be $S=I$; let an infallible predictor forecast the chosen index, $\mu_i=\delta_i$; and let
$U(s,j)=M_{sj}$. Then $M_D[i,j]=U(i,j)=M_{ij}$. The problem is fair and well defined, and deleting
actions leaves the laws of the others unchanged.
\end{proof}

In this realisation \cri{1} selects the maximal diagonal entries, and \cri{3} accepts exactly the
indices $i$ whose diagonal entry is maximal in row $i$. Rules on matrices therefore have an
unrestricted domain, as social welfare functions do in Arrow's theorem
\autocite{arrow1951}, and the framework contains every finite recommendation structure.

\subsection{Rules and axioms}

Fix an infinite universe $\mathbb L$ of labels. A \emph{matrix} is a pair $(I,M)$ with $I\subseteq\mathbb L$
finite and nonempty and $M\in\Rl^{I\times I}$; for $\varnothing\ne K\subseteq I$, $M|_K$ is the
\emph{principal restriction} to $K\times K$, which keeps labels. A bijection $\sigma\colon I\to I'$ acts by
$(\sigma\cdot M)_{\sigma(i)\sigma(j)}=M_{ij}$.

\begin{definition}[Rule]\label{def:rule}
A \emph{rule} assigns to every matrix a reflexive relation $\succeq_M$ on $I$, with strict part
$\succ_M$, and recommends its maximal elements,
\[
  G(M)=\{i\in I:\text{there is no }j\in I\text{ with }j\succ_M i\},
\]
which are required to be nonempty for every $M$.
\end{definition}

The convention makes the comparison axioms below constrain the same relation that determines
choice, and it prevents a rule that lacks an ordering from recommending nothing. A decision theory
satisfies \emph{reduction} (\ax{A0}) if its choice on $D$ depends only on the labelled matrix $M_D$,
so that it selects from $G(M_D)$ for some rule. A \emph{single-valued} deterministic selector
cannot be neutral: on the zero $2\times2$ matrix, swapping the indices fixes the input but moves
either possible answer. Choice correspondences avoid this separate obstruction.

\begin{definition}[Cross-problem axioms]\label{def:axioms}
A rule may satisfy the following conditions, for all matrices $M$ and all indices $i,j$.
\begin{description}[font=\normalfont\sffamily,leftmargin=2.4em,labelwidth=1.6em,labelsep=.8em,itemsep=.35em]
  \item[N] \emph{Neutrality.} For every bijection $\sigma\colon I\to I'$:
        $i\succeq_Mj\iff\sigma(i)\succeq_{\sigma\cdot M}\sigma(j)$.
  \item[F] \emph{Common affine invariance.} $\succeq_{aM+b\one\one^\top}=\succeq_M$ for $a>0$, $b\in\Rl$.
  \item[R] \emph{Row translations.} $\succeq_{M+b\one^\top}=\succeq_M$ for every $b\in\Rl^I$
        (adding $b_k$ to every entry of row $k$).
  \item[P] \emph{Column Pareto.} If $M_{k\sigma}\ge M_{k\tau}$ for every $k\in I$, strictly for some $k$,
        then $\sigma\succ_M\tau$.
  \item[L] \emph{Restricted Pareto.} The same, with $k$ ranging only over $\{\sigma,\tau\}$.
  \item[D] \emph{Diagonal optimality.} $G(M)\subseteq\argmax_iM_{ii}$.
  \item[J] \emph{Global ratifiability.} Every $i\in G(M)$ satisfies $M_{ii}\ge M_{ij}$ for all $j$.
  \item[B] \emph{Binary independence.} $i\succeq_Mj\iff i\succeq_{M|_{\{i,j\}}}j$.
  \item[S] \emph{Self-sovereignty.} If $M,M'\in\Rl^{I\times I}$ have the same rows $i$ and $j$, then
        $i\succeq_Mj\iff i\succeq_{M'}j$.
  \item[C] \emph{Contraction.} $G(M)\cap K\subseteq G(M|_K)$ for every nonempty $K\subseteq I$.
  \item[O] \emph{Ordering.} $\succeq_M$ is complete and transitive.
  \item[Q] \emph{Non-triviality.} For some $I$ there are $M,M'\in\Rl^{I\times I}$ with $G(M)\ne G(M')$.
\end{description}
\end{definition}

\ax{P} is the expectational counterpart of dominance; \ax{D} is \cri{1}; \ax{J} is ratifiability of
pure choices; \ax{S} is the condition of \textcite[§6]{briggs2010}, under which the comparison of two
options depends only on how each of them evaluates all options; \ax{C} is Chernoff's condition, Sen's
property~$\alpha$ \autocite{chernoff1954,sen1971}; \ax{B} is an independence of irrelevant alternatives
for binary comparisons, of the kind that \textcite{podgorski2022} imposes on overall-value theories,
required here across menu sizes. Neutrality is
stated for bijections between different label sets, because the proofs below transport choices from
one menu to another; invariances are imposed on the relation, not only on its maximal elements.
Ordering concerns the actual output relation at each input. Requiring only that each $G(M)$ be
rationalisable by some weak order would be vacuous, since any nonempty set can be placed in a top
indifference class; requiring one order to rationalise the choices from all restrictions of a
given matrix is stronger and already entails contraction \autocite{arrow1959}.

\begin{lemma}\label{lem:bimplies}
\ax{B} implies \ax{S} and \ax{C}.
\end{lemma}

\begin{proof}
If rows $i$ and $j$ coincide in $M$ and $M'$, so does the principal block on $\{i,j\}$, and \ax{B}
transfers the comparison. For \ax{C}, let $i\in G(M)\cap K$. Any $j\in K$ with $j\succ_{M|_K}i$ would
satisfy $j\succ_{M|_{\{i,j\}}}i$ and hence $j\succ_Mi$, contradicting the maximality of $i$ in $M$.
\end{proof}

The following monotonicity condition is used in Section~\ref{sec:exante}. Write $E_{ii}$ for the
matrix with a single entry $1$ in position $(i,i)$.

\begin{definition}[Diagonal monotonicity, \Mon]
If $M'=M+\delta E_{ii}$ with $\delta\ge0$, then $i\in G(M)$ implies $i\in G(M')$, and every true
comparison $i\succeq_Mj$ or $i\succ_Mj$ remains true in $M'$.
\end{definition}

\subsection{What row-translation invariance means}

The axiom \ax{R} will turn out to be decisive, and it is tempting to read it as the matrix form of
causal evaluation. It deserves a precise interpretation.

\begin{theorem}[Factorisation]\label{thm:factor}
A rule satisfies \ax{R} if and only if $\succeq_M$ depends on $M$ only through the row-normalised
matrix $Z(M)_{ij}=M_{ij}-M_{ii}$, equivalently through the within-row differences of $M$.
\end{theorem}

\begin{proof}
Row translations leave $Z$ unchanged. Conversely, if $Z(M)=Z(M')$, then
$M'_{ij}-M_{ij}=M'_{ii}-M_{ii}$ is constant along each row, so $M'$ is a row translation of $M$.
\end{proof}

On a restricted domain, \ax{R} is factorisation through the image of $Z$; the zero-diagonal
representative need not itself be admissible. With this, familiar criteria can be classified once a
matrix semantics is fixed for each.

\begin{center}
\small
\renewcommand{\arraystretch}{1.25}
\begin{tabularx}{\textwidth}{@{}>{\raggedright\arraybackslash}p{9.6em}
                                  >{\raggedright\arraybackslash}X
                                  >{\raggedright\arraybackslash}p{6.6em}@{}}
  \toprule
  criterion & matrix semantics & satisfies \ax{R}? \\
  \midrule
  fixed-prior causal decision theory & rank by $w^\top M$ for a fixed weight vector $w$ & yes \\
  pure global ratifiability & admit $i$ iff $M_{ii}\ge M_{ij}$ for all $j$ & yes, as an admissibility condition \\
  mixed ratifiability & $p\in\operatorname{BR}(p^\top M)$ & yes, as an equilibrium set \\
  deliberational causal decision theory & requires initial beliefs, update dynamics, and a final selection & depends on these \\
  evidential theory with reference-policy law & maximise $M_{ii}$, with $P(i,s)=\rho_i\mu_i(s)$ & no \\
  ex-ante policy maximisation & maximise $M_{ii}=V(i)$ & no \\
  \bottomrule
\end{tabularx}
\end{center}

Ratifiability's invariance does not establish that ratifiable options exist or supply a ranking.
For deliberational theories, fixed-belief causal comparisons with difference-based update dynamics
preserve row baselines, but arbitrary update dynamics need not; \textcite{joyce2012} distinguishes
equilibrium credences from the act subsequently chosen, and \textcite{arntzenius2008} identifies
the update rule as additional structure needed to define equilibrium. An \enquote{FDT-style} theory
is ex-ante maximisation under stipulated functional counterfactuals and so fails \ax{R} under that
stipulation \autocite{yudkowskysoares2017}.

\begin{proposition}\label{prop:regret}
\ax{R} does not characterise causal rationality: ranking by $r_i=M_{ii}-\max_jM_{ij}$ satisfies
\ax{N}, \ax{F}, \ax{R}, \ax{O}, yet on
\[
  M=\begin{pmatrix}0&10&-1\\10&0&-1\\0&0&-1\end{pmatrix}
\]
it chooses option $3$, which both other columns strictly dominate.
\end{proposition}

\begin{proof}
The scores are $(-10,-10,-1)$. Column $3$ is $(-1,-1,-1)^\top$, below both $(0,10,0)^\top$ and
$(10,0,0)^\top$ in every row.
\end{proof}

So \ax{R} means exactly invariance to row-specific baselines. It captures a feature shared by several
causal procedures, but it supplies neither causal expected-utility maximisation nor dominance.

% -----------------------------------------------------------------------------------
\section{A pairwise reduction under common affine invariance}\label{sec:pairwise}

The first cross-problem result shows that, with common affine invariance as the only invariance,
the axioms of Definition~\ref{def:axioms} behave like the within-problem criteria: every inconsistency
has a core of size one or two.

\begin{theorem}[Pairwise reduction]\label{thm:pairwise}
On the domain of all finite real square matrices, under principal restriction, the
inclusion-minimal inconsistent subsets of
$\{\ax{N},\ax{F},\ax{P},\ax{L},\ax{D},\ax{J},\ax{B},\ax{S},\ax{C},\ax{O},\ax{Q}\}$ are exactly
\[
  \{\ax{L}\},\quad\{\ax{J}\},\quad\{\ax{P},\ax{D}\},\quad\{\ax{P},\ax{B}\},\quad
  \{\ax{P},\ax{S}\},\quad\{\ax{P},\ax{C}\},
\]
and the maximal consistent subsets are exactly
\begin{center}
\small
\begin{tabular}{@{}ll@{}}
  \toprule
  maximal consistent subset & witness rule \\
  \midrule
  $\{\ax{N},\ax{F},\ax{P},\ax{O},\ax{Q}\}$ & rank by column sums $\sum_kM_{ki}$ \\
  $\{\ax{N},\ax{F},\ax{D},\ax{B},\ax{S},\ax{C},\ax{O},\ax{Q}\}$ & rank by diagonal entries $M_{ii}$ \\
  \bottomrule
\end{tabular}
\end{center}
\end{theorem}

The proof consists of five obstructions and a containment argument.

\begin{proposition}[Restricted Pareto alone]\label{prop:L}
No rule satisfies \ax{L}.
\end{proposition}

\begin{proof}
Consider
\[
  K=\begin{pmatrix}0&-1&1\\1&0&-1\\-1&1&0\end{pmatrix}.
\]
In rows $1$ and $2$, column $1$ reads $(0,1)$ and column $2$ reads $(-1,0)$, so \ax{L} forces $1\succ2$.
Likewise rows $2,3$ force $2\succ3$ and rows $1,3$ force $3\succ1$. Every option has a strict
defeater, so $G(K)=\varnothing$.
\end{proof}

If bare cyclic relations without a nonempty recommendation were allowed when \ax{O} is absent, the
obstruction would instead be the pair $\{\ax{L},\ax{O}\}$; the convention of
Definition~\ref{def:rule} makes it a singleton.

\begin{proposition}[Global ratifiability alone]\label{prop:J}
No rule satisfies \ax{J}.
\end{proposition}

\begin{proof}
In $\bigl(\begin{smallmatrix}0&1\\1&0\end{smallmatrix}\bigr)$ neither option maximises its own row.
This is Death in Damascus in reduced form.
\end{proof}

\begin{proposition}[Pareto against the diagonal]\label{prop:PD}
No rule satisfies \ax{P} and \ax{D}.
\end{proposition}

\begin{proof}
In $N_0=\bigl(\begin{smallmatrix}2&3\\0&1\end{smallmatrix}\bigr)$, column $2$ strictly dominates column
$1$, so \ax{P} gives $2\succ1$ and $G(N_0)=\{2\}$, while option $1$ uniquely maximises the diagonal.
This is Newcomb's problem in reduced form.
\end{proof}

\begin{proposition}[Pareto against self-sovereignty; \cite{briggs2010}]\label{prop:PS}
No rule satisfies \ax{P} and \ax{S}, and no rule satisfies \ax{P} and \ax{B}.
\end{proposition}

\begin{proof}
Let
\[
  M^+=\begin{pmatrix}0&0&-1\\0&0&-1\\1&0&-1\end{pmatrix},\qquad
  M^-=\begin{pmatrix}0&0&-1\\0&0&-1\\0&1&-1\end{pmatrix}.
\]
In $M^+$, column $1$ weakly dominates column $2$, strictly in row $3$, so \ax{P} gives $1\succ2$; in
$M^-$ it gives $2\succ1$. Rows $1$ and $2$ are the same in both matrices, contradicting \ax{S}; so is
the principal block on $\{1,2\}$, contradicting \ax{B}.
\end{proof}

This is an explicit instance of Briggs's theorem, which states the incompatibility of a Pareto
condition and self-sovereignty for decision-theoretic rankings and does not require neutrality or
transitivity \autocite[21--26]{briggs2010}. Causal decision theory satisfies her Pareto condition
and violates self-sovereignty; evidential decision theory and Wedgwood's benchmark theory
\autocite{wedgwood2013} satisfy self-sovereignty and violate Pareto. The conflict needs at least three
options: on two-option problems the benchmark theory satisfies both. The pair
$\{\ax{P},\ax{S}\}$ is therefore not new.

\begin{proposition}[Pareto against contraction]\label{prop:PC}
No rule satisfies \ax{P} and \ax{C}.
\end{proposition}

\begin{proof}
Return to the matrix $K$ of Proposition~\ref{prop:L}. Its binary restrictions are
\[
  K|_{\{1,2\}}=K|_{\{2,3\}}=\begin{pmatrix}0&-1\\1&0\end{pmatrix},\qquad
  K|_{\{1,3\}}=\begin{pmatrix}0&1\\-1&0\end{pmatrix},
\]
and \ax{P} forces $G(K|_{\{1,2\}})=\{1\}$, $G(K|_{\{2,3\}})=\{2\}$, $G(K|_{\{1,3\}})=\{3\}$. Each option
of $K$ is rejected in one binary restriction containing it, so by \ax{C} none can be chosen from $K$.
\end{proof}

The mechanism is the classical binary cycle that makes contraction consistency incompatible with
majority-type choice \autocite{sen1970,brandtharrenstein2011}; the modelling choice that gives it
force here is that restriction deletes the rows of the removed options together with their columns.
Section~\ref{sec:menus} shows that it disappears when the rows are retained.

\begin{proof}[Proof of Theorem~\ref{thm:pairwise}]
The six sets are inconsistent by Propositions~\ref{prop:L}--\ref{prop:PC}, and each proper subset of
each pair is consistent by the witnesses. Column sums satisfy \ax{N}, \ax{F} (scores are multiplied by
$a$ and shifted by $|I|b$), \ax{P} (a weakly dominant column with one strict entry has a strictly
larger sum), \ax{O}, and \ax{Q}. Diagonal ranking satisfies \ax{N}, \ax{F}, \ax{D}, \ax{B} (and hence
\ax{S} and \ax{C} by Lemma~\ref{lem:bimplies}), \ax{O}, and \ax{Q}. A consistent set contains neither
\ax{L} nor \ax{J}. If it contains \ax{P}, it excludes \ax{D}, \ax{B}, \ax{S}, and \ax{C}, and lies in
the column-sum set; otherwise it lies in the diagonal set.
\end{proof}

Two of the reference rules one might expect to appear do not. Maximin over rows violates the strong
Pareto axiom: in $\bigl(\begin{smallmatrix}0&0\\2&1\end{smallmatrix}\bigr)$, column $1$ weakly dominates
column $2$, but both column minima are $0$. Positive fixed-weight column averages satisfy \ax{P}, and
uniform weights add \ax{N}, but fixed unequal weights require persistent labels or additional input.
The selector $T^\nu$ of Proposition~\ref{prop:nu} is in general not a rule in the sense of
\ax{A0}: it uses structural information that $M_D$ does not retain.

With common affine invariance alone, then, there is no Arrow-style family among these axioms.

% -----------------------------------------------------------------------------------
\section{Row-baseline invariance and contraction}\label{sec:translations}

Adding row-translation invariance changes the picture. It produces a genuine impossibility among
four independent axioms, none of which concerns Pareto, ordering, or ex-ante value.

\subsection{The circulant theorem}

The proof embeds an arbitrary matrix in a matrix with a transitive symmetry group. It rests on a
distinct-differences lemma of the kind used by \textcite{finkruiz1986} to embed graphs in
circulants.

\begin{lemma}[Distinct differences]\label{lem:sidon}
Let $k\ge2$, $v_i=2^i$ for $i=0,\dots,k-1$, and $m>2(2^{k-1}-1)$. The $k(k-1)$ ordered differences
$v_j-v_i$, $i\ne j$, are distinct and nonzero modulo $m$.
\end{lemma}

\begin{proof}
For $j>i$, $v_j-v_i=2^i(2^{j-i}-1)$; the power of two and the odd part determine $i$ and $j-i$, so the
positive differences are distinct integers, and so are their negatives. Every difference has
absolute value at most $D=2^{k-1}-1$, so two distinct differences differ by a nonzero integer of
absolute value at most $2D<m$, and no difference is divisible by $m$.
\end{proof}

\begin{lemma}[Circulant embedding]\label{lem:embed}
Every matrix $Z\in\Rl^{k\times k}$ with zero diagonal is a principal submatrix of a circulant matrix
$A_{uv}=f(v-u\bmod m)$ on $\mathbb Z/m\mathbb Z$ with $f(0)=0$.
\end{lemma}

\begin{proof}
With $v_i$ and $m$ as in Lemma~\ref{lem:sidon}, set $f(0)=0$ and $f(v_j-v_i)=Z_{ij}$ for $i\ne j$, which
is well defined by the lemma, and give $f$ arbitrary values on the unused residues. Then
$A_{v_iv_j}=f(v_j-v_i)=Z_{ij}$.
\end{proof}

\begin{theorem}[Neutrality, row translations, and contraction force universal choice]\label{thm:circulant}
On the domain of all finite real square matrices, every rule satisfying \ax{N}, \ax{R}, and \ax{C}
(principal restriction) has $G(M)=I$ for every matrix $(I,M)$. Consequently
$\{\ax{N},\ax{R},\ax{C},\ax{Q}\}$ is inconsistent.
\end{theorem}

\begin{proof}
Let $(I,M)$ be given with $|I|=k\ge2$ (the case $k=1$ is trivial), and let $Z=Z(M)$, which has zero
diagonal; by \ax{R}, $G(Z)=G(M)$. By Lemma~\ref{lem:embed} and \ax{N} (relabelling the vertices of the
circulant so that the embedded ones carry the labels of $I$), $Z$ is the principal restriction to $I$
of a matrix $A$ whose symmetry group acts transitively on its indices: each cyclic shift
$u\mapsto u+1$ of $\mathbb Z/m\mathbb Z$, transported by the relabelling, maps $A$ to itself. By \ax{N},
$G(A)$ is invariant under these symmetries, and since it is nonempty it contains every index. By
\ax{C}, $G(Z)=G(A|_I)\supseteq G(A)\cap I=I$. Hence $G(M)=I$, and no two matrices on the same index set
have different choice sets.
\end{proof}

Without \ax{O} the conclusion is universal \emph{choice}, not universal indifference: incomparabilities
may remain. The symmetry step is standard in the theory of tournament solutions, where neutral,
nonempty solutions must select every vertex of a vertex-transitive tournament
\autocite{brandt2011,brandtbrillharrenstein2016}; the embedding has the precedent noted above. What
the theorem adds is the use of contraction to carry universal choice from symmetric matrices to all
matrices, and of row-baseline invariance to reach every matrix from a zero-diagonal one. No
published statement of this conclusion was found in the literature searched; its ingredients are
established, and no priority is claimed.

\begin{proposition}[Independence]\label{prop:indep}
Each three-element subset of $\{\ax{N},\ax{R},\ax{C},\ax{Q}\}$ is satisfiable, so the family is
minimally inconsistent:
\begin{center}
\small
\begin{tabular}{@{}ll@{}}
  \toprule
  omitted axiom & rule satisfying the other three \\
  \midrule
  \ax{N} & responsive fixed-label rule \\
  \ax{R} & diagonal ranking \\
  \ax{C} & row-range ranking, $r_i=\max_jM_{ij}-\min_jM_{ij}$ \\
  \ax{Q} & universal indifference \\
  \bottomrule
\end{tabular}
\end{center}
All four rules also satisfy \ax{F} and \ax{O}.
\end{proposition}

\begin{proof}
Row ranges are unchanged by row translations, multiplied by $a$ under \ax{F}, and permuted under
relabelling; the rule is nonconstant. Diagonal ranking satisfies \ax{C} by Lemma~\ref{lem:bimplies}.
The \emph{responsive fixed-label rule} fixes two labels $u,v$ and a well-order of $\mathbb L$ in which
$u,v$ form an adjacent block above all other labels. When both occur, $u$ and $v$ are compared by the
sign of $M_{uu}-M_{uv}$ (a tie if it is zero); all other comparisons follow the fixed order. The
comparison of $u$ and $v$ uses a within-row difference of their principal block, so the rule satisfies
\ax{R}, \ax{F}, and \ax{B}, hence \ax{C}; it is responsive, so \ax{Q} holds, but it is not neutral.
\end{proof}

If neutrality is regarded as part of the reduction \ax{A0} rather than as an axiom, the same result is
a minimal three-axiom family $\{\ax{R},\ax{C},\ax{Q}\}$ relative to that background.

\begin{proposition}\label{prop:NRBQ}
$\{\ax{N},\ax{R},\ax{B},\ax{Q}\}$ is minimally inconsistent, and so is $\{\ax{R},\ax{D}\}$.
\end{proposition}

\begin{proof}
Suppose $i\succ_Mj$ in some matrix. By \ax{B} the same holds in $M|_{\{i,j\}}$, and by \ax{R} in its
row normalisation $H=\bigl(\begin{smallmatrix}0&x\\y&0\end{smallmatrix}\bigr)$, with $i$ first. Let
\[
  C=\begin{pmatrix}0&x&y\\y&0&x\\x&y&0\end{pmatrix}.
\]
Its principal blocks on $(1,2)$, $(2,3)$, and $(3,1)$ all equal $H$, so \ax{N} and \ax{B} give
$1\succ2$, $2\succ3$, and $3\succ1$, and $G(C)=\varnothing$. Hence no strict comparison ever occurs and
$G(M)=I$ always, contradicting \ax{Q}. Minimality follows from the witnesses of
Proposition~\ref{prop:indep}, with row-range ranking now serving as the rule that omits \ax{B}. For $\{\ax{R},\ax{D}\}$: adding $2$ to the
second row of $\bigl(\begin{smallmatrix}1&0\\0&0\end{smallmatrix}\bigr)$ changes the unique diagonal
maximiser from option $1$ to option $2$, while \ax{R} keeps the choice fixed.
\end{proof}

Unlike Theorem~\ref{thm:circulant}, this argument uses the link between comparisons and nonempty
maximal choice fixed in Definition~\ref{def:rule}.

\begin{theorem}[Map with row translations]\label{thm:augmented}
Adding \ax{R} to the family of Theorem~\ref{thm:pairwise} introduces exactly three new minimal
inconsistent subsets,
\[
  \{\ax{R},\ax{D}\},\qquad\{\ax{N},\ax{R},\ax{B},\ax{Q}\},\qquad\{\ax{N},\ax{R},\ax{C},\ax{Q}\},
\]
and the maximal consistent subsets become
\begin{center}
\small
\begin{tabular}{@{}ll@{}}
  \toprule
  witness rule & maximal consistent subset \\
  \midrule
  column sums & $\{\ax{N},\ax{F},\ax{P},\ax{R},\ax{O},\ax{Q}\}$ \\
  diagonal ranking & $\{\ax{N},\ax{F},\ax{D},\ax{B},\ax{S},\ax{C},\ax{O},\ax{Q}\}$ \\
  row ranges & $\{\ax{N},\ax{F},\ax{R},\ax{S},\ax{O},\ax{Q}\}$ \\
  responsive fixed-label rule & $\{\ax{F},\ax{R},\ax{B},\ax{S},\ax{C},\ax{O},\ax{Q}\}$ \\
  universal indifference & $\{\ax{N},\ax{F},\ax{R},\ax{B},\ax{S},\ax{C},\ax{O}\}$ \\
  \bottomrule
\end{tabular}
\end{center}
\end{theorem}

\begin{proof}
Column sums satisfy \ax{R} because a row translation shifts every column sum by $\sum_kb_k$; row ranges
satisfy \ax{S} because each score depends only on its own row. Let a set avoid all minimal
inconsistent subsets. If it contains \ax{P}, its pairwise conflicts confine it to the first row. If it
omits \ax{P} and \ax{R}, it fits the second. If it contains \ax{R}, it omits \ax{D}; if it contains both
\ax{N} and \ax{Q}, the two four-element obstructions exclude \ax{B} and \ax{C}, leaving the third row;
without \ax{N} it fits the fourth, and without \ax{Q} the fifth.
\end{proof}

\subsection{An exact domain condition}

Theorem~\ref{thm:circulant} needs matrices of unbounded size: with a fixed maximum size $n$, the rule
that chooses everything below size $n$ and maximises row ranges at size $n$ satisfies \ax{N}, \ax{R},
\ax{C}, and \ax{Q}. The exact requirement on the domain can be stated without quantifying over
rules. Let $\mathcal D$ be a set of matrices closed under relabelling, and interpret each invariance as
applying whenever both the matrix and its transform belong to $\mathcal D$. Form a directed graph on the
\emph{pointed matrices} $\mathcal P_{\mathcal D}=\{(M,i):M\in\mathcal D,\ i\in I_M\}$, with
\begin{enumerate}[label=(\roman*)]
  \item edges in both directions between $(M,i)$ and $(\sigma\cdot M,\sigma(i))$ for admissible relabellings;
  \item edges in both directions between $(M,i)$ and $(M+b\one^\top,i)$ for admissible row translations;
  \item extension edges $(M|_K,i)\to(M,i)$ for $i\in K$, whenever $M|_K$ and $M$ are in $\mathcal D$.
\end{enumerate}
Let $\Reach(x)$ be the set of pointed matrices reachable from $x$ (including $x$), and let
$\Phi_M=\{(M,i):i\in I_M\}$ be the fibre of $M$.

\begin{theorem}[Exact domain criterion]\label{thm:DR}
On $\mathcal D$, the axioms $\ax{N},\ax{R},\ax{C},\ax{Q}$ are jointly inconsistent if and only if
\[
  \text{for every }x\in\mathcal P_{\mathcal D}\text{ there is }M\in\mathcal D\text{ with }\Phi_M\subseteq\Reach(x).
  \tag{DR}
\]
\end{theorem}

\begin{proof}
Say $(M,i)$ is \emph{rejected} if $i\notin G(M)$. Rejection propagates along relabelling and
translation edges by \ax{N} and \ax{R}, and along extension edges by the contrapositive of \ax{C}. Under
(DR), a single rejection at $x$ would reject the entire fibre of some $M$, contradicting $G(M)\neq\varnothing$;
so nothing is ever rejected, and \ax{Q} fails. Conversely, if (DR) fails at $x=(M_0,i_0)$, set
\[
  G_x(M)=\{i:(M,i)\notin\Reach(x)\}.
\]
Each $G_x(M)$ is nonempty because no fibre lies in $\Reach(x)$. Since relabelling and translation
edges are bidirectional, $\Reach(x)$ is closed under them, which gives \ax{N} and \ax{R}. If
$i\in G_x(M)\cap K$ but $i\notin G_x(M|_K)$, the extension edge would put $(M,i)$ in $\Reach(x)$; so \ax{C}
holds. Finally $i_0\notin G_x(M_0)$, so $G_x(M_0)$ is a nonempty proper subset of $I_{M_0}$, and some
relabelling of $M_0$ on the same index set moves it, which gives \ax{Q}. Tying chosen options above
rejected ones yields \ax{O} as well.
\end{proof}

The circulant condition is sufficient for (DR) but not necessary. If \ax{F} is also imposed, admissible
affine transformations are added as bidirectional edges.

\subsection{Physical restrictions on problems}

Does the richness that (DR) demands survive natural restrictions on the problems themselves?

\begin{proposition}[Bounded payoffs]\label{prop:bounded}
On the domain of $[0,1]$-valued matrices, $\{\ax{N},\ax{R},\ax{C},\ax{Q}\}$ is consistent; adding \ax{F}
makes it inconsistent again.
\end{proposition}

\begin{proof}
Fix $r\in[\tfrac12,1)$ and put $i\to j$ if $M_{ii}-M_{ij}>r$ and $M_{ji}-M_{jj}>r$. An edge forces
$M_{ii}>r\ge\tfrac12\ge1-r>M_{jj}$, so no vertex has both an incoming and an outgoing edge and the graph
is acyclic. Choose the options with no incoming edge, tied above the rejected ones, which are tied
among themselves. The conditions use only within-row differences, giving \ax{R}; edges depend only on
principal $2\times2$ blocks and persist under restriction, giving \ax{C}; the rule is neutral, ordered,
and responsive (it chooses $\{1\}$ from $\bigl(\begin{smallmatrix}1&0\\1&0\end{smallmatrix}\bigr)$ and
both options from the zero matrix); and raising a diagonal entry only removes incoming and adds
outgoing edges, so it satisfies \Mon. It fails \ax{F}, since scaling by a small $a>0$ removes every
edge.

With \ax{F}, given $M\in[0,1]^{n\times n}$ and small $\alpha>0$, pass from $M$ to $\alpha M$ by \ax{F} and
then to $A_{ij}=\tfrac12+\alpha(M_{ij}-M_{ii})$ by \ax{R}; every intermediate matrix stays in $[0,1]$.
Embed $A$ in a circulant as in Lemma~\ref{lem:embed}, filling unused residues with $\tfrac12$, and run
the proof of Theorem~\ref{thm:circulant}.
\end{proof}

\begin{proposition}[Prediction bounded away from accuracy]\label{prop:noisy}
Let $P_n=(1-\eta)I+\tfrac{\eta}{n}\one\one^\top$ with $0<\eta<1$, a prediction channel whose
probability of a correct forecast is $1-\eta+\eta/n\le1-\eta/2$ for $n\ge2$. With unrestricted real
payoffs, every $M\in\Rl^{n\times n}$ is the causal matrix of a problem whose predictor uses $P_n$. With
payoffs in $[0,1]$ and \ax{F}, the circulant argument survives, uniformly in $n$.
\end{proposition}

\begin{proof}
$P_n$ is invertible, with $P_n^{-1}=\frac1{1-\eta}I-\frac{\eta}{(1-\eta)n}\one\one^\top$. Let the state be
the forecast label, drawn from row $i$ of $P_n$ when $i$ is chosen, and let $U=P_n^{-1}M$ be the
state-by-action payoff table; then the causal matrix is $P_nU=M$. Since $P_n^{-1}\one\one^\top=\one\one^\top$,
the constant matrix $\tfrac12\one\one^\top$ is mapped to itself, and the absolute row sums of $P_n^{-1}$ are
at most $(1+\eta)/(1-\eta)$; so a neighbourhood of the constant matrix whose radius does not depend on
$n$ is mapped into $[0,1]$, and the compressed matrices of Proposition~\ref{prop:bounded} remain
realisable.
\end{proof}

An accuracy bound alone, then, does not destroy richness. Completely uninformative prediction is
different: if all induced laws coincide, the channel loses rank, every row of the causal matrix is
the same, and the problem is an ordinary causal one.

\begin{proposition}[Bounded action sets]\label{prop:binarytree}
For every $k\ge2$, the symmetric extensions required by Theorem~\ref{thm:circulant} are realisable
by problems with perfect recall and at most $k$ actions per information set, provided contraction
applies to arbitrary sets of policies.
\end{proposition}

\begin{proof}
Take a complete binary decision tree of depth $m$, with the agent observing its own past actions
and the prediction hidden until the end. It has $2^m-1$ information sets and $N=2^m$ action paths,
and each path is produced by exactly $L=2^{2^m-1-m}$ policies, since prescriptions off the path are
free. Let a predictor forecast the eventual path and let the payoff depend on the forecast and the
actual path; any path matrix $C$ is then realisable, and the policy matrix is $C\otimes J_L$, with $J_L$
the all-ones $L\times L$ matrix. If $C$ is circulant, rotating path blocks and permuting duplicate
policies within blocks act transitively on the policies. Any target matrix embeds in a circulant
whose size is a power of two (Lemma~\ref{lem:embed}), and is recovered by restricting to one
representative policy from each designated block.
\end{proof}

The proviso matters: arbitrary sets of surviving policies are generally not obtained by pruning local
action sets, which produces product-shaped sets. Restricting contraction to literal pruning is a
different axiom (Section~\ref{sec:menus}).

\begin{proposition}[A single bounded choice]\label{prop:ceiling}
With one information set and an unbounded action set, every square matrix is realisable. With one
information set and at most $k\ge2$ actions, the rule that is universally indifferent below size $k$
and ranks by $M_{ii}-\frac1k\sum_jM_{ij}$ at size $k$ satisfies \ax{N}, \ax{F}, \ax{R}, \ax{C}, \ax{O},
\ax{Q}, and \Mon.
\end{proposition}

\begin{proof}
The first claim is Theorem~\ref{thm:realise}. For the second, every proper restriction of a
size-$k$ matrix lands where every option is chosen, so \ax{C} holds; the score is a within-row
difference, invariant under \ax{F} up to scale, neutral, responsive, and increasing in $M_{ii}$.
\end{proof}

No finite size ceiling, therefore, forces triviality under the existential form of \ax{Q}. If
non-triviality is required specifically on binary menus, the three-option circulant
$\bigl(\begin{smallmatrix}0&x&y\\y&0&x\\x&y&0\end{smallmatrix}\bigr)$ already blocks it as soon as
three options are available.

% -----------------------------------------------------------------------------------
\section{The semantics of a smaller menu}\label{sec:menus}

Theorems~\ref{thm:pairwise} and~\ref{thm:circulant} suggest a dichotomy: menu-coherent evaluation
forces ex-ante evaluation, while causal evaluation cannot be menu-coherent without becoming trivial.
Both theorems, however, rest on one way of removing options from a problem: principal restriction,
which deletes the removed options' rows together with their columns. This section examines that
step. It turns out to be where the dichotomy fails.

\subsection{Three operations}

A smaller menu can be modelled in at least three ways.
\begin{description}[font=\normalfont\itshape,leftmargin=1.4em,itemsep=.35em]
  \item[Principal restriction ($\Cpr$).] Delete the removed options' rows and columns. A smaller menu
        then also removes the beliefs that the agent would have if it chose the removed options.
  \item[Column restriction ($\Ccol$).] Keep every row and delete only the columns. The input becomes a
        rectangular matrix of belief states by available options.
  \item[Fixed-reference restriction ($\Cref$).] Evaluate options by averaging rows with a reference
        distribution over belief states, as in fixed-prior causal decision theory, and keep that
        distribution when options are removed.
\end{description}

\begin{proposition}[Law-preserving removal is principal]\label{prop:principal}
Let $D'$ be obtained from $D$ by making only the policies in $K\subseteq\Pol_D$ available, and suppose
that their response laws and consequences are unchanged: $\mu'_i=\mu_i$ and $U'(s,i)=U(s,i)$ for
$i\in K$. Then $M_{D'}=M_D|_{K\times K}$.
\end{proposition}

\begin{proof}
Immediate from Definition~\ref{def:matrix}: each entry of $M_{D'}$ involves only laws and
consequences of policies in $K$.
\end{proof}

Principal restriction is thus forced by the definition of the causal matrix once law-preserving
removal is stipulated. Two qualifications apply. First, arbitrary subsets $K$ do not arise by pruning
local action sets, which produces product-shaped sets of policies; an arbitrary $K$ is realised by
treating the surviving policies as macro-actions of a new one-choice predictive problem, which
preserves the response laws but not the extensive form. Second, an intervention that actually
announces a smaller menu may change how a predictor responds; the model does not require
$\mu^{D'}_i=\mu^{D}_i$, and law preservation is an additional cross-problem assumption. Removal by
macro-actions must also be distinguished from commitment before prediction
(Definition~\ref{def:commit}), which changes the counterfactual structure and need not yield a
principal submatrix. \textcite{spencerwells2019} make the related point that converting a
simultaneous problem into a sequential one can create a different decision problem, and
\textcite{gallow2026} explicitly defines a binary sub-decision by removing the other options
\emph{and conditioning credences on the surviving disjunction}, which is yet another operation.

\subsection{Rectangular and fixed-reference domains}

\begin{definition}[Rectangular inputs]\label{def:rect}
A \emph{rectangular input} is a triple $(J,A,M)$ with $J\subseteq\mathbb L$ finite, $\varnothing\neq A\subseteq J$,
and $M\in\Rl^{J\times A}$. Here $J$ indexes the retained belief perspectives and $A$ the available
options, and option $a\in A$ has its \emph{own row} $a$. Column restriction is
$(J,A,M)|_K=(J,K,M_{J\times K})$ for $\varnothing\neq K\subseteq A$.
\end{definition}

The own-row association is what keeps diagonal optimality, self-sovereignty, and monotonicity
meaningful. The off-menu rows are retained descriptions of the old conditional or counterfactual
laws; they are additional data, not rows of the canonical causal matrix of the smaller problem,
which remains square. Column restriction can therefore describe a fair augmented problem, but it does
not follow from reduction to $M_{D'}$ alone.

The axioms extend to rectangular inputs as follows: relabellings act on $J$ (and hence on $A$);
\ax{R} adds $b_j$ to row $j$ for every $j\in J$; \ax{P} quantifies over all rows in $J$; \ax{D}
maximises $M_{aa}$ over $a\in A$; \ax{S} compares inputs with the same $J$ and $A$; \ax{Q} asks for two
inputs with the same $J$ and $A$ and different choice sets; and contraction becomes
\[
  \Ccol:\qquad G(J,A,M)\cap K\subseteq G(J,K,M_{J\times K}).
\]
Binary independence splits into two inequivalent axioms:
\begin{description}[font=\normalfont,leftmargin=2.4em,labelwidth=1.6em,labelsep=.8em,itemsep=.35em]
  \item[\Bpr] the comparison of $a,b$ agrees with their comparison in the principal square input
        $(\{a,b\},\{a,b\},M_{\{a,b\}\times\{a,b\}})$;
  \item[\Bcol] the comparison of $a,b$ agrees with their comparison in $(J,\{a,b\},M_{J\times\{a,b\}})$,
        which keeps every row.
\end{description}
Both imply \Ccol, by the argument of Lemma~\ref{lem:bimplies}; \Bpr{} implies \ax{S}, while \Bcol{} does
not, because it never compares inputs with different belief universes.

For fixed-reference restriction, supply a probability vector $w\in\Delta(J)$, evaluate
$v_a(M,w)=\sum_jw_jM_{ja}$, and restrict by $(J,A,M,w)|_K=(J,K,M_{J\times K},w)$; under relabelling, $w$
is relabelled too, and uniform weights give a rule that depends on the rectangular matrix alone. A
reference distribution over belief perspectives need not be the agent's probability over its
\emph{available} choices: if it is meant to be, eliminating choices ordinarily requires conditioning,
and keeping positive weight on eliminated choices requires a different interpretation. The next
result shows that the reference cannot be updated while preserving contraction.

\begin{theorem}[Weight preservation]\label{thm:weights}
Let $w,w'$ be probability vectors on the same set $J$, let at least two columns survive a restriction,
and suppose the update $w\mapsto w'$ does not depend on the payoff entries. Weighted averaging ranks
every payoff matrix identically under $w$ and $w'$ if and only if $w'=w$, and contraction holds for
every payoff matrix under the update only if $w'=w$.
\end{theorem}

\begin{proof}
The difference of two columns can be any $x\in\Rl^J$, so identical rankings require
$\operatorname{sign}(w^\top x)=\operatorname{sign}(w'^\top x)$ for all $x$; the two functionals are then
positive multiples of each other, and normalisation gives $w'=w$. If $w'\ne w$, choose $x$ with
$w^\top x=0<w'^\top x$ and replace it by $x-\epsilon\one$ for small $\epsilon>0$, so that $w^\top x<0<w'^\top x$.
Realise $x$ as the difference of two surviving columns $a,b$, so that $b$ is preferred under $w$ and
$a$ under $w'$, and set every removed column low enough that $b$ is a full-menu maximiser under $w$.
After restriction $b$ is no longer chosen, and contraction fails.
\end{proof}

Deleting rows and renormalising their weights therefore preserves all rankings only when the deleted
rows had total weight zero; with full-support weights, any proper row deletion changes the
reference.

\begin{center}
\small
\renewcommand{\arraystretch}{1.25}
\begin{tabularx}{\textwidth}{@{}>{\raggedright\arraybackslash}p{7.2em}
                                  >{\raggedright\arraybackslash}X
                                  >{\raggedright\arraybackslash}X
                                  >{\raggedright\arraybackslash}p{8.4em}@{}}
  \toprule
  operation & available columns & belief rows and reference & canonical law-preserving matrix? \\
  \midrule
  principal & delete unavailable options & delete the corresponding rows & yes, given law preservation \\
  column & delete unavailable options & retain every row & needs augmented reference data \\
  fixed reference & delete unavailable options & retain $J$ and $w$ & needs a supplied reference \\
  delete and renormalise & delete unavailable options & change the reference & a different belief update \\
  \bottomrule
\end{tabularx}
\end{center}

\subsection{Consistency maps under column restriction}

\begin{theorem}[Column map with principal binary independence]\label{thm:colmap}
On all finite rectangular inputs, the inclusion-minimal inconsistent subsets of
\[
  \{\ax{N},\ax{F},\ax{R},\ax{P},\ax{D},\Bpr,\ax{S},\Ccol,\ax{O},\ax{Q}\}
\]
are exactly
\[
  \{\ax{P},\ax{D}\},\quad\{\ax{P},\Bpr\},\quad\{\ax{P},\ax{S}\},\quad\{\ax{R},\ax{D}\},\quad
  \{\ax{N},\ax{R},\Bpr,\ax{Q}\},
\]
and the maximal consistent subsets are
\begin{center}
\small
\begin{tabular}{@{}ll@{}}
  \toprule
  maximal consistent subset & witness \\
  \midrule
  $\{\ax{N},\ax{F},\ax{R},\ax{P},\Ccol,\ax{O},\ax{Q}\}$ & column sums $\sum_{j\in J}M_{ja}$ \\
  $\{\ax{N},\ax{F},\ax{D},\Bpr,\ax{S},\Ccol,\ax{O},\ax{Q}\}$ & diagonal ranking \\
  $\{\ax{N},\ax{F},\ax{R},\ax{S},\Ccol,\ax{O},\ax{Q}\}$ & full-menu relative scores \\
  $\{\ax{F},\ax{R},\Bpr,\ax{S},\Ccol,\ax{O},\ax{Q}\}$ & responsive fixed-label rule \\
  $\{\ax{N},\ax{F},\ax{R},\Bpr,\ax{S},\Ccol,\ax{O}\}$ & universal indifference \\
  \bottomrule
\end{tabular}
\end{center}
All the witnesses also satisfy \Mon.
\end{theorem}

\begin{proof}
The \emph{full-menu relative-score rule} ranks by $s_a=M_{aa}-\frac1{|A|}\sum_{c\in A}M_{ac}$ when $A=J$
and ties every option when $A\subsetneq J$. It is neutral and invariant under \ax{F} and \ax{R}; on full
menus each score depends only on its own row, and inputs compared by \ax{S} have the same $J$ and $A$,
so \ax{S} holds; every proper column restriction produces a non-full menu, on which everything is
chosen, so \Ccol{} holds; and $s_a$ increases with $M_{aa}$. Column sums are unchanged on surviving
columns by column restriction, which gives \Ccol, and shift uniformly under \ax{R}. The other
witnesses are as in Theorem~\ref{thm:augmented}.

The obstructions $\{\ax{P},\ax{D}\}$, $\{\ax{P},\ax{S}\}$, and $\{\ax{R},\ax{D}\}$ are proved as in
Propositions~\ref{prop:PD}, \ref{prop:PS}, and~\ref{prop:NRBQ}, with square inputs; $\{\ax{P},\Bpr\}$
follows from $\Bpr\Rightarrow\ax{S}$; and $\{\ax{N},\ax{R},\Bpr,\ax{Q}\}$ is proved by the three-cycle
of Proposition~\ref{prop:NRBQ}, applied to square inputs. For completeness: a set containing \ax{P}
excludes \ax{D}, \Bpr, and \ax{S}, and the column-sum witness covers it. Without \ax{P} and without
\ax{R}, diagonal ranking covers it. With \ax{R}, \ax{D} is excluded and at least one of \ax{N}, \Bpr,
\ax{Q} is omitted, and the last three witnesses cover these cases.
\end{proof}

The obstruction $\{\ax{P},\ax{C}\}$ of Theorem~\ref{thm:pairwise} has disappeared, and so has
$\{\ax{N},\ax{R},\ax{C},\ax{Q}\}$. The matrix
\[
  W=\begin{pmatrix}0&-1&1\\1&0&-1\\-1&1&0\end{pmatrix}
\]
makes the difference visible. Its column sums tie. Under principal restriction, \ax{P} successively
forces $1$ over $2$, $2$ over $3$, and $3$ over $1$ in the binary sub-menus. Under column restriction,
the third row remains in every binary sub-menu and blocks those Pareto judgements, and column
averaging continues to tie.

\begin{theorem}[Column map with column binary independence]\label{thm:colmap2}
Replacing \Bpr{} by \Bcol, the minimal inconsistent subsets are exactly
\[
  \{\ax{P},\ax{D}\},\qquad\{\ax{P},\ax{S}\},\qquad\{\ax{R},\ax{D}\},
\]
and the maximal consistent subsets are
$\{\ax{N},\ax{F},\ax{R},\ax{P},\Bcol,\Ccol,\ax{O},\ax{Q}\}$ (column sums),
$\{\ax{N},\ax{F},\ax{D},\Bcol,\ax{S},\Ccol,\ax{O},\ax{Q}\}$ (diagonal ranking), and
$\{\ax{N},\ax{F},\ax{R},\Bcol,\ax{S},\Ccol,\ax{O},\ax{Q}\}$ (the binary-universe rule below).
\end{theorem}

\begin{proof}
The \emph{binary-universe rule} ranks by $s_a=M_{aa}-M_{ab}$ when $J=A=\{a,b\}$ and ties everything when
$|J|\ge3$; singleton menus are trivial. It satisfies \ax{Q} on a two-row universe, and \Bcol{} never
connects different belief universes, so it does not force the two-row behaviour to extend to three
rows. Column sums satisfy \Bcol{} because deleting other columns leaves the two sums unchanged. The
containment argument is as before: \ax{P} excludes \ax{D} and \ax{S}; without \ax{P} and \ax{R},
diagonal ranking applies; with \ax{R} and without \ax{P}, the binary-universe rule applies.
\end{proof}

The weakness here lies in the existential form of \ax{Q}. Fix a belief universe $J$ with $|J|\ge3$,
allow all its menus, and let \QJ{} require non-triviality on that universe.

\begin{theorem}[A five-axiom family]\label{thm:five}
$\{\ax{N},\ax{R},\Bcol,\ax{S},\QJ\}$ is a minimal inconsistent family.
\end{theorem}

\begin{proof}
If every binary comparison on $J$ were a tie, \Bcol{} would make every comparison on $J$ a tie, and
every menu would be chosen entire, contradicting \QJ; so some strict binary comparison occurs on $J$.
On a binary menu $\{a,b\}$, \ax{S} makes the comparison depend only on the principal block of $a$ and
$b$, and \ax{R} lets its diagonal be normalised to zero. Neutrality, applied to the relabellings of
$J$ that permute three of its elements cyclically, replicates the strict comparison around a three-cycle
of binary menus, and \Bcol{} transfers the three comparisons to the menu of those three options, all of
which are then strictly beaten. The witnesses omitting \ax{N}, \ax{R}, \Bcol, \ax{S}, \QJ{} are,
respectively, the responsive fixed-label rule, diagonal ranking, full-menu relative scores, column
sums, and universal indifference; they can additionally satisfy \ax{F}, \Ccol, \ax{O}, and \Mon.
\end{proof}

The strongest replacement thus depends on whether non-triviality is required on a universe large
enough to support the obstruction.

\subsection{What survives of the circulant theorem}

\begin{proposition}\label{prop:colaverage}
Uniform column averaging satisfies \ax{N}, \ax{F}, \ax{R}, \ax{P}, \Bcol, \Ccol, \ax{O}, \ax{Q}, and \Mon.
The same holds for weighted averaging with any fixed full-support reference $w$ under $\Cref$ (with
\ax{N} understood to relabel $w$).
\end{proposition}

\begin{proof}
Each option's score is unaffected by the other columns and by their removal; row translations shift
all scores equally; positive weights give \ax{P}; raising $M_{aa}$ raises only the score of $a$.
\end{proof}

Under column or fixed-reference restriction, then, no impossibility can be obtained from any subset
of these axioms beyond those listed in Theorems~\ref{thm:colmap} and~\ref{thm:colmap2}. In
particular, row-baseline-invariant evaluation \emph{can} be menu coherent when the belief model is
held fixed, and the first half of the proposed dichotomy fails. The impossibility of
Theorem~\ref{thm:circulant} concerns the preservation of comparisons when the belief perspectives
associated with removed options disappear, not the preservation of choices when options disappear.

% -----------------------------------------------------------------------------------
\section{What characterises ex-ante evaluation}\label{sec:exante}

The second half of the dichotomy claimed that menu coherence forces ex-ante evaluation. This section
shows that principal binary independence does force it, given symmetry, ordering, and monotonicity,
and that contraction does not.

\subsection{Binary independence: the Self-Esteem theorem}

\begin{theorem}[Relational Self-Esteem theorem]\label{thm:selfesteem}
On the domain of all finite real square matrices, a rule satisfies \ax{N}, \Bpr, and \ax{O} if and only if
there is a complete transitive relation $\succeq^*$ on $\Rl$ such that, for every $M$ and all $i,j$,
\[
  i\succeq_Mj\iff M_{ii}\succeq^*M_{jj}.
\]
\end{theorem}

\begin{proof}
Consider
\[
  H=\begin{pmatrix}p&r&p\\q&s&u\\p&t&p\end{pmatrix}.
\]
The principal block on $\{1,3\}$ is constant, so the transposition of $1$ and $3$ fixes it; by \ax{N} and
\Bpr{} the comparison of $1$ and $3$ is symmetric, and by completeness $1\sim3$. By transitivity, $1$ and
$3$ compare in the same way with $2$. By \Bpr{} the first comparison is that of the first index with the
second in $\bigl(\begin{smallmatrix}p&r\\q&s\end{smallmatrix}\bigr)$, and the second is that of the first
index with the second in the principal block on $(3,2)$, which is
$\bigl(\begin{smallmatrix}p&t\\u&s\end{smallmatrix}\bigr)$. Since $r,q,t,u$ are arbitrary, the comparison
of two options does not depend on the off-diagonal entries of their block, and by \Bpr{} and \ax{N} it is
a function of the ordered pair of diagonal values. Define $x\succeq^*y$ if the first index is weakly
preferred in $\operatorname{diag}(x,y)$. Completeness is inherited, and transitivity follows by applying
\Bpr{} to $\operatorname{diag}(x,y,z)$. Conversely, ranking by a weak order of the diagonal values satisfies
\ax{N}, \Bpr, and \ax{O}, and has nonempty maximal sets on finite menus.
\end{proof}

This is the relational form of the \enquote{Self-Esteem Theorem} of \textcite{podgorski2022}: if
the comparison of two options is a fixed function of the four values $V_X(X)$, $V_Y(Y)$, $V_X(Y)$,
$V_Y(X)$ and overall verdicts respect independence of irrelevant alternatives, the comparison
depends only on each option's value by its own lights, $V_X(X)$ and $V_Y(Y)$. Podgorski concludes
that among overall-value views only the evidential one respects that independence, and his
tournament decision theory escapes the theorem by producing no single overall value. Here the
common comparison function is supplied by \ax{N} and \Bpr, and the scalar overall value by \ax{O};
the relational version assumes no real-valued representation of $\succeq^*$.

\begin{corollary}\label{cor:three}
Adding \ax{F} leaves exactly three rules: increasing diagonal ranking, decreasing diagonal ranking,
and universal indifference. Adding \ax{Q} and \Mon{} leaves only increasing diagonal ranking.
\end{corollary}

\begin{proof}
Every pair $x>y$ is mapped to every other such pair by a positive affine transformation, so under \ax{F}
all unequal pairs are oriented alike; equal values tie by reflexivity and \ax{N}. \ax{Q} excludes
indifference. Under decreasing ranking, raising $M_{11}$ in the zero $2\times2$ matrix removes option $1$
from the choice set, violating the choice clause of \Mon.
\end{proof}

Since $M_{ii}=V_D(i)$,
\[
  \ax{N}+\ax{F}+\Bpr+\ax{O}+\ax{Q}+\Mon
\]
characterises ex-ante policy ranking under the stated matrix semantics. Under the structural axioms
\ax{N}, \ax{F}, \Bpr, \ax{O}, a single binary benchmark with unequal diagonal entries identifies which of the
three rules applies; no finite benchmark, however, certifies that the structural axioms themselves
hold everywhere.

\begin{proposition}[Exact domain threshold]\label{prop:threshold}
For domains consisting of all matrices up to a maximum size, the characterisation holds if and only if
matrices of size three are included.
\end{proposition}

\begin{proof}
With every matrix of sizes one to three, the proof of Theorem~\ref{thm:selfesteem} goes through, and
\Bpr{} propagates the conclusion to any larger sizes present. With sizes at most two, \Bpr{} is vacuous,
and ranking by $s_i=M_{ii}-M_{ij}$ satisfies \ax{N}, \ax{F}, \ax{O}, \ax{Q}, and \Mon; on
$\bigl(\begin{smallmatrix}1&100\\0&0\end{smallmatrix}\bigr)$ it chooses option $2$.
\end{proof}

Merely having some three-option problems is not enough; a sufficient richness condition is that the
domain contain all binary matrices, all the clone extensions $H$ above, and their relabellings (and,
for the affine conclusion, the relevant affine transforms). The argument also works on
$[0,1]$-valued matrices of sizes up to three, with \ax{F} interpreted internally.

\subsection{Contraction does not force the diagonal}

Can contraction replace binary independence? Under every menu semantics, the answer is no.

\begin{theorem}[A contraction-consistent family without diagonal optimality]\label{thm:lambda}
Write $d_i=M_{ii}$. For $\lambda\ge0$ let $i\to_\lambda j$ if $d_i-d_j>\lambda\lvert M_{ij}-M_{ji}\rvert$; choose
the options with no incoming edge, tied at the top, and rank the rejected options below them by their
diagonal entries. Every such rule satisfies \ax{N}, \ax{F}, \Cpr, \Ccol, \ax{O}, \ax{Q}, and \Mon. At
$\lambda=0$ it is increasing diagonal ranking; for every $\lambda>0$ it sometimes recommends an option
whose diagonal entry is not maximal.
\end{theorem}

\begin{proof}
Every edge strictly lowers the diagonal, so the edge relation is acyclic and the choice set is
nonempty. Whether $i\to_\lambda j$ depends only on $M_{ii},M_{jj},M_{ij},M_{ji}$, which survive both principal
and column restriction; so a chosen option acquires no incoming edge under restriction, which gives
both forms of contraction. Raising $d_i$ can only remove edges into $i$ and add edges out of $i$; hence
$i$'s chosen status can only improve, other options' chosen status can only worsen, and if $i$ stays
rejected its rank among the rejected improves, which gives \Mon. \ax{N} and \ax{F} are clear, since
$\lvert M_{ij}-M_{ji}\rvert$ scales with $a$ and is translation invariant. For $\lambda>0$, in
$\bigl(\begin{smallmatrix}1&0\\100/\lambda&0\end{smallmatrix}\bigr)$ there is no edge, and both options
are chosen.
\end{proof}

Even an ordered, neutral, affine-invariant, monotone, contraction-consistent rule, therefore, need not
maximise ex-ante value. The family also defeats finite certification: given finitely many benchmark
matrices, choose $\lambda>0$ so small that every option with non-maximal diagonal in every benchmark
receives an edge from a diagonal maximiser. The whole ranking then coincides with diagonal ranking on
every benchmark, while a new binary matrix with a large off-diagonal asymmetry relative to its
diagonal gap violates \ax{D}.

\begin{proposition}[Maximal cardinality]\label{prop:cardinality}
Let $\cont=\lvert\Rl\rvert$. There are $2^{\cont}$ rules satisfying \ax{N}, \ax{F}, \Cpr, \Ccol, \ax{O}, \ax{Q},
and \Mon.
\end{proposition}

\begin{proof}
For $E\subseteq\Rl$, put $i\to j$ if $d_i>d_j$, $s=M_{ji}-M_{ij}>0$, and either $d_i-d_j>s$, or $d_i-d_j=s$ and
$(d_j-M_{ij})/s\in E$; complete as in Theorem~\ref{thm:lambda}. Edges descend the diagonal and depend only
on principal blocks, so the choice set is nonempty and both contractions hold. Raising $d_i$ only
decreases the gaps $d_k-d_i$, so it never creates an edge into $i$ (at the boundary it may delete one),
and only increases the gaps $d_i-d_j$, so it never deletes an edge out of $i$; this gives \Mon. The
boundary parameter $(d_j-M_{ij})/s$ is invariant under positive affine maps. On
\[
  M_r=\begin{pmatrix}r+1&0\\1&r\end{pmatrix}
\]
we have $d_1-d_2=s=1$ and boundary parameter $r$, so the choice is $\{1\}$ if $r\in E$ and $\{1,2\}$
otherwise. Distinct sets $E$ give distinct rules, and $2^{\cont}$ is the number of all rules on finite real
matrices over a countable label universe.
\end{proof}

The class cannot, therefore, be parameterised by finitely many real numbers, or even countably many.

\subsection{Fixed-reference scores}

Suppose that only the score vector $v_a=\sum_jw_jM_{ja}$ is available, and that restriction deletes
coordinates. Even then rigidity does not return. The rule choosing the top $k\ge2$ scores, including
boundary ties, $G_k(v)=\{a:\#\{b:v_b>v_a\}<k\}$, with the chosen options tied above the others and the
others ranked by score, satisfies \ax{N}, \ax{F}, contraction, \ax{O}, \ax{Q}, and \Mon.

\begin{proposition}\label{prop:scores}
There are $2^{\cont}$ rules on finite score vectors satisfying \ax{N}, \ax{F}, contraction, \ax{O}, \ax{Q}, and
\Mon, already distinguished on vectors of length four.
\end{proposition}

\begin{proof}
Fix a set $E$ of normalised shapes of triples. For a candidate $j$ and a triple $T$ of other options,
let $g=(v_i-v_j)_{i\in T}$. Reject $j$ if some triple has $\sum g>0$, or has $\sum g=0$, $g\ne0$, and
sorted normalised shape in $E$; choose the rest, and complete the ranking as above. A score maximiser
is never rejected, since all its gaps are non-positive. Deletion removes witnesses, so contraction
holds. Raising $v_j$ lowers every gap sum for $j$, so a zero-sum rejection can only turn into a
negative sum, and raises the relevant sums for other candidates; this gives \Mon. Positive affine
maps preserve signs and normalised shapes. On $v=(0,-1,t,1-t)$ with $0<t<\tfrac12$, the option of
score $0$ has the single triple $(-1,t,1-t)$, of sum zero, and is rejected exactly when its shape lies
in $E$; these shapes are distinct for distinct $t$.
\end{proof}

With a fixed full-support $w$, raising $M_{jj}$ raises only $v_j$, so the matrix form of \Mon{} follows.
The upshot is that the dichotomy fails in both directions: contraction under any of the three
semantics is compatible with a rich family of non-ex-ante rules, and row-baseline-invariant
evaluation is compatible with contraction once the belief model is fixed. What does force ex-ante
evaluation is principal binary independence, which is a strong requirement: it imposes on every
comparison the beliefs that the two compared options would induce, and no others.

% -----------------------------------------------------------------------------------
\section{Types, parameters, and certification}\label{sec:taxonomy}

A recurring hope in the debate is that the relevant problem types are few: that a decision theory
could be settled by its behaviour on a small catalogue of representative problems. Whether this is
so depends on what counts as the same type, and on which structural axioms a theory is already
known to satisfy.

\subsection{Problems}

\emph{Boolean classifications are automatically finite.} For $k$ local criteria, the profile
$\mathcal S(D)=\{I\subseteq\{1,\dots,k\}:\text{some policy satisfies every criterion in }I\}$ takes at most
$2^{2^k}$ values. A finite classification by joint-satisfiability profiles thus exists trivially; it
need not help to solve any problem, and \cri{4}, being a relation between problems, is not one of
these local predicates.

\emph{Exact recommendation structures are not finite.} Call two problems equivalent if a bijection of
their policy spaces preserves the set of policies acceptable under each criterion.

\begin{proposition}\label{prop:infinitetypes}
Even with two actions and one always-reached information set, the continuous mixed model contains
infinitely many pairwise inequivalent problems.
\end{proposition}

\begin{proof}
For distinct $r_1,\dots,r_n\in(0,1)$, let a reward $Y\in\{0,1\}$, independent of the realised action,
satisfy $\Pr(Y=1\mid p)=1-\prod_i(p-r_i)^2$, a continuous probability since each factor lies in $[0,1]$.
The value is maximised exactly at $p\in\{r_1,\dots,r_n\}$, so exactly $n$ policies satisfy \cri{1}, while
every policy satisfies \cri{2} and \cri{3}. A bijection preserving the \cri{1}-sets cannot identify
problems with different $n$.
\end{proof}

For deterministic policies, Theorem~\ref{thm:realise} already shows that every finite recommendation
structure occurs. The permissive definition of a theory has a further, direct consequence.

\begin{proposition}[No finite benchmark suite]\label{prop:benchmark}
No finite set of problems can establish that an arbitrary decision theory satisfies \cri{1}
everywhere.
\end{proposition}

\begin{proof}
Modify $T^{\mathrm{ex}}$ only on one two-action problem outside the suite, with a strictly better
action, so as to select the worse action. The two theories agree on every benchmark, and only one
satisfies \cri{1} everywhere.
\end{proof}

This holds even if the benchmarks and the omitted problem are of the same intuitive type.
Representative testing requires an invariance requirement on theories. Note also that infinitely
many equivalence classes do not by themselves preclude a finite-dimensional parameterisation: a
single real parameter indexes infinitely many cases.

\subsection{Rules}

The question is sharper for rules on matrices, where invariance axioms are available. Finitely many
parameters must then be read cardinally: a finite-dimensional input space does not make an
arbitrary rule on it a finite-dimensional family.

\begin{proposition}\label{prop:NFO}
There are $2^{\cont}$ rules satisfying \ax{A0}, \ax{N}, \ax{F}, and \ax{O}. Finite benchmark suites cannot
certify \ax{P} or \ax{D} for them.
\end{proposition}

\begin{proof}
For $0<t<1$ let $N_t=\bigl(\begin{smallmatrix}0&-1\\t+1&t\end{smallmatrix}\bigr)$. These lie in distinct
orbits of relabelling combined with positive affine maps, and each has trivial stabiliser: without a
swap, preserving the entries $0$ and $-1$ forces the identity map; with a swap, the diagonal $(t,0)$
must be mapped to $(0,s)$, which requires a negative scale. For each $E\subseteq(0,1)$, choose option $1$
on $N_t$ when $t\in E$ and option $2$ otherwise, transport this ranking equivariantly over each orbit,
and use diagonal ranking elsewhere. This gives $2^{\cont}$ distinct rules, whereas finite tuples of reals
form a set of cardinality $\cont$. In $N_t$, column $1$ strictly dominates column $2$ and the second
diagonal entry is larger; so modifying column sums, or diagonal ranking, on one untested orbit
produces a violation of \ax{P}, or of \ax{D}, that no finite suite (even an adaptive one) detects.
\end{proof}

The construction uses the permission to have arbitrary lookup rules; it does not carry over unchanged
to classes with continuity or measurability requirements.

\begin{proposition}\label{prop:NFB}
There are $2^{\cont}$ rules satisfying \ax{N}, \ax{F}, and \ax{B} (without \ax{O}), and finite suites cannot
certify \ax{D} for them.
\end{proposition}

\begin{proof}
For $E\subseteq\Rl$ and $M_{ii}>M_{jj}$, declare $i\succ_Mj$ if $(M_{ij}-M_{jj})/(M_{ii}-M_{jj})\in E$; never prefer a
lower diagonal strictly to a higher one, and tie all other pairs. The ratio is invariant under positive
affine maps and depends only on the principal block, and strict comparisons descend the diagonal, so
choice sets are nonempty. On $M_t=\bigl(\begin{smallmatrix}1&t\\0&0\end{smallmatrix}\bigr)$ the choice is
$\{1\}$ if $t\in E$ and $\{1,2\}$ otherwise. Starting from $E=\Rl$ (diagonal ranking) and removing one
untested ratio defeats any finite suite.
\end{proof}

\begin{proposition}\label{prop:NFSO}
There are $2^{\cont}$ rules satisfying \ax{N}, \ax{F}, \ax{S}, and \ax{O} that agree on every binary
matrix.
\end{proposition}

\begin{proof}
Let $r_i=(M_{ii}-\min_jM_{ij})/(\max_jM_{ij}-\min_jM_{ij})$, with $r_i=0$ for a constant row, and for
$E\subseteq(0,1)$ rank by $s_i=\one\{r_i\in E\}$. These rules satisfy \ax{N}, \ax{F}, \ax{S}, and \ax{O}, even
under separate positive affine transformations of rows. On binary matrices every $r_i$ is $0$ or $1$,
so every score is zero; on
$\Bigl(\begin{smallmatrix}t&0&1\\1&0&1\\1&1&0\end{smallmatrix}\Bigr)$ the scores are
$(\one\{t\in E\},0,0)$.
\end{proof}

Together with Corollary~\ref{cor:three}, these results give the following picture.

\begin{center}
\small
\renewcommand{\arraystretch}{1.25}
\begin{tabularx}{\textwidth}{@{}>{\raggedright\arraybackslash}p{9em}
                                  >{\raggedright\arraybackslash}X@{}}
  \toprule
  assumptions beyond \ax{A0} & size of the family; certification \\
  \midrule
  \ax{N}, \ax{F}, \ax{O} & $2^{\cont}$; finite suites certify neither \ax{P} nor \ax{D} \\
  \ax{N}, \ax{F}, \ax{B} & $2^{\cont}$; finite suites do not certify \ax{D} \\
  \ax{N}, \ax{F}, \ax{S}, \ax{O} & $2^{\cont}$; all binary behaviour can leave the rule undetermined \\
  \ax{N}, \ax{F}, \ax{C}, \ax{O}, \ax{Q}, \Mon & $2^{\cont}$; finite suites do not certify \ax{D} (Section~\ref{sec:exante}) \\
  \ax{N}, \ax{F}, \ax{B}, \ax{O} & exactly three rules; one unequal-diagonal test distinguishes them \\
  \bottomrule
\end{tabularx}
\end{center}

\enquote{The relevant types are few} is thus true in exactly one of these regimes, and there only
conditionally: when principal binary independence and ordering are known to hold globally, a single
binary test decides among the three possible rules. Finite observation never verifies those
structural axioms themselves.

% -----------------------------------------------------------------------------------
\section{Mixed policies and evidential semantics}\label{sec:mixed}

\subsection{A canonical mixed extension}

A pure causal matrix does not determine how prediction responds to a mixed policy, and continuity
alone does not specify the extension. One precise extension takes the predictor's response to be
linear in the policy and the action to be drawn independently of the response given the announced
distribution:
\[
  \mu_p=\sum_ip_i\mu_i,\qquad g(p)=p^\top M,\qquad V(p)=p^\top Mp ,
\]
where $g(p)_j$ is the causal value of option $j$ under the beliefs induced by $p$. This differs from
a predictor of the realised sample, which can instead yield the linear value $\sum_ip_iM_{ii}$. Define
\begin{description}[font=\normalfont,leftmargin=2.4em,labelwidth=1.6em,labelsep=.8em,itemsep=.35em]
  \item[$\ax{D}_m$] $p$ maximises $p^\top Mp$;
  \item[$\ax{J}_m$] every action in the support of $p$ maximises $p^\top M$;
  \item[$\ax{P}_m$] $p$ assigns probability zero to every weakly dominated column.
\end{description}

\begin{proposition}\label{prop:mixedPJ}
In the canonical extension, $\ax{P}_m$ and $\ax{J}_m$ are jointly satisfiable on every matrix.
\end{proposition}

\begin{proof}
For $0<\eta<1/|I|$ let $\Delta_\eta=\{p:p_i\ge\eta\text{ for all }i\}$. The correspondence
$p\mapsto\argmax_{q\in\Delta_\eta}p^\top Mq$ on the compact convex set $\Delta_\eta$ has nonempty convex
values and a closed graph, so Kakutani's theorem gives
$p^\eta\in\argmax_{q\in\Delta_\eta}(p^\eta)^\top Mq$. If column $a$ weakly dominates column $b$, full support
makes $(p^\eta)^\top Me_a>(p^\eta)^\top Me_b$, so $p^\eta_b=\eta$; otherwise moving mass from $b$ to $a$ would
improve the response. Let $p$ be a limit point as $\eta\downarrow0$. It assigns zero mass to every
dominated column. For any $q$, put $q^\eta=(1-|I|\eta)q+\eta\one\in\Delta_\eta$; then
$(p^\eta)^\top Mp^\eta\ge(p^\eta)^\top Mq^\eta$, and in the limit $p^\top Mp\ge p^\top Mq$. Hence $p$ is supported on
best responses to itself.
\end{proof}

\begin{proposition}\label{prop:mixedconflict}
$\{\ax{P}_m,\ax{D}_m\}$ and $\{\ax{J}_m,\ax{D}_m\}$ are inconsistent.
\end{proposition}

\begin{proof}
For $M=\bigl(\begin{smallmatrix}2&3\\0&1\end{smallmatrix}\bigr)$ and $p=(x,1-x)$, $V(p)=1+x$, so $\ax{D}_m$ selects
$x=1$, while column $2$ dominates and is the unique best response to every $p$, so $\ax{P}_m$ and
$\ax{J}_m$ select $x=0$.
\end{proof}

\begin{proposition}\label{prop:mixedN}
For rules that select a single lottery, $\ax{N}$ and $\ax{D}_m$ are inconsistent.
\end{proposition}

\begin{proof}
For $M=I_2$, the swap fixes $M$, so a neutral single lottery is $(\tfrac12,\tfrac12)$; but
$V(p)=p_1^2+p_2^2$ is maximised only at the two pure endpoints.
\end{proof}

A neutral \emph{correspondence} of maximising lotteries exists; a single neutral maximising lottery
need not. The pure consistency maps therefore cannot simply be relabelled as mixed maps.

\subsection{Contraction for lotteries}

For a rule $L$ that selects one lottery, the natural strengthening of contraction is
conditional-renormalisation consistency:
\[
  \ax{C}_m:\qquad L(M|_K)=\frac{L(M)|_K}{L(M)(K)}\quad\text{whenever }L(M)(K)>0 .
\]

\begin{proposition}\label{prop:mixedcirculant}
$\{\ax{N},\ax{R},\ax{C}_m,\ax{Q}\}$ is a minimal inconsistent family for single-lottery rules.
\end{proposition}

\begin{proof}
In the circulant extension of Theorem~\ref{thm:circulant}, neutrality forces the uniform lottery;
conditioning forces the uniform lottery on every embedded normalised matrix; row translations then
force the uniform lottery everywhere. For minimality, take the uniform distribution over the
choice set of, respectively, the responsive fixed-label rule, diagonal ranking, row-range ranking,
and universal indifference. The first two maximise a binary-independent weak order, so that
$G(M|_K)=G(M)\cap K$ whenever the right-hand side is nonempty, which is exactly what $\ax{C}_m$ requires of
uniform lotteries; row-range ranking is the witness that omits $\ax{C}_m$.
\end{proof}

Support contraction alone is much weaker: $p_i\propto1+r_i/r_{\max}$, with $r_i$ the row ranges and the
uniform lottery when $r_{\max}=0$, has full support everywhere, satisfies \ax{N}, \ax{F}, \ax{R}, and is
nonconstant. A complete mixed consistency map requires choosing among genuinely different output and
contraction semantics; the results above are the extensions established under the stated ones.

\subsection{Two evidential semantics}\label{sec:evsem}

The matrix does not supply an evidential criterion until the conditioning experiment is identified.
Two natural choices give different answers.

\emph{Reference-policy law.} With a full-support reference $\rho$ and $P(I=i,s)=\rho_i\mu_i(s)$,
conditioning on the policy index $i$ yields $\mu_i$, and the evidential value of $i$ is
$\E[U(s,i)\mid I=i]=M_{ii}$. This semantics gives diagonal optimality, and it is the matrix form of
\cri{6}$^\kappa$ on single-choice problems.

\emph{Conditioning on the realised action under a known mixed policy.} In the canonical extension the
environment has law $\mu_p$, and observing a positive-probability realised action does not change it,
so its conditional value is $g_i(p)$. With an explicit completion for zero-probability alternatives,
the resulting support-optimality condition is ratifiability.

The two condition on different variables under different joint laws, and neither can be silently
substituted for the other. This is one precise sense in which \enquote{evidential decision theory}
names a family rather than a criterion, as the literature on the tickle defence and on ratifiability
has long made clear \autocite{eells1982,jeffrey1983}.


\part{Procedures that read descriptions}\label{part:programs}
% -----------------------------------------------------------------------------------
\section{Problems as programs}\label{sec:programs}

Every result so far has been proved on problems that \emph{supply} their counterfactuals. A problem
of Section~\ref{sec:model} comes with the law $\mu_\pi$ for every policy $\pi$, and with structural
plans that fix what each response would be on every branch. The decision theories that the debate
compares disagree precisely about where such laws come from: causal, evidential, and functional
theories differ over which parts of the world would have been different had the agent chosen
otherwise. A model that hands this dependence to the agent has already settled what the theories
dispute. This is why ex-ante optimality came so cheaply in Theorem~\ref{thm:universal}, and why the
deeper results of Part~\ref{part:across} concerned the shape of an evaluation rule rather than the
identification of alternatives.

This part removes the supply. An agent is now a computable procedure that receives a description of
its situation---by default the source code of the environment---and must itself determine what
would happen under each of its possible outputs. The question is whether some such procedure is
optimal on every fair problem. Answering it requires saying what \enquote{optimal} means once the
comparison is no longer given by a supplied law, and that is where most of the work lies.

\subsection{Agents and environments}

Fix an effective language of \emph{descriptions}, finite strings from which a finite nonempty set
$A(x)$ of legal actions can be computed for each description $x$. For the negative results a fixed
binary action set suffices, so nothing depends on how defaults are chosen when action sets vary.

\begin{definition}[Agents]\label{def:agent}
An \emph{agent} is a partial computable function $T$ from descriptions to actions with
$T(x)\in A(x)$ whenever $T(x)$ is defined. It is \emph{globally total} if it halts on every
well-formed description, and \emph{promise-total} on a class $\cls$ of environments if it halts on
every description belonging to $\cls$.
\end{definition}

The two notions differ. \enquote{A total computable function} normally means the first;
\enquote{an agent that halts on every problem in the class} permits the second. Both receive
answers below, and the distinction turns out to separate a positive theorem from a negative one.

\begin{definition}[Program environment]\label{def:progenv}
A \emph{program environment} $E$ is a program with a finite set $H_E$ of information sets, each
$h\in H_E$ carrying a finite nonempty action set $A_{E,h}$ and an \emph{owned query} $q(E,h)$, a
description from which $E$ and $h$ can be recovered. Its policy space is
$\Pol_E=\prod_{h\in H_E}A_{E,h}$, and its root description $d(E)$ is one of its owned queries. Given
access to an agent, $E$ produces a payoff. For a table $\pi\in\Pol_E$, $E^\pi$ denotes the execution
in which every owned query is answered by the corresponding entry of $\pi$.
\end{definition}

Four stipulations make owned queries behave like the information sets of
Definition~\ref{def:problem}: a query identifies its parent environment and its information set;
distinct environments have disjoint owned-query namespaces; repeated queries receive the same
answer; and, when the fairness condition $\cls_{\mathrm{self}}$ below is imposed, the environment
consults the agent only through its owned queries. The disjointness matters: if two environments
may both call one query their own, the conflict of Proposition~\ref{prop:foreign} below reappears
inside a class that is ostensibly self-only.

Three representations of payoffs must be kept apart. An execution may return an \emph{exact
rational}, or an explicit finite distribution with rational probabilities and payoffs; it may return
a \emph{computable real}, given by a program producing rational approximations of any required
accuracy; or it may be a random execution returning a \emph{sampled} payoff whose expectation is
the value. The positive results use the first; the strongest negative result
(Theorem~\ref{thm:exactsel}) uses only sampled payoffs in $\{0,1\}$, so it does not exploit exotic
real-valued outputs.

\subsection{Replacement and anchored evaluation}

To ask whether $T$ is optimal on $E$ one must say what $E$ would pay another agent $A$. Some ways of
saying it define a single payoff function $U(E,A)$ common to all agents. Others first run or inspect
a reference agent $T$, freeze part of the resulting situation, and only then consider alternatives;
these define an \emph{anchored kernel} $U_K(E,T;A)$, where $K$ names the convention. Every kernel is
required to satisfy the baseline identity
\begin{equation}\label{eq:baseline}
  U_K(E,T;T)=U(E,T),
\end{equation}
and the corresponding notion of optimality is
\begin{equation}\label{eq:anchored}
  U(E,T)\ge U_K(E,T;A)\qquad\text{for every admissible agent }A.
\end{equation}
In general there is no function $W$ with $U_K(E,T;A)=W(E,A)$ for every reference $T$. Some versions
of the question therefore ask for \emph{the best agent in a common environment}, others for \emph{an
agent that is a fixed point of its own counterfactual evaluations}. The difference is the
program-level form of the difference between \cri{1} and \cri{3}.

\begin{definition}[Replacement conventions]\label{def:conventions}
Let $E$ be an environment with a designated \emph{action port}, the place at which the agent's
actual action on $d(E)$ is received.
\begin{description}[font=\normalfont,leftmargin=2.4em,labelwidth=1.6em,labelsep=.8em,itemsep=.35em]
  \item[$O_{\mathrm{const}}$] Replace the action at the port by $A$'s action, and freeze every other
        response of the environment at its value under the reference $T$; fixed embedded algorithms
        are left unchanged.
  \item[$O_{\mathrm{oracle}}$] Replace every designated oracle call, the port included, by a call to
        $A$.
  \item[$O_{\mathrm{syn}}$] Replace specified, typed occurrences of program components whose syntax
        equals the reference program.
  \item[$O_{\mathrm{ext}}$] Replace specified program components that pass a stated test of
        extensional equivalence with the reference agent (Section~\ref{sec:extensional}).
\end{description}
Under $O_{\mathrm{syn}}$ and $O_{\mathrm{ext}}$ the description $d(E)$ supplied to the agent is held
fixed during replacement.
\end{definition}

The last clause is a choice. Supplying the rewritten source $d(E[A])$ instead would change the
agent's observation as well as the program occupying the agent's role; that is a legitimate model,
but a different one (Proposition~\ref{prop:rawsource} shows one of its consequences). Raw
substring replacement is not an adequate reading of $O_{\mathrm{syn}}$, since it can alter quoted
data, comments, or unrelated identifiers; typed occurrences or explicit component interfaces are
needed. Logical counterfactuals, $O_{\mathrm{log}}$, form a family of semantics rather than a fifth
replacement operation, and are treated in Section~\ref{sec:logical}.

\subsection{Fairness}

The fairness condition of Definition~\ref{def:fair} splits into several conditions, which overlap
rather than partition the environments.
\begin{description}[font=\normalfont,leftmargin=2.4em,labelwidth=1.6em,labelsep=.8em,itemsep=.35em]
  \item[$\cls_{\mathrm{self}}$] Dependence on the evaluated agent is confined to its answers on the
        owned queries $q(E,h)$.
  \item[$\cls_{\mathrm{any}}$] The environment may consult the agent on arbitrary descriptions.
  \item[$\cls_{\mathrm{ext}}$] Replacing agents by extensionally equal agents preserves the payoff
        kernel: $T\equiv T'$ and $A\equiv A'$ imply $U_K(E,T;A)=U_K(E,T';A')$.
  \item[$\cls_{\mathrm{const}}$] Fixed algorithms, including copies of a particular agent, may occur
        as constants of the environment.
\end{description}
Two locality requirements should also be separated. \emph{Executed-query locality} says that actual
executions call agents only on owned queries. \emph{Kernel locality} says that the whole
counterfactual kernel, including the rule that decides what to replace, depends on the reference
agent only through its answers on owned queries. Global program-equivalence matching
(Section~\ref{sec:extensional}) satisfies the first while violating the second.

In Section~\ref{sec:model} a fixed algorithm inside a problem was harmless: its output on the
relevant input is a number that is part of the problem, and the problem is a finite object
specified by such numbers. In a program environment the algorithm's code is part of the description
and its output is not; the agent must compute it, and may be the very thing being computed. This is
what makes $\cls_{\mathrm{const}}$ consequential here (Theorem~\ref{thm:fixedcopy}).

\subsection{Programs and causal matrices}

The conventions $O_{\mathrm{oracle}}$ and $O_{\mathrm{const}}$ have exact counterparts in the
causal matrix of Section~\ref{sec:matrix}. Let $E$ be a finite oracle environment in which every
execution with table answers halts with an exact rational payoff, and let the owned queries be
answered at two kinds of position: the action ports, and the remaining oracle calls. For
$\sigma,\tau\in\Pol_E$ let $M_E[\sigma,\tau]$ be the payoff of the execution in which the oracle
calls are answered by $\sigma$ and the ports by $\tau$.

\begin{proposition}[Programs induce causal matrices]\label{prop:progmatrix}
For agents $T$ and $A$ write $\pi_T$ and $\pi_A$ for the tables of their answers on the owned
queries. Then
\[
  U_{\mathrm{oracle}}(E,A)=M_E[\pi_A,\pi_A],\qquad
  U_{\mathrm{const}}(E,T;A)=M_E[\pi_T,\pi_A].
\]
$M_E$ is the causal matrix of the finite problem whose structural states are the oracle tables,
with $\mu_\pi=\delta_\pi$ and $U(\sigma,\tau)=M_E[\sigma,\tau]$, and every finite real square matrix
indexed by a product of finite action sets arises in this way.
\end{proposition}

\begin{proof}
Under $O_{\mathrm{oracle}}$ every owned query, port or not, is answered by $A$, which is the
execution indexed by $(\pi_A,\pi_A)$. Under $O_{\mathrm{const}}$ the ports receive $A$'s answers
while every other response is frozen at its value under $T$; since the only responses depending on
the agent are the oracle calls, this is the execution indexed by $(\pi_T,\pi_A)$. The problem
described has causal matrix $\sum_\rho\mu_\sigma(\rho)U(\rho,\tau)=M_E[\sigma,\tau]$. For the last
claim, an environment that calls the oracle on all owned queries, then reads the ports, and pays a
prescribed entry $M[\sigma,\tau]$ of a given matrix realises $M$, exactly as in
Theorem~\ref{thm:realise}.
\end{proof}

Oracle substitution thus recovers the diagonal of the matrix, ex-ante value, and port-only
intervention recovers a row, causal evaluation at the baseline $\pi_T$. Both require a designation
of which positions are ports and which are calls; a program in which the agent's code occurs inside
other programs, with no designated interface, does not carry that designation, and the agent must
construct it. Section~\ref{sec:operators} shows that the two kinds of evaluation reappear, in this
form, as the only two coherent ways of doing so.

\subsection{Comparing the conventions}

Replacing more components does not make optimality logically stronger or weaker.

\begin{proposition}[No implications between conventions]\label{prop:nonimplication}
For any two of $O_{\mathrm{const}}$, $O_{\mathrm{oracle}}$, $O_{\mathrm{syn}}$, $O_{\mathrm{ext}}$ there
are terminating environments on which a given agent is optimal under the first and not under the
second, and environments on which the reverse holds.
\end{proposition}

\begin{proof}
Let a typed environment contain an action port, an additional oracle call, a literal copy of the
reference agent's code, and an independently encoded copy computing the same function, each
queried at $d(E)$; let the reference agent be the constant $0$ and the alternative the constant $1$.
The four conventions then produce the following patterns of the four values.
\begin{center}
\small
\begin{tabular}{@{}lcccc@{}}
  \toprule
  evaluation & port & oracle call & literal copy & encoded copy \\
  \midrule
  baseline                & 0 & 0 & 0 & 0 \\
  $O_{\mathrm{const}}$ alternative  & 1 & 0 & 0 & 0 \\
  $O_{\mathrm{oracle}}$ alternative & 1 & 1 & 0 & 0 \\
  $O_{\mathrm{syn}}$ alternative    & 1 & 1 & 1 & 0 \\
  $O_{\mathrm{ext}}$ alternative    & 1 & 1 & 1 & 1 \\
  \bottomrule
\end{tabular}
\end{center}
For two conventions producing distinct patterns $v$ and $w$ from the baseline pattern $z$, set the
payoff to $1$ at $z$, $0$ at $v$, $2$ at $w$, and $0$ elsewhere. The reference agent is optimal under
the first convention and not under the second; exchanging the payoffs at $v$ and $w$ gives the
reverse.
\end{proof}

The conventions coincide only under conditions that make their replacement sets coincide.
Oracle and syntactic substitution agree under a slot-preserving compilation of oracle calls into
designated code occurrences (Section~\ref{sec:syntactic}); oracle and extensional substitution agree
when the extensional test selects exactly the designated dependencies; syntactic and extensional
substitution agree when their replacement masks coincide. For logical semantics no comparison is
possible until a particular semantics is fixed: a guarded provability semantics is partial
(Definition~\ref{def:guarded}), a proof of the value of each policy-substituted execution simply
evaluates an independently supplied substitution (Proposition~\ref{prop:proofopt}), and the scores
of a logical inductor are finite-time estimates. Treating these as one fifth payoff function would
conceal the main issue.

% -----------------------------------------------------------------------------------
\section{Computable optimisation and its limits}\label{sec:computation}

This section asks what changes when the decision theory of Theorem~\ref{thm:universal} must be a
single algorithm reading descriptions. With a fixed substitution semantics and a termination promise,
nothing essential changes: one algorithm is optimal everywhere. Three kinds of obstruction appear
once the promise or the class is enlarged, and none of them concerns counterfactual inference.

\subsection{A uniform optimiser}

\begin{definition}[Promise class]\label{def:promise}
$\mathcal P_{\mathbb Q}$ is the class of program environments $E$ such that $E^\pi$ halts and returns an
exact rational value $V_E(\pi)$ for every $\pi\in\Pol_E$.
\end{definition}

The definition uses table oracles, so the validity of an environment is not defined in terms of
whichever optimiser is eventually constructed.

\begin{theorem}[Uniform optimisation on the promise class]\label{thm:promiseopt}
There is an agent that is promise-total on $\mathcal P_{\mathbb Q}$ and optimal under
$O_{\mathrm{oracle}}$ on every $E\in\mathcal P_{\mathbb Q}$.
\end{theorem}

\begin{proof}
On input $q(E,h)$, let $T^*$ recover $E$ and its finite policy space, run $E^\pi$ for every
$\pi\in\Pol_E$, select the first maximiser $\pi^*_E$ of $V_E$ in a fixed effective order, and return
$\pi^*_E(h)$. Every simulation halts by the promise, there are finitely many, and exact rational
comparison terminates. All owned queries of $E$ recover the same environment and use the same
tie-breaking, so the answers form the single policy $\pi^*_E$. Any agent $A$ induces the policy
$\pi^A_E(h)=A(q(E,h))$, and under oracle replacement
\[
  U(E,A)=V_E(\pi^A_E)\le\max_{\pi\in\Pol_E}V_E(\pi)=V_E(\pi^*_E)=U(E,T^*).\qedhere
\]
\end{proof}

$T^*$ never asks its live self what it would do during a simulated execution; it supplies table
answers instead. This is Theorem~\ref{thm:universal} made uniform: the lookup table is replaced by
an algorithm, and no knowledge of which parts of $E$ are predictors is needed, because oracle
substitution designates them. By Proposition~\ref{prop:progmatrix} the theorem maximises the
diagonal of $M_E$.

An explicit finite policy header is sufficient for the argument but not necessary. A finitely
branching oracle machine that halts for \emph{every} total oracle has a finite tree of consistent
oracle answers, since otherwise Kőnig's lemma \autocite{konig1927} gives an infinite path, whose
answers extend to a total oracle on which the machine runs forever; exhaustive exploration of the
tree then yields a finite optimisation problem. Halting against every \emph{computable} oracle does
not justify this compactness argument, since an infinite computable tree can lack a computable
infinite path.

Dependence of termination on the oracle is harmless as long as every policy terminates. The theorem
fails when some alternatives are undefined: if action $0$ returns $0$ and action $1$ simulates
machine $e$ and returns $1$ if it halts, the environment lies outside $\mathcal P_{\mathbb Q}$ unless
$e$ halts. Adding the convention that divergence pays $-1$ would make an exact optimiser choose $1$
exactly when $e$ halts, and so decide the halting problem. The obstruction then comes from the
added divergence semantics, not from the finite-policy argument.

\subsection{Two baseline impossibilities}

Two impossibilities are available with no counterfactual subtlety at all. Any result that claims
to be about counterfactual inference must avoid them.

\begin{proposition}[Foreign queries]\label{prop:foreign}
There are two terminating environments such that no agent, computable or not, is optimal on both
under $O_{\mathrm{oracle}}$.
\end{proposition}

\begin{proof}
Fix a binary query $q$ distinct from both environments' own descriptions, and let $E_0$ pay
$\one[A(q)=0]$ and $E_1$ pay $\one[A(q)=1]$. Each has optimal value $1$, attained by a constant agent,
and no agent attains $1$ in both.
\end{proof}

This is a cross-problem consistency contradiction requiring neither self-reference nor
undecidability; it is excluded by $\cls_{\mathrm{self}}$ with disjoint namespaces, and it
disappears under $O_{\mathrm{const}}$, which freezes the foreign response.

\begin{theorem}[The fixed-copy diagonal]\label{thm:fixedcopy}
If admissible agents are globally total and arbitrary fixed computable algorithms may occur in
environments, no agent is optimal on every environment, even among deterministic, binary,
rational-valued environments that consult the evaluated agent only on their own description.
\end{theorem}

\begin{proof}
Fix a globally total candidate $T$, and construct an environment $E_T$ containing the fixed algorithm
$g_0=T$. With $d=d(E_T)$, it computes $c=g_0(d)$, receives the evaluated agent's action $a$ at the
port, and pays $\one[a\neq c]$. The algorithm $g_0$ does not change when the evaluated agent changes.
Since $T$ is globally total, $g_0(d)$ halts, so $E_T$ terminates under either action. $T$ takes the
action $c$ and receives $0$; a globally total agent taking $1-c$ at $d$ receives $1$.
\end{proof}

In a programming formalism with no built-in parameter for the environment's own description,
Kleene's second recursion theorem \autocite{kleene1938} supplies the self-reference; see
\textcite{moschovakis2010} for a modern statement and proof. It produces a program with a specified
behaviour involving its own index; it does not prove that the behaviour terminates, and global
totality is what supplies termination. The proof holds under $O_{\mathrm{const}}$, and under
$O_{\mathrm{oracle}}$ whenever $g_0$ is an ordinary fixed program rather than a designated oracle
call. The transparent-and-opaque-box construction of \textcite{skalse2021} is the same mechanism with
different payoffs.

Remark~\ref{rem:skalse} found that construction ill defined in the model of
Section~\ref{sec:model}, and the contrast is instructive. There a problem was a finite object whose
data included the constant's output, so the construction required a problem $D_b$ with $b=T(D_b)$,
and an ex-ante optimiser has no such fixed point. Here the description of $E_T$ encodes a
computation involving $T$ without supplying its result, and no fixed point is needed: whatever $T$
outputs on $d(E_T)$, the copy outputs the same. Once that output $c$ is known, the environment is
the ordinary one-shot problem with payoffs $\one[a\ne c]$, and its maximiser is $1-c$. The agent
cannot uniformly select it while remaining globally total, because finding $c$ means running $T$ on
the very description it is reading.

There is no conflict with Theorem~\ref{thm:promiseopt}, because $T^*$ is only promise-total. Its
attempted diagonal environment contains a call to $T^*(d(E_{T^*}))$, whose evaluation invokes the
same optimisation again; the regress violates the table-termination promise, and the environment
is not in $\mathcal P_{\mathbb Q}$. The positive theorem supplies no globally total extension of $T^*$.

A third baseline is the program form of Death in Damascus.

\begin{proposition}[Self-prediction held fixed]\label{prop:doublecall}
Let an oracle environment call the agent twice on $d$, treating one answer as a prediction $p$ and the
other as the action $a$, and pay $\one[a\neq p]$. Under $O_{\mathrm{oracle}}$ every agent receives $0$
and every agent is optimal. Under $O_{\mathrm{const}}$ every agent has an improving deviation.
\end{proposition}

\begin{proof}
Coherent substitution answers both calls alike. An intervention changing only the action, with the
prediction frozen at the reference agent's answer $p$, earns $1$ with $a=1-p$.
\end{proof}

By Proposition~\ref{prop:progmatrix}, this is the matrix
$\bigl(\begin{smallmatrix}0&1\\1&0\end{smallmatrix}\bigr)$, in which no diagonal entry is maximal in
its row; it is Proposition~\ref{prop:did} again, a ratifiability obstruction and not a computability
result.

\subsection{Exact selection}

The following result avoids every baseline: there are no predictors, no copies of the agent, no
foreign queries, and the agent need halt only on valid inputs.

\begin{theorem}[No computable exact optimiser]\label{thm:exactsel}
No promise-total agent is an exact expected-payoff optimiser on all almost-surely terminating
binary-choice environments, even if realised payoffs lie in $\{0,1\}$, the agent is consulted only at
the action port on its own description, every expected value is uniformly computable, and every
environment has finite expected running time.
\end{theorem}

\begin{proof}
Let $(\varphi_e)$ be an effective enumeration of the partial computable functions, and let
$K_b=\{e:\varphi_e(e)\text{ halts with output }b\}$ for $b\in\{0,1\}$. The environment $E_e$ receives
the action $a\in\{0,1\}$, draws $N\in\mathbb N$ with $\Pr(N=n)=2^{-n-1}$, simulates $\varphi_e(e)$ for $N$
steps, and pays $1$ if the simulation halts within them with output $a$, and $0$ otherwise. Every
simulation is bounded, $E_e$ terminates almost surely, and its expected running time is finite. If
$e\in K_b$ with halting time $t$, then $V_e(b)=\Pr(N\ge t)=2^{-t}>0$ and $V_e(1-b)=0$; if neither
output occurs, both values are $0$.

Suppose a promise-total exact optimiser $T$ existed. Then $f(e)=T(d(E_e))$ is a total computable
function with $f(e)=b$ whenever $e\in K_b$. Let $e$ be an index of the total program $x\mapsto1-f(x)$.
Then $\varphi_e(e)=1-f(e)$, so $e\in K_{1-f(e)}$, and the condition requires $f(e)=1-f(e)$.
\end{proof}

The proof is the recursive inseparability of $K_0$ and $K_1$. What fails is \emph{exact selection},
not the approximation of expected values: after $n$ simulation steps the unresolved part of each
expectation is at most $2^{-n}$, so the values are uniformly computable reals with an explicit error
bound. What cannot be computed uniformly is which action belongs to the exact argmax when
arbitrarily small advantages and genuine ties are both admitted. The obstruction disappears as soon
as one of the two is excluded.

\begin{proposition}[Positive replacements]\label{prop:approxopt}
On a class of environments with finite policy spaces and uniformly computable values, there are
promise-total agents that are
\begin{enumerate}[label=(\roman*)]
  \item $\varepsilon$-optimal everywhere, for any fixed rational $\varepsilon>0$;
  \item exactly optimal on every environment with a unique maximising policy;
  \item exactly optimal on every environment whose values are exact rationals.
\end{enumerate}
\end{proposition}

\begin{proof}
(i) Approximate every value within $\varepsilon/3$ and maximise the approximations. (ii) Refine
enclosing intervals until one policy's lower bound exceeds every other policy's upper bound, which
happens in finite time when the maximiser is unique. (iii) is Theorem~\ref{thm:promiseopt}.
\end{proof}

Theorem~\ref{thm:exactsel} is therefore a true and non-trivial version of the thesis that no
computable rule is optimal on all fair problems, on an especially fair class. Its obstruction is
classical recursion theory, not a conflict among decision theories.

\subsection{Resource bounds}

Finite clocks or circuit descriptions can support globally total optimisers, because every required
evaluation then has a finite bound. They do not make optimisation as cheap as one execution.

\begin{proposition}\label{prop:sat}
If an agent running in time polynomial in the length of the description is exactly optimal on all
environments with $n$ binary information sets and a Boolean circuit as payoff, then
$\mathrm{P}=\mathrm{NP}$.
\end{proposition}

\begin{proof}
Given a Boolean formula with $n$ variables, form the environment whose information sets are the
variables and whose payoff is $1$ exactly on satisfying assignments. Query the agent on its $n$ owned
queries and evaluate the formula on the resulting policy. By optimality the policy is satisfying
whenever a satisfying assignment exists, so this decides satisfiability in polynomial time
\autocite{cook1971}.
\end{proof}

Per-input runtime bounds do not decide program equivalence either; the construction of
Theorem~\ref{thm:haltred} uses components bounded by $|x|$ steps on input $x$. The bounded variants
of Löb's theorem proved by \textcite{critch2019} yield positive cooperation results for specified
proof-responsive programs under proof-length assumptions; they do not establish optimality on a
class of environments of the present kind.

\subsection{What a fixed semantics settles}\label{sec:fixedsemantics}

Theorem~\ref{thm:promiseopt} extends at once to any environment class that supplies a
\emph{reference-independent} effective evaluator: a computable $V(E,\pi)$ on a finite policy space,
with exactly comparable values, whether the evaluator works by source analysis, by simulation, or by
proof. Finite circuits, explicitly clocked programs with stipulated time-out behaviour, and
descriptions carrying checkable certificates of the required evaluations all have such evaluators,
and admit optimisers that are globally total on the certified language, returning a default on
invalid strings. One might conclude that once a rule for propagating the dependence of the world on
the agent's output is fixed, optimal choice is a finite computation, and the only remaining question
is which rule.

The conclusion is false for anchored kernels. A fixed, computable kernel makes each comparison
computable, but it need not admit an agent whose actual output maximises it. In
Proposition~\ref{prop:doublecall} under $O_{\mathrm{const}}$ every baseline yields a finite,
unambiguous ranking, and every baseline is ranked below its alternative. Maximising a supplied table
$V_E$ is a finite computation; being a fixed point of a baseline-dependent kernel is a separate
problem, the ratifiability problem of Part~\ref{part:single} in program form.

Two further remarks fix the scope of \enquote{largest solvable class}. For reference-independent
problems with disjoint owned queries, every decidable class on which some globally total agent is
optimal can be enlarged by a single further problem, by altering the agent on that problem's owned
queries only; under $\cls_{\mathrm{self}}$ this changes nothing on the original class, so there is
no largest such class. Conversely, the class of all terminating, rational-valued, finite-policy
templates is a maximal promise domain for that representation, so an unqualified claim that there is
no largest solvable class would also be false.

% -----------------------------------------------------------------------------------
\section{Inferring dependence from source}\label{sec:inference}

Theorem~\ref{thm:promiseopt} works because oracle substitution designates the positions at which the
agent is consulted. In an embedded environment, a closed program whose source may contain
predictors, simulations, or copies of the agent with no designated interface, the agent must
identify those positions itself. This section examines the three natural ways of doing so:
recognising the agent's code, recognising programs that compute the same function, and reasoning
about the agent's output in a formal theory. Each has a positive theorem on a restricted domain and a
characteristic failure outside it.

\subsection{Designated slots and literal copies}\label{sec:syntactic}

Suppose an environment is a template with designated agent slots, that syntactic replacement
changes exactly those slots, that every slot uses the same owned-query namespace, that all other
programs are constants, and that every finite-policy replacement terminates with an exact rational
payoff. Compiling the oracle calls of Theorem~\ref{thm:promiseopt} to these slots preserves every
counterfactual execution, so
\[
  U_{\mathrm{syn}}(E,T;A)=U_{\mathrm{oracle}}(E,A),
\]
and the uniform optimiser applies. This is a substantive positive answer: the rule evaluates
arbitrary code implementing the environment without analysing every consequence of its action; it
executes the stipulated replacements. It is also a restricted answer, since the representation
identifies the slots.

Without designated slots, literal recognition fails. Fix a globally total $T$ and let $g$ be a closed
program with $g(x)=T(x)$ for every $x$, encoded independently in a typed, atomic component so that its
representation contains no eligible occurrence of $T$'s code. (Wrapping $T$ is not enough, since the
original would remain an eligible nested occurrence.) Let the environment pay $\one[a\neq g(d(E))]$.
Literal substitution changes the live agent and leaves $g$ unchanged, so $T$ receives $0$ while an
agent taking the opposite action receives $1$. Every computation terminates. This is the fixed-copy
diagonal of Theorem~\ref{thm:fixedcopy} again, now evading a syntactic matcher. Literal matching also
violates $\cls_{\mathrm{ext}}$: two extensionally equal reference agents with different source code
produce different replacement masks. A slot-based representation avoids the problem by preserving
the agent's role independently of the code that currently occupies it; code equality and role
identity should not be conflated.

\subsection{Extensional matching}\label{sec:extensional}

The natural repair is to replace every component that computes the same function as the agent.
To make this precise, let $E$ contain a finite list of typed, atomic components $g_j$, fix comparison
domains $X_j$ and a reference function $t$, and let
\[
  S_t(E)=\{j: g_j|_{X_j}=t|_{X_j}\}.
\]
To evaluate $A$, replace exactly the components in $S_t(E)$ by $A$ and execute; write
$U_{\mathrm{ext}}(E,t;A)$ for the result. Matching occurs once, before replacement; component
boundaries and the description supplied to the agent stay fixed; and every permitted counterfactual
execution is required to terminate, which termination of the original environment does not imply.
The principal choices are \emph{local} equality, on a specified finite set or a frozen baseline
trace, and \emph{global} equality, on the entire input domain.

\begin{proposition}[Matching on the counterfactual trace]\label{prop:tracematch}
Requiring replacement exactly when a component agrees with $t$ on the inputs used in the resulting
counterfactual execution can leave no consistent replacement set, although every candidate
execution terminates.
\end{proposition}

\begin{proof}
Let $t(0)=t(1)=0$, $g(0)=0$, $g(1)=1$, and $A(0)=1$, and let the environment call $g(0)$ and then call
$g(1)$ exactly when the first result is $1$. Without replacement only the input $0$ is used, where $g$
agrees with $t$, so replacement is required. With replacement, $A(0)=1$ triggers the call on $1$, where
$g$ disagrees with $t$, so replacement is forbidden.
\end{proof}

\enquote{The inputs it is run on} must therefore name the execution that determines those inputs.
Equality of total programs on a fixed finite set is decidable by running them, so local matching
needs no appeal to Rice's theorem. It fails for a different reason.

\begin{theorem}[Local equality can leave no optimal agent]\label{thm:localmatch}
Under frozen matching on the single input $d(E)$, there is a terminating environment satisfying
$\cls_{\mathrm{self}}$ and $\cls_{\mathrm{ext}}$ on which no deterministic agent satisfies anchored
optimality~\eqref{eq:anchored}.
\end{theorem}

\begin{proof}
Let $E$ contain the components $g_0\equiv0$ and $g_1\equiv1$, call both at $d(E)$, query the agent
there, and pay $1$ exactly when the two component outputs agree. For a reference agent with
$t(d(E))=a$, exactly $g_a$ matches. With $t$ itself substituted, the outputs are $a$ and $1-a$, so
$U_{\mathrm{ext}}(E,t;t)=0$. With an agent $A$ taking $1-a$, the matched component $g_a$ becomes $A$
and both outputs equal $1-a$, so $U_{\mathrm{ext}}(E,t;A)=1$.
\end{proof}

This is not an undecidability theorem. It is a pathology of identifying computations through
accidental agreement on one input, and it should not be presented as a theorem about functional
decision theory without that identification.

\begin{proposition}[Rematching does not compose]\label{prop:rematch}
Let $R_{t\to a}$ replace by $a$ the components that match $t$. With components $g_0\equiv0$,
$g_1\equiv1$ and reference $t\equiv0$, and with matches recomputed after each replacement,
$R_{a\to b}(R_{t\to a}(E))\neq R_{t\to b}(E)$ for $a\equiv1$ and $b\equiv0$.
\end{proposition}

\begin{proof}
$R_{t\to a}$ turns $g_0$ into $1$, giving the pair $(1,1)$. Both components now match $a$, and $R_{a\to b}$
gives $(0,0)$. Direct replacement of $t$ by $b=t$ leaves $(0,1)$.
\end{proof}

Keeping an original dependency label attached to the first component restores composition, but the
label is additional structure; equality of the current functions does not recover it. Global
matching avoids accidental agreement, but it identifies the wrong thing.

\begin{theorem}[Global equality leaves the diagonal intact]\label{thm:complement}
For globally total agents, global-equivalence replacement of atomic components admits no universal
optimiser when fixed algorithms are allowed.
\end{theorem}

\begin{proof}
Fix $T$, compile an atomic program $h_T(x)=1-T(x)$, and let the environment pay
$\one[a=h_T(d(E))]$. Since $h_T$ disagrees with $T$ everywhere, it is not replaced. $T$ receives $0$,
an agent taking the opposite action receives $1$, and all calls terminate.
\end{proof}

The diagnosis is that \enquote{depends functionally on $T$} is not \enquote{is extensionally equal
to $T$}. If the intended counterfactual turns $1-T$ into $1-A$, it transports dependence through a
program context, which is a stronger operation than equality-based replacement. Global matching
also brings in undecidability, and the expected route needs one correction. For a uniformly total
family $g_e$ and a fixed total $T$, non-equivalence has a finite witness, so $\{e:g_e=T\}$ is co-c.e.;
one cannot uniformly make equality hold exactly when a machine halts, but one can make it hold exactly
when it does not.

\begin{theorem}[Global matching can decide halting]\label{thm:haltred}
Assume atomic global matching, frozen descriptions, and globally total binary agents, and let $P$ be a
decidable class of descriptions that contains all the environments constructed below and omits
arbitrarily long strings. No agent satisfies anchored optimality on all environments in $P$.
\end{theorem}

\begin{proof}
Fix a candidate $T$ and let
\[
  g_e(x)=
  \begin{cases}
    1-T(x), & x\notin P\text{ and machine }e\text{ halts within }|x|\text{ steps},\\
    T(x),   & \text{otherwise}.
  \end{cases}
\]
Each $g_e$ is total, and $g_e=T$ exactly when $e$ does not halt. Let $E_e$, with $d_e=d(E_e)\in P$,
receive the action $a$, compute $b=g_e(d_e)$, and pay $u(a,b)=2(1-b)+a$:
\begin{center}
\small
\begin{tabular}{@{}lcc@{}}
  \toprule
  & action $0$ & action $1$ \\
  \midrule
  $b=0$ & 2 & 3 \\
  $b=1$ & 0 & 1 \\
  \bottomrule
\end{tabular}
\end{center}
Because $d_e\in P$, $g_e(d_e)=T(d_e)$ whether or not $e$ halts. If $e$ does not halt, global matching
replaces $g_e$ by the candidate $A$, so $b=A(d_e)$ and the values of the two actions are $2$ and $1$;
optimality requires $T(d_e)=0$. If $e$ halts, $g_e$ stays fixed, and with $b$ held fixed action $1$
pays one unit more; optimality requires $T(d_e)=1$. The computable map $e\mapsto T(d_e)$ would decide
the halting problem.
\end{proof}

The payoff table is Newcomb's, and the reduction works by making it undecidable whether its
predictor counts as linked to the agent. Every execution terminates and every payoff is rational,
but the result uses global equality, atomic components compiled from $T$, and a class already
admitting fixed-constant diagonals; it satisfies executed-query locality and violates kernel
locality. It is a valid reduction, not a semantics-independent impossibility arising from the mere
presence of embedded programs. Under local matching $g_e$ always matches at $d_e$, so action $0$ always
wins and the reduction disappears. Two further qualifications bound any positive subclass. Clocks do
not decide global equivalence, since even a program that simulates machine $e$ for only $|x|$ steps
can equal zero everywhere exactly when $e$ never halts. And \enquote{replace when a theory $S$ proves
equality} changes the semantics: it agrees with global matching only where certified matching
information is complete, since positive certificates do not show that uncertified components are
unequal.

\subsection{Proofs and logical counterfactuals}\label{sec:logical}

Three questions must be separated: whether proofs can establish the values of independently
specified counterfactual programs; whether implications about the actual agent can serve as
counterfactual values; and whether logical uncertainty supplies useful finite-time estimates. The
answers are yes, no, and not by itself.

\begin{proposition}[Ground proof-based optimisation]\label{prop:proofopt}
Let $S$ be a sound, effectively axiomatised arithmetic theory that verifies finite program traces.
The agent that, on an environment in $\mathcal P_{\mathbb Q}$, enumerates $S$-proofs until it has, for
every $\pi\in\Pol_E$, a proof of $\mathrm{Eval}(E,\pi)=r_\pi$ for some rational $r_\pi$, and then
maximises the $r_\pi$, is promise-total and optimal under $O_{\mathrm{oracle}}$.
\end{proposition}

\begin{proof}
Each terminating execution $E^\pi$ has a finite trace, the trace has a proof in $S$, and soundness
makes every proved value correct. Only finitely many proofs are needed.
\end{proof}

The incompleteness of a fixed theory does not prevent it from verifying every particular terminating
computation; what it may fail to prove is a uniform termination statement or a global equivalence.
The proposition evaluates the same counterfactual programs as Theorem~\ref{thm:promiseopt} and infers
no new semantics from proof search.

Material implication is different. Write $H_a$ for $T(d(E))=a$ and $Y_u$ for $U(E,T)=u$. If
$S\vdash\neg H_a$, then $S\vdash H_a\to Y_u$ for every $u$, so the relation
\[
  \mathcal V(E,T,a)=\{u: (H_a\to Y_u)\text{ is true}\}
\]
contains only the actual payoff at the actual action and every rational at every other action. It is
a relation, not an operator; any numerical counterfactual drawn from it needs a selection principle
that the implication does not supply. For an actually terminating computable agent, a sufficiently
strong theory verifies the agent's execution and so refutes every non-actual antecedent, and
conditioning a coherent probability over its models on such an antecedent is undefined. Proof search
can also select a spurious value, as in the \enquote{five-and-ten} problem
\autocite[§2.1]{demskigarrabrant2019}, which can be realised by a total program.

\begin{theorem}[A computable five-and-ten failure]\label{thm:fiveten}
Let $E_0$ pay the action chosen from $\{5,10\}$, and let $S$ be an effectively axiomatised theory
satisfying the derivability conditions for Löb's theorem \autocite{lob1955}. There is a globally
total agent $T$ that, on $d(E_0)$, takes $5$ because it has found a proof that taking $10$ yields $0$.
It receives $5$, while the constant agent $10$ receives $10$.
\end{theorem}

\begin{proof}
By the recursion theorem let $T$ return $10$ on descriptions other than $d(E_0)$, and on $d(E_0)$
enumerate $S$-proofs until it finds a proof of
\[
  \psi=\bigl(T(d(E_0))=5\to U(E_0,T)=5\bigr)\wedge\bigl(T(d(E_0))=10\to U(E_0,T)=0\bigr),
\]
and then return $5$. Let $q$ be the sentence $T(d(E_0))=5$. The definitions of $T$ and $E_0$ give
$S\vdash q\to\psi$, since under $q$ the first conjunct holds and the second is vacuous. Hence
$S\vdash\Prov_Sq\to\Prov_S\psi$. Formalising the search gives $S\vdash\Prov_S\psi\to q$. So
$S\vdash\Prov_Sq\to q$, and Löb's theorem yields $S\vdash q$, whence $S\vdash\psi$. A proof of $\psi$
therefore exists, the search halts, and $T$ returns $5$.
\end{proof}

This refutes the displayed implication-based selection scheme, not every logical-counterfactual
procedure. \textcite[§2]{bensontilsen2014} analyses the mechanism, and his classification of
outcomes by proof-search order concerns a specified transparent-universe model with a provability
oracle; it should not be quoted as a theorem about every computable proof-based agent. A different
response, the \enquote{chicken rule}, has the agent take any action it can prove it will not take
\autocite{slepnev2011}; within Benson-Tilsen's model it is the only universally successful search
strategy \autocite[Corollary~4.6]{bensontilsen2014}, and it blocks spurious proofs at the cost of
counterfactuals that can differ from what another agent would actually obtain
\autocite{fallenstein2015chicken}.

A precise logical semantics that avoids spurious values does so by becoming partial.

\begin{definition}[Guarded modal counterfactuals]\label{def:guarded}
Let $(A_i)$ and $(U_j)$ be closed modal formulas, provably mutually exclusive and exhaustive, describing
the actions and outcomes of an agent and a problem. Action $i$ \emph{leads to} outcome $u_j$ if
\[
  \mathrm{GL}\vdash A_i\to U_j\qquad\text{and}\qquad\mathrm{GL}\nvdash\neg A_i,
\]
and has no value otherwise \autocite{fallenstein2015third}.
\end{definition}

The outcome is unique when defined and is undefined if the action is provably impossible or no
outcome is determined; the semantics \enquote{fails visibly} rather than assigning a spurious value.
Fallenstein uses it to state the optimality of modal updateless decision theory on problems where
every action has a value \autocite{fallenstein2015third}. This is optimality relative to the specified
counterfactuals, not relative to physical agent replacement, and the two can diverge: on a variant of
five-and-ten where every action has a value, these counterfactuals certify as optimal an action that
another agent would beat \autocite{fallenstein2015chicken}. In provability logic $\mathrm{GL}$
\autocite{boolos1993} the guard is informative. For ordinary programs it is not.

\begin{proposition}[Guarded counterfactuals for terminating programs]\label{prop:guardedarith}
Replace $\mathrm{GL}$ by an arithmetic theory $S$ that verifies terminating computations and proves the
determinism of its interpreter. For an agent that terminates on $d$, every non-actual action has no
value.
\end{proposition}

\begin{proof}
If the actual action is $b$, the finite execution trace gives $S\vdash T(d)=b$, and determinism gives
$S\vdash T(d)\neq a$ for every $a\ne b$, so every non-actual action fails the guard.
\end{proof}

This reconstruction is not attributed to the source of Definition~\ref{def:guarded}; it explains
why positive results for modal agents, or agents with provability oracles, do not transfer directly
to ordinary terminating programs.

The modal setting has its own diagonal and its own optimality theorem, and their quantifiers matter.
Consider fixed, fully modalised decision schemata $T(\mathbf u)$ and problems that are \emph{provably
extensional}, in the sense that $\mathrm{GL}\vdash(\mathbf a\leftrightarrow\mathbf b)\to
(P(\mathbf a)\leftrightarrow P(\mathbf b))$. Every such fixed theory with at least two actions and
two outcomes fails on some provably extensional problem \autocite{fallenstein2014evil}: the problem
rewards any action other than the one the theory takes when facing the constant second-best outcome,
modal fixed-point uniqueness makes that action the theory's own, and a different constant action
wins. This is the fixed-constant diagonal in modal form. Conversely, for every finite provably
extensional problem $P$ there is a true sentence $\varphi$ such that modal updateless decision theory
using $\mathrm{PA}+\varphi$ selects an optimal action on $P$
\autocite{fallenstein2014optimality,fallenstein2014improving}. The quantifiers are
$\forall P\,\exists\varphi$: there is no single proof system that works for all problems, and the
general procedure of \textcite{soaresfallenstein2015} that queries proofs uses a halting oracle. $\mathrm{GL}$
itself is decidable, so none of this shows that restricted modal computation is uncomputable.

Logical induction \autocite{garrabrant2016li} supplies computable, improving assessments of logical
sentences. It does not specify which counterfactual program or intervention to evaluate, and it
imposes no exact constraint on any finite stage: by its Theorem~4.6.1, changing finitely many market
days preserves the logical induction criterion. A decision procedure can use an inductor to estimate
the values of policy-substituted programs, but Theorem~\ref{thm:exactsel} still applies to any
resulting computable rule. Decision procedures built on logical induction typically add
\emph{exploration}, taking an action whenever its probability falls below a threshold so that the
conditionals on it remain defined \autocite{garrabrant2017obstacles}; Section~\ref{sec:candidates}
shows that what such a procedure computes depends on the exploration scheme it is given.

\subsection{What inference supplies}

The pattern across these results is uniform. Wherever a substitution is designated---by an oracle
interface, by slots, or by an evaluator whose results a proof certifies---optimisation is
computable on the corresponding termination promise. Wherever the agent must decide for itself
which parts of a closed program would change with its output, each recognition rule fails in its
own way: literal matching misses encoded copies, local matching manufactures dependence from
accidental agreement, global matching misses transformed dependence and can be undecidable, and
implication-based reasoning about the actual agent yields either spurious values or none. None of
these failures is a theorem that no coherent rule exists. They show that the rule is the open
problem, which the next section takes up axiomatically.

% -----------------------------------------------------------------------------------
\section{Axioms for counterfactual operators}\label{sec:operators}

Section~\ref{sec:inference} leaves one question: is there a coherent way to assign outcomes to the
non-actual outputs of a computation? This is the problem of \emph{logical counterfactuals}. This
section states candidate desiderata precisely and determines every consistent and inconsistent
combination of them. There is a sharp impossibility among three individually satisfiable axioms,
and it identifies exactly which notion of program equivalence is responsible. Under a different
notion, all the desiderata hold together.

\subsection{Preliminary repairs}

Three features of a naive formulation produce failures before the substantive question is reached.
They are recorded here and excluded below by stipulation.

\emph{Transported contexts need not terminate.} Let a counterfactual operator be required to
\emph{transport} dependence: a component $C[T]$ that computes a context $C$ applied to the agent is
evaluated, under the alternative, as $C$ applied to the altered agent. It is natural to promise only
that the host environment terminates when its components are replaced by tables. That promise is too
weak. Let $T\equiv0$, let $C[P](x)$ return $0$ if $P(x)=0$ and diverge if $P(x)=1$, treat $g=C[T]$ as one
opaque component, and let the host return $g(d)$. The host terminates under every whole-component table
replacement, and $g$ is the total constant $0$, yet the transported execution $C[T^{d\leftarrow1}](d)$
diverges. On this domain totality and execution-faithful transport are already inconsistent. The
domain is therefore required to be \emph{intervention-total}: every supplied computational context
terminates under every permitted substitution, as can be ensured by finite tables, a terminating
language, certificates, or an explicit promise. Atomic boundaries matter, since replacing an opaque
outer component by a table can hide divergence inside it.

\emph{Source-changing replacement changes inputs.}

\begin{proposition}\label{prop:rawsource}
If the description supplied to the agent is the literal source and replacing a component changes it,
baseline correctness is inconsistent with invariance under replacement of a component by an
extensionally equal one.
\end{proposition}

\begin{proof}
Let $E$ ignore the agent's action and return $g(d(E))$ with $g(x)=|x|\bmod2$. Replacing $g$ by an
extensionally identical implementation padded by one character changes the parity of the source,
and hence the actual payoff, while a constant agent takes the same action in both versions.
\end{proof}

The components are equal but are called on different inputs. Below, comparisons of implementations
preserve the distinguished inputs, and counterfactual transformations retain the original
observation.

\emph{Single-input transport changes one input.} An intervention at $d$ replaces $T$ by
\[
  T^{d\leftarrow a}(x)=\begin{cases}a,&x=d,\\T(x),&x\ne d,\end{cases}
\]
and strong transport evaluates $C[T^{d\leftarrow a}]$. A context $C[T]=T\circ h$ does not in general
become the constant $a$: if $h(d)\neq d$, then $C[T^{d\leftarrow a}](d)=T(h(d))$. Replacing the whole agent
by a constant is a different, policy-level intervention.

\subsection{The finite context domain}

The main theorems use a finite domain in which termination and program equivalence raise no
difficulty, so that what remains is the semantic question alone.

\begin{definition}[Context environment]\label{def:context}
A \emph{context environment} is $E=(u;c_1,\dots,c_k)$, where the action set is $A=\{0,1\}$, each
component $c_i\colon A\to B_i$ is a table with values in a finite set $B_i$, and the host
$u\colon A\times\prod_iB_i\to\mathbb Q$ is a finite payoff table. The argument of $c_i$ is the
\emph{live parameter}: the agent's output at the distinguished input, as seen by component $i$. $E$ is
\emph{independent} if every $c_i$ is constant. Write
\[
  V_E(a)=u\bigl(a,c_1(a),\dots,c_k(a)\bigr),\qquad
  W_{E,b}(a)=u\bigl(a,c_1(b),\dots,c_k(b)\bigr).
\]
\end{definition}

$V_E$ is coherent substitution through every context; $W_{E,b}$ changes the action at the port and
freezes every component at its value under the baseline $b$. The definition of independence is not
circular: it is invariance under every substitution of the declared live parameter, not the absence
of an undefined relation of \enquote{depending on $T$}. The domain is the causal-matrix domain
with its factorisation made explicit.

\begin{proposition}\label{prop:contextmatrix}
Let $M_E[b,a]=W_{E,b}(a)$. Then $M_E$ is the causal matrix of the one-shot problem in which an
infallible predictor forecasts the agent's output $b$ and the components respond to the forecast, and
$M_E[a,a]=V_E(a)$. Every real $2\times2$ matrix with rational entries is $M_E$ for some context
environment, and $E$ is independent exactly when all rows of $M_E$ coincide.
\end{proposition}

\begin{proof}
With structural states $b\in A$, $\mu_a=\delta_a$, and $U(b,a)=u(a,c(b))$, the causal matrix is
$\sum_s\mu_b(s)U(s,a)=W_{E,b}(a)$. Given $M$, take one identity component $c(z)=z$ and $u(a,x)=M[x,a]$. If
every $c_i$ is constant, $W_{E,b}$ does not depend on $b$; if some row differs from another, some
component's value differs between the two baselines.
\end{proof}

Strictly, the converse in the last clause holds for the matrix, not for the environment: distinct
component values may be invisible to the host. Nothing below uses the converse.

\begin{definition}[Counterfactual kernel]\label{def:kernel}
A \emph{kernel} assigns to every context environment $E$, baseline $b\in A$, and action $a\in A$ a value
$\kappa(E,b,a)\in\mathbb Q\cup\{\bot\}$, where $\bot$ means undefined. The baseline is the agent's actual
output, $b=T(d(E))$. A deterministic agent $\tau$ on context environments is
\emph{$\kappa$-ratifiable} at $E$ if $\kappa(E,\tau(E),\cdot)$ is defined at every action and is
maximised at $\tau(E)$.
\end{definition}

That the kernel depends on $T$ only through the local baseline is an explicit stipulation;
the program extension in Section~\ref{sec:programmap} keeps the same interface. Agent extensionality
is then automatic, because agents computing the same function have the same baseline. Two
notions of extensionality for the environment must be distinguished.

\begin{definition}[Grounded and parametric equivalence]
$E\sim_bE'$ (\emph{grounded equivalence at $b$}) if $E$ and $E'$ have the same host and interfaces and
$c_i(b)=c'_i(b)$ for every $i$. $E\approx E'$ (\emph{parametric equivalence}) if they have the same
host and interfaces and $c_i=c'_i$ as functions on $A$.
\end{definition}

Grounded equivalence has a direct reading for programs. Once the actual agent is inserted, a
component is a closed computation returning $c_i(b)$, and two closed computations can be
extensionally equal although their open contexts would respond differently to another action. The
function $x\mapsto T(x)$ and the constant function with value $T(d)$ agree at $d$, and after
insertion of $T$ both are closed programs with output $T(d)$. Parametric equivalence identifies
contexts only when they agree under every substitution, and keeps the information that grounded
equivalence discards.

\begin{definition}[Desiderata]\label{def:desiderata}
A kernel may satisfy the following conditions, for all context environments $E$, $E'$ and all
$a,b\in A$.
\begin{description}[font=\normalfont\sffamily,leftmargin=3.2em,labelwidth=2.4em,labelsep=.8em,itemsep=.35em]
  \item[CF1] \emph{Baseline.} $\kappa(E,b,b)=V_E(b)$.
  \item[CF2] \emph{Totality.} $\kappa(E,b,a)\ne\bot$.
  \item[CF3] \emph{Independent calibration.} If $E$ is independent, $\kappa(E,b,a)=V_E(a)$.
  \item[CF4] \emph{Transport.} If $\kappa(E,b,a)$ is defined, it equals $V_E(a)$.
  \item[CF5] \emph{Grounded extensionality.} If $E\sim_bE'$, then $\kappa(E,b,a)=\kappa(E',b,a)$, including
        equality of definedness.
  \item[CF6] \emph{Conservativity.} Every defined value has the form $u(a,z_1,\dots,z_k)$ with $z_i=c_i(b)$
        for every constant context $c_i$.
  \item[CF7] \emph{Composition.} A transformed environment is the record $(E,b;a)$, whose evaluator
        returns $\kappa(E,b,a)$, and a further intervention overwrites the target:
        $(E,b;a)[a\leftarrow z]=(E,b;z)$.
  \item[CF8] \emph{Computability.} $\kappa$, with $\bot$ returned explicitly, is computable.
  \item[CF9] \emph{Realisability.} Every $\kappa(E,b,a)$ is defined and is the payoff that some agent taking
        $a$ obtains under coherent substitution of the live parameter.
  \item[CF10] \emph{Ratifiability.} Some computable $\tau$ is $\kappa$-ratifiable at every $E$.
\end{description}
\end{definition}

Three remarks fix the reading. \ax{CF4} transports the whole supplied context and preserves the
host, and it is conditional on definedness; a version requiring the transported values to exist
would incorporate \ax{CF2}. Every agent taking $a$ obtains exactly $V_E(a)$ under slot substitution,
and constant agents supply the alternatives, so
\[
  \ax{CF9}\iff\ax{CF2}+\ax{CF4}.
\]
And \ax{CF7} is deliberately weak: every kernel has such a lift, since the record exists even when its
evaluator returns $\bot$. A history-free requirement, or one demanding fresh rematching after each
intervention, would be a different axiom (Section~\ref{sec:candidates}). \ax{CF6} constrains only the
constant contexts; how the others respond is left open. On the finite domain \ax{CF8} holds for every
kernel defined by a finite recipe from the tables, and \ax{CF7} for every kernel; both are recorded
because they become substantive on program domains.

\subsection{Four kernels}

\begin{definition}[Witness kernels]
\[
  \kappa_{\mathrm{slot}}(E,b,a)=V_E(a),\qquad
  \kappa_{\mathrm{port}}(E,b,a)=W_{E,b}(a),
\]
\[
  \kappa_{\mathrm{base}}(E,b,a)=\begin{cases}V_E(b),&a=b,\\\bot,&a\ne b,\end{cases}\qquad
  \kappa_{\mathrm{flat}}(E,b,a)=V_E(b).
\]
\end{definition}

By Proposition~\ref{prop:contextmatrix}, $\kappa_{\mathrm{slot}}$ reads the diagonal of the causal
matrix and $\kappa_{\mathrm{port}}$ reads row $b$: they are ex-ante and causal evaluation, now
appearing as counterfactual semantics. $\kappa_{\mathrm{base}}$ values only the actual output, as
partial semantics do; $\kappa_{\mathrm{flat}}$ assigns the actual payoff to every alternative, so
every agent is ratifiable under it. All four satisfy \ax{CF1}, \ax{CF7}, and \ax{CF8}. Their other
properties are as follows.

\begin{lemma}\label{lem:witnesses}
$\kappa_{\mathrm{slot}}$ satisfies every desideratum except \ax{CF5}, and the agent choosing the first
maximiser of $V_E$ ratifies it. $\kappa_{\mathrm{port}}$ satisfies \ax{CF2}, \ax{CF3}, \ax{CF5}, and
\ax{CF6}. $\kappa_{\mathrm{base}}$ satisfies \ax{CF4}, \ax{CF5}, and \ax{CF6}. $\kappa_{\mathrm{flat}}$
satisfies \ax{CF2}, \ax{CF5}, and \ax{CF10}.
\end{lemma}

\begin{proof}
$\kappa_{\mathrm{slot}}$ is total and equals $V_E$, which gives \ax{CF2}--\ax{CF4} and \ax{CF9}; its
values $u(a,c(a))$ have $c_i(a)=c_i(b)$ for constant contexts, giving \ax{CF6}; the first maximiser of
$V_E$ is computable from the tables. $\kappa_{\mathrm{port}}$, $\kappa_{\mathrm{base}}$, and
$\kappa_{\mathrm{flat}}$ depend on $E$ only through $u$ and the baseline outputs $c_i(b)$, which grounded
equivalence preserves, so all three satisfy \ax{CF5}. $\kappa_{\mathrm{port}}$ is total; on an
independent environment $W_{E,b}=V_E$; and its values have the form required by \ax{CF6}.
$\kappa_{\mathrm{base}}$ is defined only at $a=b$, where its value $V_E(b)=u(b,c(b))$ satisfies \ax{CF4}
and \ax{CF6}. $\kappa_{\mathrm{flat}}$ is total and constant in $a$, so every agent ratifies it.
\end{proof}

Nothing more holds: $\kappa_{\mathrm{port}}$ violates \ax{CF4}, \ax{CF9}, and \ax{CF10} on the
environments of Theorems~\ref{thm:transport} and~\ref{thm:calibration} below; $\kappa_{\mathrm{base}}$ is
not total and so violates \ax{CF2}, \ax{CF3}, \ax{CF9}, and \ax{CF10}; and $\kappa_{\mathrm{flat}}$
violates \ax{CF3}, \ax{CF4}, \ax{CF6}, and \ax{CF9} on any independent environment whose two
actions have different payoffs.

\subsection{The projection lemma and two theorems}

For a context environment $E$ and a baseline $b$, let $E^b$ replace every component by the constant
with its baseline value, $c^b_i(a)=c_i(b)$. Then $E\sim_bE^b$, $E^b$ is independent, and
$V_{E^b}(a)=W_{E,b}(a)$.

\begin{lemma}[Projection]\label{lem:projection}
\ax{CF3} and \ax{CF5} together imply $\kappa=\kappa_{\mathrm{port}}$. \ax{CF5} and \ax{CF6} together
imply $\kappa(E,b,a)=W_{E,b}(a)$ whenever $\kappa(E,b,a)$ is defined.
\end{lemma}

\begin{proof}
\ax{CF3} defines the values at $E^b$ as $V_{E^b}(a)=W_{E,b}(a)$, and \ax{CF5} transfers them, with
their definedness, to $E$. For the second claim, transfer a defined value to $E^b$ by \ax{CF5}. Every
component of $E^b$ is constant, so \ax{CF6} fixes every component value at $c_i(b)$, and the value is
$u(a,c(b))=W_{E,b}(a)$.
\end{proof}

This is the main structural fact: \emph{calibration together with the erasure of dependency
information forces causal evaluation}. The two semantic requirements reproduce, without any
appeal to decision-theoretic intuition, the frozen-environment counterfactual that \cri{3} uses.

\begin{theorem}[Transport and grounded equivalence]\label{thm:transport}
$\{\ax{CF2},\ax{CF4},\ax{CF5}\}$ is minimally inconsistent.
\end{theorem}

\begin{proof}
Let $u(a,x)=x$, and compare $c(z)=z$ with $c'(z)=0$ at the baseline $0$. The instantiated component
outputs agree, so $E\sim_0E'$, but $V_E(1)=1$ and $V_{E'}(1)=0$. \ax{CF2} defines both values,
\ax{CF4} makes them unequal, and \ax{CF5} makes them equal. Omitting \ax{CF2}, \ax{CF4}, or
\ax{CF5}, the remaining pair is satisfied by $\kappa_{\mathrm{base}}$, $\kappa_{\mathrm{port}}$, or
$\kappa_{\mathrm{slot}}$ respectively.
\end{proof}

No computability or undecidability enters; the obstruction is the loss of information under an
equivalence relation. The second theorem connects that loss with the instability of
Part~\ref{part:single}.

\begin{theorem}[Calibration and ratifiability]\label{thm:calibration}
$\{\ax{CF3},\ax{CF5},\ax{CF10}\}$ is minimally inconsistent, and each of its proper subsets is satisfied
by a total, computable kernel satisfying \ax{CF1}, \ax{CF2}, \ax{CF7}, and \ax{CF8}.
\end{theorem}

\begin{proof}
Take the copied-action environment $u(a,x)=\one[a\neq x]$, $c(z)=z$. For either baseline,
$W_{E,b}(b)=0$ and $W_{E,b}(1-b)=1$. By Lemma~\ref{lem:projection} \ax{CF3} and \ax{CF5} force this table,
so every actual action recommends its opposite, contradicting \ax{CF10}. Omitting \ax{CF3}, \ax{CF5},
or \ax{CF10}, the remaining pair is satisfied by $\kappa_{\mathrm{flat}}$, $\kappa_{\mathrm{slot}}$, or
$\kappa_{\mathrm{port}}$, all total and computable.
\end{proof}

The witness environment is Death in Damascus with an infallible predictor
(Proposition~\ref{prop:did}). That no pure action is ratifiable there is familiar; what the theorem
adds is a derivation of the causal kernel from two semantic requirements that make no reference to
causation, and the exact minimality of the triple. Unlike the within-problem maps of
Section~\ref{sec:commit}, it is a minimal inconsistent family of three individually satisfiable
conditions.

\subsection{The complete map}

\begin{theorem}[Consistency map]\label{thm:cfmap}
On the finite context domain, the inclusion-minimal inconsistent subsets of
$\{\ax{CF1},\dots,\ax{CF10}\}$ are exactly
\begin{gather*}
  \{\ax{CF5},\ax{CF9}\},\quad
  \{\ax{CF2},\ax{CF4},\ax{CF5}\},\quad
  \{\ax{CF3},\ax{CF4},\ax{CF5}\},\\
  \{\ax{CF3},\ax{CF5},\ax{CF10}\},\quad
  \{\ax{CF4},\ax{CF5},\ax{CF10}\},\quad
  \{\ax{CF5},\ax{CF6},\ax{CF10}\},
\end{gather*}
and the maximal consistent subsets are exactly
\begin{center}
\small
\begin{tabular}{@{}ll@{}}
  \toprule
  witness & maximal consistent subset \\
  \midrule
  $\kappa_{\mathrm{slot}}$ & all except \ax{CF5} \\
  $\kappa_{\mathrm{port}}$ & $\{\ax{CF1},\ax{CF2},\ax{CF3},\ax{CF5},\ax{CF6},\ax{CF7},\ax{CF8}\}$ \\
  $\kappa_{\mathrm{base}}$ & $\{\ax{CF1},\ax{CF4},\ax{CF5},\ax{CF6},\ax{CF7},\ax{CF8}\}$ \\
  $\kappa_{\mathrm{flat}}$ & $\{\ax{CF1},\ax{CF2},\ax{CF5},\ax{CF7},\ax{CF8},\ax{CF10}\}$ \\
  \bottomrule
\end{tabular}
\end{center}
\end{theorem}

\begin{proof}
\emph{Inconsistency.} $\{\ax{CF5},\ax{CF9}\}$: \ax{CF9} gives \ax{CF2} and \ax{CF4}, and
Theorem~\ref{thm:transport} applies. $\{\ax{CF3},\ax{CF4},\ax{CF5}\}$: in the environment of
Theorem~\ref{thm:transport}, \ax{CF3} sets the value of action $1$ in the constant-context environment
$E'$ to $0$, \ax{CF5} transfers it to $E$, and \ax{CF4} requires $1$. $\{\ax{CF4},\ax{CF5},\ax{CF10}\}$:
take $u(a,x)=x$, $c(z)=z$, and let \ax{CF10} select $b$. The alternative $1-b$ is defined; \ax{CF5}
transfers definedness to $E^b$, whose component is the constant $b$, and \ax{CF4} requires the values
$1-b$ at $E$ and $b$ at $E^b$, which \ax{CF5} makes equal. $\{\ax{CF5},\ax{CF6},\ax{CF10}\}$: by
Lemma~\ref{lem:projection}, defined values are those of $\kappa_{\mathrm{port}}$, and \ax{CF10} defines
both values at the selected baseline in the environment of Theorem~\ref{thm:calibration}. The other two
sets are Theorems~\ref{thm:transport} and~\ref{thm:calibration}.

\emph{Minimality.} Deleting the members of $\{\ax{CF3},\ax{CF4},\ax{CF5}\}$ in turn leaves pairs
satisfied by $\kappa_{\mathrm{base}}$, $\kappa_{\mathrm{port}}$, and $\kappa_{\mathrm{slot}}$; for
$\{\ax{CF4},\ax{CF5},\ax{CF10}\}$ by $\kappa_{\mathrm{flat}}$, $\kappa_{\mathrm{slot}}$, and
$\kappa_{\mathrm{base}}$; for $\{\ax{CF5},\ax{CF6},\ax{CF10}\}$ by $\kappa_{\mathrm{slot}}$,
$\kappa_{\mathrm{flat}}$, and $\kappa_{\mathrm{port}}$. \ax{CF5} and \ax{CF9} are satisfied singly by
$\kappa_{\mathrm{port}}$ and $\kappa_{\mathrm{slot}}$. The remaining two cases are in
Theorems~\ref{thm:transport} and~\ref{thm:calibration}.

\emph{Exhaustiveness.} Let $\Sigma$ contain none of the six sets. If $\ax{CF5}\notin\Sigma$,
$\kappa_{\mathrm{slot}}$ witnesses $\Sigma$. Otherwise $\ax{CF9}\notin\Sigma$. If $\ax{CF3}\in\Sigma$, then
\ax{CF4} and \ax{CF10} are excluded and $\kappa_{\mathrm{port}}$ witnesses $\Sigma$. If $\ax{CF4}\in\Sigma$ and
$\ax{CF3}\notin\Sigma$, then \ax{CF2} and \ax{CF10} are excluded and $\kappa_{\mathrm{base}}$ witnesses
$\Sigma$. If neither \ax{CF3} nor \ax{CF4} is in $\Sigma$ but \ax{CF10} is, \ax{CF6} is excluded and
$\kappa_{\mathrm{flat}}$ witnesses $\Sigma$; if \ax{CF10} is also absent, $\kappa_{\mathrm{port}}$ does.
Each displayed set is maximal, since adding any further desideratum completes one of the six cores.
\end{proof}

Every core contains \ax{CF5}, and deleting \ax{CF5} leaves a single kernel satisfying everything else.

\begin{corollary}[A coherent semantics]\label{cor:parametric}
If \ax{CF5} is replaced by \emph{parametric extensionality}, $E\approx E'\Rightarrow\kappa(E,b,a)=\kappa(E',b,a)$,
then $\kappa_{\mathrm{slot}}$ satisfies all ten desiderata, and the agent choosing the first maximiser
of $V_E$ ratifies it.
\end{corollary}

\begin{proof}
$V_E$ depends on $E$ only through $u$ and the functions $c_i$, which parametric equivalence preserves;
the rest is Lemma~\ref{lem:witnesses}.
\end{proof}

The impossibility therefore concerns the erasure of dependency information, not extensionality in
every sense. The same positive model arises from a supplied dependency graph: a finite acyclic graph
of typed variables with total structural equations, a designated node for the decision computation,
and fixed exogenous data. Intervening on that node and recomputing its descendants, in the manner of
structural causal models \autocite{pearl2009}, retains the independent mechanisms and yields the
finite table $V_E$. The graphical formulation of functional decision theory intervenes in exactly
this way on a node representing the agent's computation \autocite[§§5--6]{yudkowskysoares2017}; what
the formulation does not provide, and what no axiom here supplies, is a procedure that extracts the
graph from an arbitrary closed program.

In matrix terms the map says the following. Grounded extensionality allows the kernel to see only
the host and the baseline outputs of the components, that is, the row of the causal matrix belonging
to the actual output, together with the payoff table. Calibration then pins the kernel to that row
(Lemma~\ref{lem:projection}), and Theorem~\ref{thm:calibration} is the observation that the diagonal
entry need not be maximal in its row. Parametric extensionality allows the kernel to see the whole
matrix, and then the diagonal is available. The choice between causal and ex-ante evaluation of
Part~\ref{part:single}, which there was a choice between two laws, reappears here as a choice between
two notions of when two computations are the same.

\subsection{Other conventions}\label{sec:otherconventions}

The map depends on how undefined values and realisability are treated, and these conventions are not
cosmetic. Suppose realisability is weakened to \ax{CF9$'$}: \emph{whenever} $\kappa(E,b,a)$ is defined,
it is attainable under slot substitution, so that \ax{CF9$'$} is equivalent to \ax{CF4}. Suppose
ratifiability is weakened to \ax{CF10$'$}: some computable $\tau$ has $\kappa(E,\tau(E),\tau(E))$ defined
and maximal among the defined values.

\begin{proposition}\label{prop:cfmapconditional}
With \ax{CF9$'$} and \ax{CF10$'$} in place of \ax{CF9} and \ax{CF10}, the minimal inconsistent subsets
are exactly
\begin{gather*}
  \{\ax{CF2},\ax{CF4},\ax{CF5}\},\quad \{\ax{CF2},\ax{CF5},\ax{CF9$'$}\},\quad
  \{\ax{CF3},\ax{CF4},\ax{CF5}\},\\ \{\ax{CF3},\ax{CF5},\ax{CF9$'$}\},\quad
  \{\ax{CF3},\ax{CF5},\ax{CF10$'$}\},\quad \{\ax{CF2},\ax{CF5},\ax{CF6},\ax{CF10$'$}\},
\end{gather*}
and the maximal consistent subsets are those of Theorem~\ref{thm:cfmap} with the primed axioms
substituted, except that $\kappa_{\mathrm{base}}$ now witnesses
\[
  \{\ax{CF1},\ax{CF4},\ax{CF5},\ax{CF6},\ax{CF7},\ax{CF8},\ax{CF9$'$},\ax{CF10$'$}\}.
\]
\end{proposition}

\begin{proof}
$\kappa_{\mathrm{base}}$ satisfies \ax{CF9$'$} and, with any computable $\tau$, \ax{CF10$'$}, so the cores
$\{\ax{CF4},\ax{CF5},\ax{CF10}\}$ and $\{\ax{CF5},\ax{CF6},\ax{CF10}\}$ do not survive. The cores containing
\ax{CF4} are duplicated with \ax{CF9$'$}, which is equivalent to it. \ax{CF3} and \ax{CF5} make $\kappa$
total and equal to $\kappa_{\mathrm{port}}$, which fails \ax{CF10$'$} in the copied-action environment.
\ax{CF2}, \ax{CF5}, and \ax{CF6} likewise give $\kappa_{\mathrm{port}}$ by Lemma~\ref{lem:projection}. The
four-element core is minimal: deleting \ax{CF2}, \ax{CF5}, \ax{CF6}, or \ax{CF10$'$} leaves a set satisfied
by $\kappa_{\mathrm{base}}$, $\kappa_{\mathrm{slot}}$, $\kappa_{\mathrm{flat}}$, or $\kappa_{\mathrm{port}}$.
For exhaustiveness, let $\Sigma$ be core-free and contain \ax{CF5} (otherwise $\kappa_{\mathrm{slot}}$
witnesses it). If $\ax{CF3}\in\Sigma$, $\kappa_{\mathrm{port}}$ witnesses it. If \ax{CF4} or \ax{CF9$'$} is in
$\Sigma$, \ax{CF2} is excluded and $\kappa_{\mathrm{base}}$ witnesses it. Otherwise, if \ax{CF10$'$} is in
$\Sigma$, either \ax{CF6} is absent and $\kappa_{\mathrm{flat}}$ witnesses $\Sigma$, or \ax{CF2} is absent
and $\kappa_{\mathrm{base}}$ does; if \ax{CF10$'$} is absent, $\kappa_{\mathrm{port}}$ does.
\end{proof}

The four-element core reflects the splitting of \enquote{defined and calibrated} into separate
requirements. Theorem~\ref{thm:calibration} is less sensitive to such bookkeeping, since its witnesses
are all total.

The word \enquote{realisable} also needs a benchmark. If realisability freezes every embedded
component at its reference implementation, an alternative taking $a$ obtains $W_{E,b}(a)$, and full
realisability in that sense forces $\kappa=\kappa_{\mathrm{port}}$; slot realisability forces
$\kappa=\kappa_{\mathrm{slot}}$. Neither follows from the word, and the benchmark must be supplied.

\subsection{Program domains and global totality}\label{sec:programmap}

The finite theorem extends to a program domain. Let the components be programs $c_i(d,a)\in B_i$,
promised to terminate on every string $d$ and every legal $a$, with a finite rational host, and with the
environment's description supplied as $d$. A fixed algorithm is allowed as a component $c(d,a)=F(d)$,
and it is independent in the declared sense, being invariant in the live parameter, even when $F$
happens to be the agent under evaluation. When \ax{CF10} asks only for a promise-total $\tau$, the four
kernels remain promise-computable, the slot optimiser halts on every valid description, and
Theorem~\ref{thm:cfmap} holds unchanged; no algorithm deciding membership in the promise class is
needed. Requiring $\tau$ to be globally total changes the map.

\begin{proposition}\label{prop:globalmap}
On the program domain, with \ax{CF10} requiring a globally total $\tau$, each of
$\{\ax{CF3},\ax{CF10}\}$, $\{\ax{CF4},\ax{CF10}\}$, $\{\ax{CF6},\ax{CF10}\}$, and $\{\ax{CF9},\ax{CF10}\}$ is
inconsistent. The minimal inconsistent subsets are these four together with $\{\ax{CF5},\ax{CF9}\}$,
$\{\ax{CF2},\ax{CF4},\ax{CF5}\}$, and $\{\ax{CF3},\ax{CF4},\ax{CF5}\}$, and the maximal consistent subsets are
$\{\ax{CF1},\ax{CF2},\ax{CF3},\ax{CF4},\ax{CF6},\ax{CF7},\ax{CF8},\ax{CF9}\}$ (witnessed by
$\kappa_{\mathrm{slot}}$, without a globally total optimiser) and the sets of $\kappa_{\mathrm{port}}$,
$\kappa_{\mathrm{base}}$, and $\kappa_{\mathrm{flat}}$ in Theorem~\ref{thm:cfmap}.
\end{proposition}

\begin{proof}
For a globally total candidate $\tau$, the component $c(d,a)=\tau(d)$ is admissible. With the host
$u(a,x)=\one[a\ne x]$ and $b=\tau(d(E_\tau))$, the component is constant in the live parameter, so
\ax{CF3}, \ax{CF4}, \ax{CF6}, and \ax{CF9} each force the values $0$ at $b$ and $1$ at $1-b$, contradicting
\ax{CF10}. The finite-domain cores remain, and those containing \ax{CF10} are now subsumed. If \ax{CF10} is
in a core-free set, \ax{CF3}, \ax{CF4}, \ax{CF6}, and \ax{CF9} are excluded, and the set lies in that of
$\kappa_{\mathrm{flat}}$, which every globally total agent ratifies; without \ax{CF10}, the case analysis of
Theorem~\ref{thm:cfmap} applies.
\end{proof}

This is the fixed-copy diagonal of Theorem~\ref{thm:fixedcopy}, factored through the desiderata, and
there is no conflict with promise-total optimisation: embedding the promise-total optimiser as $F$
need not yield a component that terminates on every string, so the diagonal need not belong to the
domain. If \ax{CF3} uses a syntactically narrower independence class that excludes embedded agent
algorithms, $\{\ax{CF3},\ax{CF10}\}$ need not be inconsistent; the proposition uses invariance in the
live parameter.

\subsection{The candidate semantics located}\label{sec:candidates}

The semantics of Section~\ref{sec:inference} can now be placed.

\emph{Transport through compositions.} For postcomposition alone, $h_0\circ T=h_1\circ T$ means that $h_0$
and $h_1$ agree on the image of $T$, which determines their values at every legal action exactly when
$T$ is surjective. So the relevant property is surjectivity, not injectivity, and in any case no
function from infinitely many descriptions to a finite action set is injective. Arbitrary contexts
restore the ambiguity for every agent: with $t=T(d)$, the contexts $C_0[P](x)=0$ and
$C_1[P](x)=\one[P(d)\ne t]$ are both the constant $0$ at $P=T$ and differ under a change of the output at
$d$. This is the environment of Theorem~\ref{thm:transport} in program form. With tagged contexts,
conditional transport, and grounded extensionality the core is $\{\ax{CF2},\ax{CF4},\ax{CF5}\}$; with
transport that asserts existence it is $\{\ax{CF4},\ax{CF5}\}$; with parametric extensionality the
contexts are not identified and $\kappa_{\mathrm{slot}}$ works.

\emph{Local matching} records each component's actual queried inputs, selects those components that
agree with $T$ there, and freezes the selection. On an appropriate termination promise it satisfies
\ax{CF1}, \ax{CF2}, agent extensionality, implementation extensionality at fixed inputs, and
\ax{CF8}. In the environment of Theorem~\ref{thm:localmatch} it violates \ax{CF3} and \ax{CF6} under
fixed-component independence, transport for transformed copies, \ax{CF9} under the ordinary
fixed-component benchmark, and \ax{CF10}. That does not make $\{\ax{CF1},\ax{CF2},\ax{CF5},\ax{CF8},\ax{CF10}\}$
inconsistent, which $\kappa_{\mathrm{flat}}$ satisfies.

\emph{Global matching} misses $1-T$, which is never equal to $T$, while transport requires its value
to become $1-a$. Its kernel can also be uncomputable with every component total: with $T\equiv0$ and
the components $g_e$ of Theorem~\ref{thm:haltred}, the environment returning $g_e(d(E_e))$ has actual
payoff $0$, and its value at action $1$ is $1$ if $e$ never halts and $0$ if it halts.

\emph{Fresh rematching} violates \ax{CF7} (Proposition~\ref{prop:rematch}: the final payoff is $1$ where
direct intervention gives $0$), but no desideratum requires it; retaining the original matching set, or
a supplied graph, avoids the failure. Literal syntactic substitution needs a one-pass convention for
the same reason, since the replacement code $T^{d\leftarrow a}$ mentions $T$ away from $d$ and repeated
rescanning can expand without end.

\emph{Material implication} yields the relation $\mathcal V$ of Section~\ref{sec:logical}, not a
kernel. In an independent five-and-ten the attainable payoff for ten is ten, so a selector assigning
$0$ violates \ax{CF3} and \ax{CF9}, but a selector from the relation need not: an additional rule could
choose the intervention value. A kernel assigning $(5,0)$ to an agent that takes five ratifies that
agent, so the failure is one of calibration and realisability, not of \ax{CF10}.

\emph{Fixed copies} are transported only if the equivalence convention erases their role. If
\ax{CF5} identifies a designated copy with a fixed program computing the same closed function, the
transport owed to the first passes to the second, against conservativity at a non-actual action;
if extensionality preserves the open interface, nothing passes. Under the fixed-component
realisation benchmark the global fixed-copy diagonal makes that realisability inconsistent with
\ax{CF10} alone.

\emph{Guarded logical semantics} for terminating programs is $\kappa_{\mathrm{base}}$, by
Proposition~\ref{prop:guardedarith}: it ratifies no agent under \ax{CF10} on a binary problem, and
every agent under \ax{CF10$'$}.

\emph{Logical induction with exploration} computes, at a stage $n$, values such as
\[
  \kappa_n(a)=\frac1n\sum_{i=0}^{n-1}\mathbb P_n\bigl(\bar U>i/n\mid A_n=a\bigr),
\]
using the capped conditional of \textcite[Definitions~4.7.1 and~4.8.2]{garrabrant2016li}, with value $1$
when the conditioning probability is zero. Nothing in the logical induction criterion makes these
values calibrated at a given stage: setting every price at one stage to zero preserves the criterion
by finite-perturbation closure, the capped conditional then assigns $1$ to every conditioned
proposition, and a constant-zero payoff receives $\kappa_n(0)=\kappa_n(1)=1$, violating \ax{CF1},
\ax{CF3}, and \ax{CF9}. The scope of exploration is a further input: with $u(a,b)=a+2b$ at the baseline
$a=b=0$, exploration that changes only the port yields the values $(0,1)$, while exploration that
changes the port and a designated copied bit together yields $(0,3)$. Conditioning computes the
consequences of a supplied exploration mechanism; it does not choose the mechanism.

\subsection{Robustness}

\begin{center}
\small
\renewcommand{\arraystretch}{1.2}
\begin{tabularx}{\textwidth}{@{}>{\raggedright\arraybackslash}p{10em}>{\raggedright\arraybackslash}X@{}}
  \toprule
  variation & effect \\
  \midrule
  globally total or promise-total agents & the finite obstruction survives both; arbitrary fixed total
     algorithms add the pairwise cores of Proposition~\ref{prop:globalmap} \\
  fixed algorithms admitted & fixed-target diagonals return when alternatives leave the target fixed;
     they are not needed for the finite obstruction \\
  fixed algorithms excluded & requires a concrete grammar: declared agent interfaces, no unrestricted
     evaluation or source reflection, intervention-total contexts; \enquote{nothing derived from $T$}
     is not a definition, since every $g$ is $C[T]$ for a context ignoring its argument \\
  informative agents & injectivity is unavailable; surjectivity removes only the $h\circ T$ ambiguity \\
  distribution-valued kernels & exact transport can demand the distinct point masses $\delta_0$ and
     $\delta_1$ for one grounded representation; uncertainty does not remove the contradiction \\
  mixed agents & some ratifiability failures acquire mixed fixed points, depending on prediction and
     sampling semantics (Section~\ref{sec:mixedexist}); the transport contradiction is untouched \\
  policy-level interventions & every one-input witness embeds; finite policy tables support the slot
     theorem; infinite policy spaces need separate assumptions \\
  bounded or certified languages & can support a globally total optimiser with a default on invalid
     descriptions; the unrestricted fixed-copy diagonal need not belong to them \\
  \bottomrule
\end{tabularx}
\end{center}

% -----------------------------------------------------------------------------------
\section{The witness problems revisited, and counterpossibles}\label{sec:counterpossibles}

\subsection{Verdicts of the candidate semantics}

The witness problems of Section~\ref{sec:witnesses} can now be given as programs and evaluated
under each candidate semantics. Let $0$ stand for one-boxing or cooperating and $1$ for two-boxing or
defecting, and use the Newcomb and twin payoffs $u(a,b)=2(1-b)+a$, where $b$ is the predictor's
output or the twin's action. With $t=T(d)$ the reference agent's output, write the vectors of values
in the action order $(0,1)$ as
\[
  P_t=\bigl(2(1-t),\,2(1-t)+1\bigr),\quad S=(2,1),\quad D_t=(t,\,1-t),\quad Z=(0,0),\quad H=(5,10).
\]
$P_t$ freezes the predictor; $S$ transports a shared copy; $D_t$ and $Z$ are the frozen and transported
vectors of the mismatch diagonal $\one[a\ne b]$ of Theorem~\ref{thm:fixedcopy}; and $H$ gives the
payoffs of five-and-ten. In Newcomb and the twin the dependence is a designated oracle call; in the
diagonal the copy of the agent is a fixed component.

\begin{center}
\small
\renewcommand{\arraystretch}{1.2}
\begin{tabular}{@{}lcccc@{}}
  \toprule
  semantics & Newcomb & twin & five-and-ten & fixed-copy diagonal \\
  \midrule
  port only & $P_t$: two-box & $P_t$: defect & $H$: ten & $D_t$: opposite of $t$ \\
  designated slots & $S$: one-box & $S$: cooperate & $H$: ten & $D_t$: opposite of $t$ \\
  slots and literal copies & $S$ & $S$ & $H$ & $Z$: tie \\
  matching & $S$ & $S$ & $H$ & $Z$: tie \\
  supplied graph & $S$ or $P_t$ & $S$ or $P_t$ & $H$ & $Z$ or $D_t$ \\
  guarded, arithmetic & actual only & actual only & actual only & actual only \\
  logical induction & none & none & none & none \\
  \bottomrule
\end{tabular}
\end{center}

The first two rows are $\kappa_{\mathrm{port}}$ and $\kappa_{\mathrm{slot}}$, that is, $O_{\mathrm{const}}$
and $O_{\mathrm{oracle}}$; matching is local or global; logical induction is at an unspecified stage.
For the dependency graph the first entry applies when the graph links the predictor or copy to the
decision node and the second when it treats it as fixed. \enquote{Actual only} means that every
non-actual action has no value (Proposition~\ref{prop:guardedarith}), so there is no recommendation
under \ax{CF10}. For logical induction, nothing in the criterion fixes a numerical verdict at a given
stage; the all-zero-price stage of Section~\ref{sec:candidates} assigns the normalised upper bound to
both actions, giving $(3,3)$, $(3,3)$, $(10,10)$, and $(1,1)$.

The table makes visible what Section~\ref{sec:witnesses} took as given. In Newcomb's problem and the
twin prisoner's dilemma the conflict between \cri{1} and \cri{2} is the difference between the second
row and the first: whether the predictor's output, or the twin's action, is transported with the
agent's output or frozen at its baseline. A semantics that transports it one-boxes and cooperates;
one that freezes it two-boxes and defects. The model of Section~\ref{sec:model} made this choice by
stipulating $\mu_\pi$; a program does not make it, and the rows below the second show the ways an agent
may try to make it from source. On the fixed-copy diagonal, conversely, transporting the copy turns a
lost problem into a tie, which is why semantics that recognise copies escape
Theorem~\ref{thm:fixedcopy} and why Theorem~\ref{thm:complement} needs a copy that is transformed rather
than equal.

\subsection{Counterpossibles}

Every non-actual output of a deterministic computation is, relative to the computation's complete
description, an impossibility. Logical counterfactuals are therefore \emph{counterpossibles},
conditionals with impossible antecedents, and the philosophical literature on them bears on the
results of Section~\ref{sec:operators}.

On the orthodox semantics of \textcite{lewis1973counterfactuals}, a counterfactual with an impossible
antecedent is vacuously true. Vacuism is the counterfactual analogue of material implication in
Section~\ref{sec:logical}: it assigns every consequent to every non-actual output, and so provides a
relation, not a kernel. Non-vacuists answer with impossible worlds, in which impossible antecedents
hold without everything holding \autocite{nolan1997,brogaardsalerno2013}. \textcite{williamson2015,williamson2018}
defends vacuism, and one of his arguments has the structure of Theorem~\ref{thm:transport}. Suppose the
intension of $A\mathrel{\Prov\!\!\to}B$ is a function of the intensions of $A$ and of $A\wedge B$. If $A$ is
impossible, both intensions are empty for every $B$, so all counterfactuals with antecedent $A$ have the
same intension; since $A\mathrel{\Prov\!\!\to}A$ is true everywhere, every counterpossible is true.
\textcite{bertoetal2018} reply by admitting impossible worlds whose valuations discriminate more finely,
and discuss the conflict between reflexivity, closure under classical consequence, and non-trivial
treatment of contradictions, together with failures of substitution inside counterfactual contexts.

The analogy is exact at one level: an identification appropriate to actual or possible behaviour need
not preserve the distinctions needed under counterpossible suppositions. Williamson's argument
identifies antecedents by their possible-worlds intension, which is empty for every impossibility;
grounded extensionality identifies components by their instantiated behaviour, which is the same for
the identity context and the constant context at the baseline. In both cases the identification
collapses exactly the distinctions that non-vacuous counterfactuals need.

Impossible worlds can represent every kernel of Section~\ref{sec:operators}. Given a baseline-correct
kernel on finite environments, retain an actual world with the actual action and payoff; for every
defined alternative $a$ introduce a world recording action $a$ and payoff $\kappa(E,b,a)$, selected when
the antecedent \enquote{the output is $a$} is evaluated; and for an undefined alternative select several
worlds with incompatible payoffs, so that no value is determined. The selected worlds are impossible
relative to the complete description of the computation. Retaining the origin of each world and
overwriting its target realises \ax{CF7}. The construction represents $\kappa_{\mathrm{slot}}$,
$\kappa_{\mathrm{port}}$, $\kappa_{\mathrm{base}}$, and $\kappa_{\mathrm{flat}}$ alike, and therefore
derives none of them: impossible worlds can record unattainable values, which \ax{CF9} separately
excludes, and the selection relation encodes a counterfactual policy chosen in advance. Impossible-world
semantics makes room to preserve distinctions; it does not recover dependence information that a
closed program's actual behaviour has lost.

Two further precedents are close. \textcite{lewis1976} proved a triviality theorem about representing
conditional probabilities as probabilities of fixed conditional propositions; it does not prohibit
numerical conditional expectations or any kernel of Section~\ref{sec:operators}, but it is the classical
example of a semantic identification that cannot respect the required distinctions.
\textcite{taylor2020} discusses changing the agent's source code together with its hypothetical policy,
and the underdetermination of such changes, as an alternative to logical counterfactuals; it is closely
related to the slot semantics but contains no axiomatic impossibility. The discussion of embedded
agency by \textcite{demskigarrabrant2019} makes the point that an externally supplied agent slot already
contributes disputed structure, without a theorem about numerical kernels.

\subsection{What an optimality theorem requires}

The results of this part combine into a list of the assumptions that any optimality theorem for an
agent reading descriptions must state, because each of them changes some answer above:
\begin{enumerate}[label=(\roman*)]
  \item the atomic interfaces at which the agent is consulted;
  \item which contexts are intervention-total;
  \item whether implementations are compared on preserved inputs;
  \item whether fixed algorithms are admitted as constants;
  \item whether agents must halt globally or only on a promise class;
  \item the substitution benchmark by which alternatives are realised;
  \item how undefined alternatives are treated in ratifiability.
\end{enumerate}
They are mathematical assumptions, not implementation details that may be dropped once the theorem is
stated. Supplying them is supplying what Section~\ref{sec:model} supplied with $\mu_\pi$.


% -----------------------------------------------------------------------------------
\section{Conclusion}\label{sec:conclusion}

The question was whether a decision theory can be optimal on every fair decision problem, and the
expectation was that conflicting criteria of optimality would rule it out. The answer has four
parts, the first three for problems that supply their counterfactuals and the fourth for problems
that do not.

First, \emph{ex-ante optimality is always attainable}. On every class of well-defined finite
problems, choosing an ex-ante maximiser in each problem is a decision theory
(Theorem~\ref{thm:universal}). Because lookup tables count as decision theories, an impossibility
for criteria that judge each problem separately is never more than a single problem on which their
acceptable sets are disjoint (Lemma~\ref{lem:pointwise}).

Second, \emph{within a problem the conflicts are shallow}. Among ex-ante optimality, structural
causal dominance, ratifiability, existence, commitment stability, and a reference-policy evidential
criterion, every inconsistency has a core of one or two criteria, and the maximal consistent
combinations have explicit total witnesses (Theorems~\ref{thm:map15} and~\ref{thm:map16}). The
central conflict is Nozick's. It is possible exactly when the law of the states preceding a choice
depends on the prescription (Theorem~\ref{thm:characterisation}); it survives imperfect prediction,
randomisation, and every restriction on prediction targets short of excluding anticipation; and it
obeys an exact trade-off between ex-ante and causal regret (Proposition~\ref{prop:tradeoff}). The
conflict between dominance and commitment stability, by contrast, is confined on its witness to
infallible prediction, whereas its expectational counterpart persists whenever one-boxing has the
higher ex-ante value (Proposition~\ref{prop:noisycommit}).

Third, \emph{across problems, the depth depends on the semantics of menus}. On the causal matrix, with
common affine invariance, the cross-problem axioms reduce to singleton and pairwise cores
(Theorem~\ref{thm:pairwise}). Row-translation invariance---evaluation insensitive to the baseline
of each belief state---produces a minimal four-axiom impossibility: neutrality, row translations,
and contraction force universal choice (Theorem~\ref{thm:circulant}), exactly on domains satisfying
(DR) (Theorem~\ref{thm:DR}), a condition that survives imperfect prediction, bounded action sets,
and, with affine invariance, bounded payoffs. Principal binary independence, with neutrality, affine
invariance, ordering, non-triviality, and monotonicity, characterises ex-ante evaluation
(Corollary~\ref{cor:three}). Both results presuppose principal restriction, on which removing an
option also removes the belief state it would induce. Under column or fixed-reference restriction,
fixed-prior causal evaluation satisfies every axiom considered except diagonal optimality,
self-sovereignty, and principal binary independence (Proposition~\ref{prop:colaverage}), and contraction under any semantics is compatible with
$2^{\cont}$ rules that do not maximise ex-ante value (Proposition~\ref{prop:cardinality}).

Fourth, \emph{when the dependence must be inferred, the obstruction is the semantics of alternatives}.
For agents that read source code, a designated substitution interface and a termination promise give
a single optimal algorithm (Theorem~\ref{thm:promiseopt}). Three impossibilities hold without any
counterfactual subtlety: queries about other problems (Proposition~\ref{prop:foreign}), fixed copies of
the agent against agents that halt on every input (Theorem~\ref{thm:fixedcopy}), and exact selection
among computable expected values, which is uncomputable even with no predictor
(Theorem~\ref{thm:exactsel}) but becomes computable as soon as either ties or arbitrarily small
advantages are excluded (Proposition~\ref{prop:approxopt}). The ways of recognising the agent in code
each fail: literal matching misses encoded copies, local matching creates dependence from accidental
agreement (Theorem~\ref{thm:localmatch}), global matching misses transformed copies and can make
optimality decide the halting problem (Theorems~\ref{thm:complement} and~\ref{thm:haltred}), and
proof-based counterfactuals about the actual agent are spurious (Theorem~\ref{thm:fiveten}) or undefined
(Proposition~\ref{prop:guardedarith}). On the axiomatic level, every inconsistency among the desiderata
for counterfactual operators contains grounded extensionality (Theorem~\ref{thm:cfmap}). Together with
calibration it forces causal evaluation (Lemma~\ref{lem:projection}); with calibration and ratifiability
it forms a minimal inconsistent triple whose witness is Death in Damascus
(Theorem~\ref{thm:calibration}); and replacing it by parametric extensionality leaves ex-ante evaluation
satisfying every desideratum (Corollary~\ref{cor:parametric}). The fixed-copy diagonal reappears as
four pairwise cores when the optimiser must halt on every input (Proposition~\ref{prop:globalmap}).

The thesis that the paper can defend is therefore the following.

\begin{quote}
\itshape
On any class of well-defined, fair problems containing a sufficiently predictive Newcomb problem, no
assignment of policies satisfies both ex-ante optimality and causal dominance, although an ex-ante
optimal assignment exists. Within problems, every inconsistency among the natural criteria has a
core of size at most two. Across problems, on the unrestricted causal-matrix domain, principal binary
independence together with neutrality, ordering, affine invariance, non-triviality, and monotonicity
characterises ex-ante ranking, and principal contraction together with neutrality is incompatible
with non-trivial, baseline-invariant choice on every domain satisfying (DR). Neither cross-problem
result survives when a smaller menu retains the belief model of the larger one. For computable agents
that read descriptions, optimality is attainable on the promise domain of any designated substitution
semantics, and fails, apart from fixed-copy diagonals and the uncomputability of exact selection, only
through the semantics itself: no total counterfactual operator can both transport a computation
through alternative outputs and identify every presentation with the same actual behaviour, and
calibration with such an identification is inconsistent with ratifiability. In every case the
obstruction lies in how the agent's alternatives are specified, not in choosing among them.
\end{quote}

\begin{table}[htbp]
  \centering\small
  \renewcommand{\arraystretch}{1.25}
  \setlength{\tabcolsep}{5pt}
  \captionabove{The principal results, with their status. \enquote{Not found} means not found in the
                literature searched; no priority is claimed.}
  \label{tab:results}
  \begin{tabularx}{\textwidth}{@{}>{\raggedright\arraybackslash}X
                                    >{\raggedright\arraybackslash}p{6.8em}
                                    >{\raggedright\arraybackslash}p{8.6em}@{}}
    \toprule
    result & where & status \\
    \midrule
    ex-ante optimality vs dominance (Newcomb) & Prop.~\ref{prop:newcomb} & \textcite{nozick1969} \\
    no ratifiable pure act (Death in Damascus) & Prop.~\ref{prop:did} & \textcite{gibbardharper1978}; ratifiability literature \\
    universal ex-ante selection; pointwise reduction & Thm.~\ref{thm:universal}, Lem.~\ref{lem:pointwise} & elementary \\
    characterisation of the dominance conflict & Thm.~\ref{thm:characterisation} & not found \\
    Bomb: ill-posedness and repair & Props.~\ref{prop:bombill}, \ref{prop:bomb} & case from \textcite{macaskill2019} \\
    within-problem maps & Thms.~\ref{thm:map15}, \ref{thm:map16} & not found \\
    noisy commitment and couplings & Props.~\ref{prop:noisycommit}, \ref{prop:strassen} & uses \textcite{strassen1965} \\
    regret trade-off & Prop.~\ref{prop:tradeoff} & not found \\
    mixed existence under continuity & Thm.~\ref{thm:mixedexist} & \textcite{halpern2020}; \textcite{belletal2021} \\
    realisability of every causal matrix & Thm.~\ref{thm:realise} & elementary \\
    \ax{R} as factorisation; \ax{R} is not causal rationality & Thm.~\ref{thm:factor}, Prop.~\ref{prop:regret} & not found \\
    Pareto vs self-sovereignty & Prop.~\ref{prop:PS} & \textcite{briggs2010} \\
    pairwise reduction under \ax{F} & Thm.~\ref{thm:pairwise} & not found \\
    circulant theorem; map with \ax{R} & Thms.~\ref{thm:circulant}, \ref{thm:augmented} & not found; ingredients in \textcite{finkruiz1986}, \textcite{brandtbrillharrenstein2016} \\
    exact domain criterion (DR); physical restrictions & Thm.~\ref{thm:DR}, Props.~\ref{prop:bounded}--\ref{prop:ceiling} & not found \\
    menu semantics; weight preservation; column maps & Sec.~\ref{sec:menus} & not found; cf.\ \textcite{gallow2026} \\
    binary characterisation of ex-ante ranking & Thm.~\ref{thm:selfesteem}, Cor.~\ref{cor:three} & core in \textcite{podgorski2022} \\
    contraction without the diagonal; $2^{\cont}$ families & Thm.~\ref{thm:lambda}, Props.~\ref{prop:cardinality}--\ref{prop:NFSO} & not found \\
    mixed and evidential extensions & Sec.~\ref{sec:mixed} & not found \\
    uniform optimiser on the promise class & Thm.~\ref{thm:promiseopt}, Prop.~\ref{prop:proofopt} & elementary \\
    foreign queries; fixed-copy diagonal & Prop.~\ref{prop:foreign}, Thm.~\ref{thm:fixedcopy} & \textcite{skalse2021}; modal form in \textcite{fallenstein2014evil} \\
    no computable exact optimiser & Thm.~\ref{thm:exactsel} & recursive inseparability; application not found \\
    matching semantics and their failures & Thms.~\ref{thm:localmatch}--\ref{thm:haltred}, Prop.~\ref{prop:rematch} & not found \\
    computable five-and-ten & Thm.~\ref{thm:fiveten} & mechanism in \textcite{slepnev2012}, \textcite{bensontilsen2014} \\
    guarded counterfactuals for programs & Prop.~\ref{prop:guardedarith} & semantics from \textcite{fallenstein2015third} \\
    projection lemma; counterfactual-operator maps & Sec.~\ref{sec:operators} & not found; cf.\ \textcite{williamson2015} \\
    \bottomrule
  \end{tabularx}
\end{table}

Several questions remain open. The first is normative rather than mathematical: which semantics of
menu restriction is the right one. Principal restriction is forced by the causal matrix once
law-preserving removal is assumed, but an announced restriction may change a predictor's behaviour,
and a fixed belief model is the natural reading of fixed-prior causal theories. The cross-problem
impossibilities stand or fall with that choice, and the paper has shown only how they depend on it.
The second is whether structural dominance and commitment stability are jointly satisfiable on a
whole class of problems with imperfect prediction: Proposition~\ref{prop:noisycommit} removes the
Newcomb witness, and Proposition~\ref{prop:strassen} decides dominance among commitments, but no
general consistency theorem has been proved. The third concerns the reduced form. The causal matrix
discards structural states and the local structure of multi-stage problems, so that column Pareto is
not structural dominance and global ratifiability is not local ratifiability. The context environments
of Section~\ref{sec:operators} are one richer reduced form, which keeps the factorisation of the matrix
through the components, and on it the semantic choice that the matrix leaves implicit becomes an axiom;
a cross-problem theory of rules on that form has not been developed. The fourth concerns
regularity: the cardinality results use the freedom to define rules arbitrarily on orbits, and the
classification of rules subject to continuity or measurability is open. The fifth is a complete
consistency map for mixed policies, which requires a choice among several inequivalent semantics for
output, prediction, and contraction. The sixth is whether parametric extensionality can be
implemented for unrestricted programs: Corollary~\ref{cor:parametric} assumes the contexts are
supplied, and extracting them from a closed program is the problem of extracting a dependency graph,
for which no procedure is known and which global matching shows cannot be equality of computed
functions. The seventh is whether a counterfactual semantics can be chosen on grounds other than the
desiderata considered here, for instance by a learning process that is calibrated in the limit, and
what such a choice would do to the map. Finally, the status in the literature of the circulant
theorem and of the consistency maps has been checked only by a bounded search; their ingredients are
classical, and an earlier statement may exist.


\appendix
% -----------------------------------------------------------------------------------
\section{Register of stipulations}\label{app:stipulations}

The results are theorems about the definitions displayed in the text. The following choices are made
there, and each could be made differently; where a different choice is known to change a result, the
section that discusses it is given.

\begin{enumerate}[label=\arabic*.,leftmargin=2em,itemsep=.3em]
  \item Policies are deterministic unless a mixed variant is stated (Sections~\ref{sec:robust},
        \ref{sec:mixed}).
  \item Every admissible policy induces a well-defined probability law (Definition~\ref{def:welldefined});
        self-referential problems that lack one are excluded (Remark~\ref{rem:skalse}).
  \item No information set occurs twice on a path (Definition~\ref{def:problem}).
  \item Structural environment plans, including responses on branches not taken, are part of the problem
        (Definition~\ref{def:structural}); how outcomes on alternative branches are coupled is fixed by
        them.
  \item Causal deviations freeze the baseline distribution over plans, and a local deviation keeps the
        policy's other prescriptions (definition of $Q_\pi$).
  \item Dominance quantifies over structural states, including those of probability zero under the
        evaluated policy (Section~\ref{sec:criteria}); with supported states only, \cri{2} becomes
        unsatisfiable in Death in Damascus.
  \item Positive reach is measured under the evaluated policy.
  \item Predictor error is given by a kernel; the Newcomb calculations use symmetric error.
  \item A predictor targets specified policy coordinates; \enquote{the same computation} is represented
        by a stipulated dependence (Sections~\ref{sec:twin}, \ref{sec:bomb}).
  \item Commitment is irrevocable, precedes the environment's response, offers the full policy set, and
        uses a specified coupling and structural state set (Definition~\ref{def:commit},
        Section~\ref{sec:couplings}).
  \item Commitment extensions are tagged and well founded, and the class is closed under them.
  \item \cri{5} is the totality already built into the definition of a theory.
  \item The evidential criterion uses the reference-policy law and the null-event convention displayed in
        Section~\ref{sec:evidential}; other evidential semantics change the map (Section~\ref{sec:evsem}).
  \item Mixed variants specify sampling, predictor access, information, counterfactual coupling, and
        continuity (Sections~\ref{sec:mixedexist}, \ref{sec:mixed}).
  \item Cross-problem rules output a reflexive relation whose maximal elements are required to be nonempty
        (Definition~\ref{def:rule}); neutrality acts by bijections between label sets, and invariances act
        on the relation.
  \item Ordering concerns the output relation at each input, not rationalisability of each choice set.
  \item Menu restriction is principal unless stated otherwise; principal restriction is forced by
        law-preserving removal, and column and fixed-reference restriction require additional data
        (Section~\ref{sec:menus}).
  \item Non-triviality is existential unless a universe is specified (Theorem~\ref{thm:five}).
  \item On restricted domains, an invariance applies when both a matrix and its transform are admissible
        (Section~\ref{sec:translations}).
  \item Equivalence of problems and parameterisation of rules are made precise before taxonomy claims are
        stated (Section~\ref{sec:taxonomy}).
\end{enumerate}


\begin{center}
  \vspace{1.2\baselineskip}
  {\large ⁂}
\end{center}

\printbibliography[heading=bibintoc,title={References}]

\end{document}
